\documentclass[11pt,a4paper]{article}
\usepackage[T1]{fontenc}
\usepackage[utf8]{inputenc}
\usepackage{amsmath,amsthm,mathtools}
\usepackage{newtxtext,newtxmath}
\usepackage[margin=26mm,headheight=15pt]{geometry}
\usepackage{microtype}
\usepackage{xcolor,booktabs,array,tabularx,longtable,enumitem}
\usepackage{fancyhdr}
\usepackage[most]{tcolorbox}
\usepackage{hyperref}
\usepackage{xurl}
\usepackage[nameinlink,noabbrev]{cleveref}
\definecolor{Ink}{HTML}{20344B}
\definecolor{Accent}{HTML}{176B78}
\definecolor{Pale}{HTML}{F0F5F7}
\hypersetup{colorlinks=true,linkcolor=Accent,citecolor=Accent,urlcolor=Accent,
 pdftitle={Lexicographic orders of the well-orderings of the reals},
 pdfauthor={Research report merged from three manuscripts prepared for Vladimir Reshetnikov}}
\pagestyle{fancy}\fancyhf{}
\fancyhead[L]{\small\textsc{Lexicographic orders of real well-orderings}}
\fancyhead[R]{\small\thepage}
\renewcommand{\headrulewidth}{0.3pt}
\setlength{\parskip}{3pt}
\setlength{\parindent}{1.25em}
\setlist{itemsep=3pt,topsep=5pt}
\newtheorem{theorem}{Theorem}[section]
\newtheorem{lemma}[theorem]{Lemma}
\newtheorem{proposition}[theorem]{Proposition}
\newtheorem{corollary}[theorem]{Corollary}
\theoremstyle{definition}
\newtheorem{definition}[theorem]{Definition}
\newtheorem{example}[theorem]{Example}
\newtheorem{question}[theorem]{Research question}
\theoremstyle{remark}
\newtheorem{remark}[theorem]{Remark}
\newcommand{\R}{\mathbb R}
\newcommand{\N}{\mathbb N}
\newcommand{\Q}{\mathbb Q}
\newcommand{\W}{\mathcal W}
\newcommand{\A}{\mathcal A}
\newcommand{\B}{\mathbb B}
\newcommand{\two}{\{0,1\}}
\newcommand{\lex}{\mathrm{lex}}
\newcommand{\restr}{\mathbin{\upharpoonright}}
\newcommand{\cat}{\mathbin{\frown}}
\newcommand{\ran}{\operatorname{ran}}
\newcommand{\dom}{\operatorname{dom}}
\newcommand{\cf}{\operatorname{cf}}
\newcommand{\ci}{\operatorname{ci}}
\newcommand{\dens}{\operatorname{d}}
\newcommand{\cell}{\operatorname{cell}}
\newcommand{\bl}{\ell_2}
\newcommand{\emb}{\hookrightarrow}
\newcommand{\equi}{\mathrel{\simeq_{\!e}}}
\newcommand{\rev}{\mathrm{op}}
\newcommand{\kplus}{\kappa^+}
\newcommand{\powlex}[2]{\bigl({#1}^{#2},<_{\lex}\bigr)}
% Added in the batch-76 merge (sources 06 and 07); see the notation table.
\newcommand{\weight}{\operatorname{w}}
\newcommand{\repr}{\operatorname{repr}_{\R}}
\newcommand{\Reg}{\operatorname{Reg}}
\newcommand{\Wq}{\W/{\sim_{\mathrm{fin}}}}
\newcommand{\chm}{\chi^-}
\newcommand{\chp}{\chi^+}
\newcommand{\con}{\mathfrak c}
\newcommand{\GapChar}{\operatorname{GapChar}}
\newcommand{\Jinj}{\mathcal I}
\newcommand{\JU}{\mathcal J}
\newcommand{\merge}{\textsf{[merge]}}
\crefname{question}{research question}{research questions}
\Crefname{question}{Research question}{Research questions}
\newtcolorbox{resultbox}{colback=Pale,colframe=Accent,boxrule=.6pt,
 arc=1mm,left=3mm,right=3mm,top=2mm,bottom=2mm,breakable}
\newcolumntype{Y}{>{\raggedright\arraybackslash}X}
\title{\color{Ink}\LARGE Lexicographic orders of the\\well-orderings of the reals\\[7pt]
\large Ordinal spectra, coding length and representation rank,\\finite-tail geometry, characters, gaps, and self-absorption}
\author{Research report merged from three independent manuscripts\\prepared for Vladimir Reshetnikov (sources 02, 06 and 07 of batch 76)}
\date{Sources dated 2 October 2026; merged 2 October 2026}
\begin{document}
\maketitle
\begin{abstract}
Let $\kappa=2^{\aleph_0}$, regarded as an initial ordinal. A well-ordering of
$\R$ has a unique exhaustive enumeration by an ordinal of cardinality $\kappa$.
We compare two such enumerations at their first unequal real entry. The
resulting linear order $\W$ has cardinality $2^\kappa$, contains every linear
order of cardinality at most $\kappa$, and contains an ordinal $\beta$ exactly
when $\beta<\kappa^+$. Nevertheless, its least binary lexicographic coding
length is $\kappa^+$: it contains every binary lexicographic cube of smaller
length but does not contain the cube of length $\kappa^+$.

The restriction $\W_\kappa$ to enumerations of length exactly $\kappa$ is
quite different. It is equimorphic with $\two^\kappa$ in lexicographic order,
as is the actual ultrafilter-enumeration index in the Kanovei--Shelah
construction. We prove the exact fixed-type classification
$\W_{\lambda+n}\equi\B_\lambda\times_{\lex}n!$ for limit $\lambda$ of cardinality $\kappa$,
characterize every adjacent
pair in $\W$, compute several order-topological invariants, and prove the
sharp absorption theorem $\W^\gamma_{\lex}\emb\W$ if and only if
$\gamma<\kappa^+$. We also analyze Dedekind completion and formulate further
embedding, definability, and formalization problems.

\merge\ This report is built from three manuscripts written independently on
the same question and delivered together in batch~76: source~02 (the base,
whose abstract is the two paragraphs above), source~06 and source~07. They
prove the same core three times; each shared theorem is printed once, with the
other sources' arguments marked as second or third routes where the argument
differs. Sources~06 and~07 add universality in every fixed-type slice; a dense
set of $2^\kappa$ isolated points with a perfect, nowhere dense remainder; the
exact point-character spectrum; the exact gap spectrum of $\W$, in which every
gap has a countable side, with exactly $2^\kappa$ gaps; the gap spectrum of
$\W_\kappa$, which has a $(\cf\kappa,\cf\kappa)$ gap, so that $\W_\kappa$ has no
convex copy in $\W$; the real representation rank $\repr(\W)=\kappa^+$; the
canonical lexicographic-sum decomposition of $\W$ over injections; a dense
finite-condensation quotient equimorphic with $\W$, with a connected
completion; bounded surreal sign orders; further gaps of the Kanovei--Shelah
index; and a well-ordering-based ultrafilter index. Each source answers a
question of another: in ZFC, $\W$ does not embed into $\W_\kappa$, which
source~06 had proved only under $2^{<\kappa}<2^\kappa$; and sources~06 and~07
settle source~02's questions on cut and character spectra and on the quotient.
\end{abstract}
\begin{resultbox}
\textbf{Conventions and status.} All main theorems are in ZFC. All embeddings
are order embeddings, not necessarily continuous maps. Equimorphism means
embeddability in both directions, not isomorphism. Proofs are supplied below;
they have not been independently refereed or checked by a proof assistant.
The results are presented as a research synthesis with explicit derivations,
not as a claim of established historical priority.

\smallskip
\merge\ \textbf{Status of the merged report.} AI-assisted (all three sources
are ``prepared for Vladimir Reshetnikov''; none carries AI wording in its
author line), unrefereed, and not formalized. None of the three sources claims
historical priority, peer review, or machine verification; their finite
checks are illustrations of finite instances, not evidence for the
transfinite assertions.
\end{resultbox}
\clearpage
\setcounter{tocdepth}{1}
\tableofcontents
\clearpage

\section{About this report}\label{lwo:sec:about}

\subsection{Three independent manuscripts}\label{lwo:sec:sources}
\merge\ This section, the notation table, the crosswalk of
\cref{lwo:app:crosswalk} and every paragraph tagged \merge\ were added when
the three manuscripts of \cref{lwo:tab:sources} were merged. Everything
untagged in \cref{lwo:sec:three,lwo:sec:summary,lwo:sec:enumerations,lwo:sec:cubes,lwo:sec:ordinal,lwo:sec:fixed,lwo:sec:length,lwo:sec:all-strata,lwo:sec:absorption,lwo:sec:adjacency,lwo:sec:topology,lwo:sec:KS-topology,lwo:sec:completion,lwo:sec:scope,lwo:sec:research}
and \cref{lwo:app:checks,lwo:app:sources} is source~02 as delivered, with its
labels prefixed by \texttt{lwo:}.
\Cref{lwo:sp:sec:slices,lwo:sp:sec:isolation,lwo:sp:sec:characters,lwo:sp:sec:gaps,lwo:rk:sec:rank,lwo:sp:sec:minimal,lwo:sp:sec:quotient,lwo:sp:sec:surreal,lwo:rk:sec:KS,lwo:rk:sec:greedy,lwo:sec:conclusions}
print sources~06 and~07; each opens with a tagged note naming its source,
and every statement in them names the source it comes from. A reference such as ``06:171'' means line
171 of the main \texttt{.tex} file of source~06 as delivered (archive in the
arrival commit \texttt{6914ccca6}; the manuscripts of sources~06 and~07 are
not shipped with this report).

\begin{table}[ht]
\centering\small
\renewcommand{\arraystretch}{1.15}
\begin{tabularx}{\textwidth}{@{}l>{\raggedright\arraybackslash}p{.27\textwidth}Y>{\raggedright\arraybackslash}p{.13\textwidth}@{}}
\toprule
Source & Archive (batch 76) & Delivered title & Size\\
\midrule
02 (base) & \nolinkurl{lexicographic_well_orderings_of_reals.zip}
 & \emph{Lexicographic Orders of Well-Orderings of the Continuum: ordinal
 spectra, lexicographic complexity, finite-tail geometry, and self-absorption}
 & 1,826 lines; 24 pp.\\
06 & \texttt{Lexicographic\_\allowbreak Well\_\allowbreak Orderings\_\allowbreak of\_\allowbreak Reals~(1).zip}
 & \emph{The Lexicographic Order of All Well-Orderings of the Real Line:
 universality, exact ordinal spectra, finite tails, local characters, and
 Dedekind gaps}
 & 1,522 lines; 23 pp.\\
07 & \nolinkurl{lexicographic_well_orderings_research.zip}
 & \emph{Lexicographic Orders of Real Well-Orderings: ordinal universality,
 representation rank, products, and cut geometry}
 & 1,521 lines; 21 pp.\\
\bottomrule
\end{tabularx}
\caption{\merge\ The three sources. The number of a source is its batch-76
manuscript number. All three are dated 2 October 2026, are ``prepared for
Vladimir Reshetnikov'', and pin the ProveIt revision
\texttt{6ea60e3677bd847a00baf023c3c532c83884530e} (source~02: \cref{lwo:app:sources};
source~06: its Section~1.1 and audit; source~07: its Section~13.2 and audit).
They arrived in commit \texttt{6914ccca6} and were placed in commit
\texttt{2a04b60f2}.}
\label{lwo:tab:sources}
\end{table}

The three texts are \emph{independent}: they share no member file, their
word 8-gram overlap is $0.65$--$0.75\,\%$ pairwise, and their titles,
notation and proofs were written separately. None is a version of another,
and none supersedes another. They nevertheless prove the same core---prefix
freeness, cardinality $2^\kappa$, universality for orders of size at most
$\kappa$, the exact ordinal spectrum below $\kappa^+$, the equimorphism
$\W_\kappa\equi\B_\kappa$, the adjacency criterion, the finite blocks of
size $n!$, and the topology of $\W$ and of $\W_\kappa$---three times.
Sources~02 and~07 further share the coding rank $\kappa^+$, the real powers,
power absorption, the fixed-length square and the Kanovei--Shelah index;
sources~06 and~07 share the point-character spectrum, the gap spectrum of
$\W$ and the dense quotient. Source~02 is the base because it has the widest
scope: every stratum $\W_\alpha$, the Kanovei--Shelah index $\A$, general
alphabets in its coding-length theorem, and completion invariants.

\subsection{How the merge is printed}\label{lwo:sec:how}
\merge\ Each theorem proved by more than one source is stated once. Where
another source proves it by the same argument, a short note after the proof
says so (``also source~06, Theorem~5.2, 06:465, same argument''). Where the
argument differs, the other proof is printed as a marked \emph{second route}
(or \emph{third route}), credited to its source. No result, proof, example,
remark, question or limitation of any source is dropped: the crosswalk
(\cref{lwo:tab:crosswalk}) lists every numbered statement of the three
sources and where it stands.

Two cross-repository duplications are also recorded at their points of use.
Source~02's strict cube hierarchy (\cref{lwo:thm:cube-strict}) re-proves
Corollary \texttt{cor:strict} of the ProveIt report \emph{Every ordinal as a
point-separating game value}; that corollary is cited first and source~02's
cylinder proof is a second route. Source~07's strictness of real powers
(\cref{lwo:rk:thm:strict}) is the ordinal case of a classical theorem of
Kuhlmann \cite[Corollary~2.4]{Kuhlmann}, as source~07 itself says.

\subsection{Questions one source answers for another}\label{lwo:sec:answered}
\merge\ Each source left open a question that another source answers. These
answers are printed as theorems, crediting the answering sources, and the
questions are re-scoped in \cref{lwo:sec:research}.
\begin{enumerate}[label=(\roman*)]
\item Source~06 proves that $\W$ does not embed into $\W_\kappa$ only under
$2^{<\kappa}<2^\kappa$ (06:1053, its Corollary~9.5) and asks for an
unconditional comparison (06:1332, its Research question~14.1). Sources~02
and~07 prove it in ZFC: \cref{lwo:cor:unbounded} (source~02) and source~07's
Corollary~5.4 (07:525), which also excludes every fixed real power
$\R^\gamma_{\lex}$ with $\gamma<\kappa^+$.
\item Source~02 asks for the exact spectra of cuts and point characters
(02:1635--1646). Sources~06 and~07 compute the point-character spectrum and
the gap spectrum of $\W$ (\cref{lwo:sp:thm:point-spectrum,lwo:sp:thm:gapspectrum};
06:913, 06:802; 07:861, 07:941), and source~07 the gap spectrum of $\W_\kappa$
(\cref{lwo:rk:thm:Pgaps}; 07:996).
\item Source~02 asks whether the finite-condensation quotient $\Wq$ is
equimorphic with $\W$ (02:1604--1607). Sources~06 and~07 prove that it is
(\cref{lwo:sp:thm:quotient}; 06:1133, 07:1072).
\end{enumerate}

\subsection{Notation}\label{lwo:sec:notation}
\merge\ The three sources use different letters for the same objects and the
same letters for different objects. The merged text uses the symbols of
\cref{lwo:tab:notation}; every symbol taken from source~06 or~07 and renamed
here is listed in its last column. No normalization changes: every renaming
is a change of letter only.

{\small
\renewcommand{\arraystretch}{1.15}
\begin{longtable}{@{}>{\raggedright\arraybackslash}p{.29\textwidth}>{\raggedright\arraybackslash}p{.17\textwidth}>{\raggedright\arraybackslash}p{.46\textwidth}@{}}
\caption{\merge\ Notions and symbols. ``False reading'' marks a tempting
confusion.}\label{lwo:tab:notation}\\
\toprule
Notion & Here & Sources, and readings to avoid\\
\midrule
\endfirsthead
\toprule
Notion & Here & Sources, and readings to avoid\\
\midrule
\endhead
\bottomrule
\endfoot
Continuum as an initial ordinal & $\kappa=\con=2^{\aleph_0}$ &
02 and 06: $\kappa$ (06 also writes $\con$); 07: $\kappa$ (macro \texttt{\textbackslash k}).
False reading: $\kappa^+$ is the successor \emph{cardinal}, not the ordinal $\kappa+1$.\\
All well-orderings of $\R$ & $\W=\W(\R)$ &
02, 06: $\W$; 07: $\W$ (macro \texttt{\textbackslash WO}). The order compares
enumerations at the first unequal entry in the \emph{usual} order of $\R$; it is
not an order of ordinal types.\\
Enumerations of an ordered set $X$ & $\W(X)$, $\W_\alpha(X)$ &
02: $|X|=\lambda$; 06, 07: $|X|=\nu$. Kept as printed in each source's text.\\
Stratum of length $\alpha$ & $\W_\alpha$ & 02: $\W_\alpha$; 06: $\W_\beta$.\\
Minimal-type slice & $\W_\kappa$ &
06: $\mathcal P$ (\texttt{\textbackslash PP}); 07: $\mathcal P$ (\texttt{\textbackslash Perm}).
False reading: source~07's $\mathcal P$ is \emph{not} the power set
$\mathcal P(\N)$, which source~07 also uses.\\
Kanovei--Shelah index & $\A$, fibres $\A_U$ & 02: $\A$; 07: $\mathcal A$
(\texttt{\textbackslash KS}); ultrafilter: 02 $U$, 07 $D$, here $U$.\\
Binary cube & $\B_\theta=(\two^\theta,<_\lex)$ &
06: $({}^{\theta}2,<_{\lex})$; 07: $2^\theta_{\lex}$. False reading: $2^\kappa$
without a subscript is a cardinal.\\
Real power & $\R^\gamma_{\lex}$ & 06: $({}^{\gamma}\R,<_{\lex})$; 07: $\R^\gamma_{\lex}$.
All functions, no finite-support restriction.\\
Binary coding length & $\bl(L)$ &
02 only: $\min\{\theta:L\emb\B_\theta\}$. It equals the game value
$\mathrm{ps}(L^{\mathrm{dd}})$ of the point-separating report
(its Theorem \texttt{thm:embedding}).\\
Real representation rank & $\repr(L)$ &
07 only, after Giarlotta: $\min\{\alpha:L\emb\R^\alpha_{\lex}\}$. Both
invariants are kept: $\repr(L)\le\bl(L)\le\omega\cdot\repr(L)$; they agree
on $\W$ ($\kappa^+$) and on $\W_\kappa$ and $\A$ ($\kappa$), but not in
general ($\repr(\R)=1$, $\bl(\R)=\omega$).\\
Equimorphism & $L\equi M$ & 02: $\simeq_e$; 07: $\equiv_{\mathrm{emb}}$; 06:
``mutually order-embeddable''. False reading: not isomorphism.\\
Order embedding & $L\emb M$ & 06: $\hookrightarrow_{\mathrm{ord}}$. Not
continuous and not convex unless said.\\
Reverse order & $L^{\rev}$ & 06: $L^{*}$; 07: $L^{\mathrm{op}}$.\\
Prefix cylinder & $C_s$, $C_p$ & 02: $C_s$; 06: cone $C(p)$; 07: $[p]$.\\
Residual set after a prefix $p$ & $Y_p=X\setminus\ran p$ &
02: $Y$; 06: $A_p$; 07: $S\setminus\ran p$. False reading: source~06's $A_p$
is unrelated to the index $\A$.\\
First disagreement & $\delta$, $\Delta(e,f)$ & 02, 07: $\delta$; 06: $\Delta(e,f)$.\\
Lower/upper branch positions & $D^-(e)$, $D^+(e)$ &
06 and 07, same meaning. False reading: unrelated to source~02's
$D_{\kappa,m}=\B_\kappa\times_\lex\{0,\dots,m-1\}$.\\
One-sided point characters & $\chm(e)$, $\chp(e)$ &
06: $\chi^\mp$ (\texttt{\textbackslash cfm}, \texttt{\textbackslash cfp}); 07:
$\chi^\mp$ (\texttt{\textbackslash chm}, \texttt{\textbackslash chp}); 02: ``left
and right local cofinalities''.\\
$\cf\kappa$ & $\cf\kappa$ & 07: $\theta$. Renamed because $\theta$ is an ordinal
variable here ($\B_\theta$).\\
$2^{<\kappa}=\sup_{\lambda<\kappa}2^\lambda$ & $2^{<\kappa}$ & 07: $\tau$.\\
Density, cellularity, weight & $\dens$, $\cell$, $\weight$ & 02: $\dens$,
$\cell$; 06, 07: $d,c,w$. Cellularity is written $\cell$, never $c$, to avoid $\con$.\\
Infinite regular cardinals $\le\kappa$ & $\Reg(\le\kappa)$ & 06: ``infinite
regular $\rho\le\kappa$''; 07: $\operatorname{Reg}_\infty(\le\kappa)$.\\
Gap and its character & $(L,U)$, $(\cf L,\ci U)$; $\GapChar$ &
06: $(L,U)$; 07: $(A,B)$; 02 (\cref{lwo:prop:real-gap}): $(L,R)$. A gap has no
endpoint on either side; a jump is not a gap.\\
Finite-condensation quotient & $\Wq$ & 02: $\Wq$; 06: $D$; 07:
$\W_{\mathrm{fc}}$. False reading: source~06's $D$ is not $D^\pm(e)$.\\
Injections $\kappa\to\R$ & $\Jinj$ & 06: $J$ (not necessarily onto).\\
Well-ordering-based ultrafilter index & $\JU$ & 07: $J$ (bijections
$\kappa\to\mathcal P(\N)$). False reading: sources~06 and~07 both write $J$
for these two different orders.\\
Bounded surreal sign order & $S_{<\theta}$ & 06 only.\\
\end{longtable}
}

\subsection{Non-claims of the three sources}\label{lwo:sec:nonclaims}
\merge\ Every limitation stated by a source is kept where the source states
it; together they are:
\begin{itemize}
\item No source claims historical priority (02: \cref{lwo:app:sources};
06: its abstract and Section~1.1; 07: its abstract and Section~13.2).
Source~07 identifies its real-power strictness as Kuhlmann's classical
theorem; source~06 attributes the characteristic-function coding to the
standard literature (Moore's Section~2); source~02 cites Knoblauch for the
terminology of intrinsic length.
\item The proofs are written ZFC arguments. None has been refereed or checked
in Lean or Rocq. The finite checks of all three sources test finite instances
only (\cref{lwo:app:checks}).
\item Repository and literature searches were targeted, not exhaustive
(02: \cref{lwo:app:sources}; 06 and~07: their shipped audits).
\item The research questions are proposed continuations, not an authenticated
list of published open problems.
\item Source~06 does not classify all gaps of $\W_\kappa$ through its
missing-branch construction, does not describe the automorphism group, and
does not claim that full choice is necessary (its audit, Section~5).
\item No source gives an intrinsic criterion for an arbitrary linear order to
embed into $\W$; absence of $\kappa^+$ and its reverse is necessary and is not
claimed sufficient (02, 06 and~07 all say so).
\end{itemize}
Placement of this report in the Cardinals collection, beside the Lean
development of the Cardinals project, confers no formal status: no statement
of this report has a Lean or Rocq counterpart.

\section{Three orders that must be distinguished}\label{lwo:sec:three}

\subsection{The literal order of all well-orderings}
Throughout the real case, put
\[
                 \kappa=|\R|=2^{\aleph_0},
\]
and identify this cardinal with its initial ordinal. For an ordinal $\alpha$
of cardinality $\kappa$, let
\[
 \W_\alpha=\{e:\alpha\longrightarrow\R:e\text{ is a bijection}\},
 \qquad
 \W=\bigcup_{\kappa\leq\alpha<\kappa^+}\W_\alpha.
\]
A well-order relation on $\R$ determines exactly one member of $\W$: its
increasing enumeration in that well-order. Conversely, $e\in\W_\alpha$
defines the relation
\[
       r\mathrel{\prec_e}s\quad\Longleftrightarrow\quad e^{-1}(r)<e^{-1}(s).
\]
Thus $\W$ really represents all well-orderings, not merely their ordinal
isomorphism types. The order in which we compare their entries is the
\emph{usual} order of the real numbers, not one of the well-orders being
compared.

For distinct $e,f\in\W$, write
\begin{equation}\label{lwo:eq:lex}
 e<_{\lex}f
 \quad\Longleftrightarrow\quad
 e(\delta)<f(\delta),\qquad
 \delta=\min\{\xi<\min(\dom e,\dom f):e(\xi)\ne f(\xi)\}.
\end{equation}
The existence of $\delta$ is proved in \cref{lwo:lem:prefix-free}. In particular,
there is no choice of an ``end-of-word'' convention hidden in this definition.

One could instead sort by ordinal length first. That would be a different
order: choosing one enumeration of each possible length would then give an
increasing sequence of type $\kappa^+$. This is impossible for
\eqref{lwo:eq:lex}. Another different construction is to compare characteristic
functions of well-order relations using a chosen well-order of $\R^2$.
Neither alternative is used here.

\subsection{The fixed-length restriction}
The suborder $\W_\kappa$ contains only those well-orderings whose ordinal type
is the initial ordinal $\kappa$. It is a substantial suborder, but it is not
all of $\W$. For example, types $\kappa+1$, $\kappa+2$, and
$\kappa\cdot2$ also occur in $\W$.

This distinction cannot be dismissed on cardinality grounds. Both
$\W_\kappa$ and $\W$ have cardinality $2^\kappa$, yet $\W$ does not embed
into $\W_\kappa$. Indeed, no union of strata with a bounded set of ordinal
lengths can contain an embedded copy of $\W$.

\subsection{The index in the checked Kanovei--Shelah paper}
Kanovei and Shelah's 2004 paper constructs a definable, countably saturated
elementary extension of the reals. Its Section~1 uses the index
\begin{equation}\label{lwo:eq:KS}
 \A=\{a:\kappa\longrightarrow\mathcal P(\N):
                     \ran a\text{ is an ultrafilter on }\N\},
\end{equation}
ordered lexicographically; subsets of $\N$ are themselves compared by their
binary characteristic sequences \cite[Section~1]{KS}. These maps need not be
injective. The definition includes principal ultrafilters. For an ultrafilter
$U$, write $\A_U=\{a\in\A:\ran a=U\}$.

Thus the checked source uses fixed-length ultrafilter enumerations, not
literally $\W$. The two constructions share the principle of ordering
\emph{presentations} without first selecting a presentation. We study both
rather than identifying them.

\subsection{Connection to ProveIt and to existing terminology}
The ProveIt report \emph{Every ordinal as a point-separating game value}
contains a binary-cylinder argument and an order-embedding interpretation of
lexicographic coding \cite{ProveIt}. We use that line of thought, but give the
needed cube argument independently in \cref{lwo:sec:cubes}. None of that report's
claims about compact $T_1$ spaces is needed here.

Minimum ordinal lengths of binary lexicographic representations are also
studied under the name \emph{intrinsic length} in work of Knoblauch
\cite{Knoblauch}. We fix our precise convention in \cref{lwo:def:binary-length};
no classification theorem from that literature is assumed. The targeted source
search did not establish a prior full classification of the specific order
$\W$ considered here. This is a limitation of the bibliography, not a claim
that such a classification is absent from the literature.

\section{Main conclusions at a glance}\label{lwo:sec:summary}

Let $\B_\theta=(\two^\theta,<_{\lex})$, for every ordinal $\theta$.
An expression such as $\R^\gamma_{\lex}$ always has ordinal-indexed
coordinates, read from left to right. We use $L\equi M$ for equimorphism.

\begin{theorem}[Main embedding spectrum]\label{lwo:thm:main-spectrum}
The following statements hold.
\begin{enumerate}[label=(\roman*)]
\item $|\W|=|\W_\kappa|=|\A|=2^\kappa$. Every linear order of cardinality
at most $\kappa$ embeds into each of these orders.
\item For every ordinal $\beta$ and each
$L\in\{\W,\W_\kappa,\A\}$,
\[
                \beta\emb L\quad\Longleftrightarrow\quad\beta<\kappa^+.
\]
The same equivalence holds for reverse ordinals.
\item $\W_\kappa\equi\A\equi\A_U\equi\B_\kappa$ for every ultrafilter
$U$ on $\N$.
\item $\B_\theta\emb\W$ exactly when $\theta<\kappa^+$, whereas
$\B_\theta\emb\W_\kappa$ exactly when $\theta\leq\kappa$.
The latter criterion also holds with $\A$ in place of $\W_\kappa$.
\item $\R^\gamma_{\lex}\emb\W$ exactly when $\gamma<\kappa^+$.
It embeds into $\W_\kappa$ or $\A$ exactly when $\gamma\leq\kappa$.
\item $\W^\gamma_{\lex}\emb\W$ exactly when $\gamma<\kappa^+$.
For $0<\gamma<\kappa^+$ the two orders are equimorphic. In contrast,
$\W_\kappa\times_{\lex}\W_\kappa\not\emb\W_\kappa$, and
$\A\times_{\lex}\A\not\emb\A$.
\end{enumerate}
\end{theorem}

These assertions are proved in \cref{lwo:sec:ordinal,lwo:sec:fixed,lwo:sec:length,lwo:sec:absorption}.
\merge\ Items (i)--(iii) are also the opening theorems of sources~06
(Theorem~1.1, 06:102) and~07 (Theorem~2.1, 07:187); items (iv)--(vi) are
also in source~07's Theorem~2.1, in real-power form.
The exact refinement for every fixed type, proved in \cref{lwo:sec:all-strata}, is
\[
 \W_{\lambda+n}\equi\B_\lambda\times_{\lex}\{0,\ldots,n!-1\},
 \qquad |\lambda|=\kappa,\quad\lambda\text{ limit},\quad n<\omega.
\]
Moreover, two strata compare by their limit parts first and by their final
factorial block sizes second; see \cref{lwo:cor:strata-comparison}.
The structural distinction is summarized in \cref{lwo:tab:invariants}. Here
$2^{<\kappa}=\sup_{\lambda<\kappa}2^\lambda$, with cardinal exponents;
$\cell$ is the supremum of the sizes of pairwise disjoint families of
nonempty open sets. All topologies in the table are order topologies.

\begin{table}[ht]
\centering\small
\renewcommand{\arraystretch}{1.12}
\begin{tabularx}{\textwidth}{@{}Y*{3}{>{\centering\arraybackslash}p{.14\textwidth}}@{}}
\toprule
Invariant & $\W_\kappa$ & $\A$ & $\W$\\
\midrule
Cardinality & $2^\kappa$ & $2^\kappa$ & $2^\kappa$\\
Ordinals that embed & $\beta<\kappa^+$ & $\beta<\kappa^+$ & $\beta<\kappa^+$\\
Binary coding length $\bl$ & $\kappa$ & $\kappa$ & $\kappa^+$\\
Dense linear order & yes & yes & no\\
Least or greatest element & neither & neither & neither\\
Cofinality & $\omega$ & $\cf(\kappa)$ & $\omega$\\
Coinitiality & $\omega$ & $\omega$ & $\omega$\\
Density and cellularity & $2^{<\kappa}$ & $2^{<\kappa}$ & $2^\kappa$\\
Isolated points & none & none & $2^\kappa$\\
Self-dual & yes & no & yes\\
Unfilled $(\omega,\omega)$ cut & yes & yes & yes\\
Absorbs its lexicographic square & no & no & yes\\
\midrule
\multicolumn{4}{@{}l}{\merge\ \emph{Added from sources~06 and~07}}\\
Weight & $2^{<\kappa}$ & $2^{<\kappa}$ & $2^\kappa$\\
Real representation rank $\repr$ & $\kappa$ & $\kappa$ & $\kappa^+$\\
Point characters $(\chm,\chp)$ & $(\cf\kappa,\cf\kappa)$ & $(\cf\kappa,\cf\kappa)$ & \eqref{lwo:sp:eq:pointpairs}\\
Gap characters & \cref{lwo:rk:thm:Pgaps} & include $(\omega,\omega)$, $(\cf\kappa,\cf\kappa)$ & \cref{lwo:sp:thm:gapspectrum}\\
Uncountable balanced gaps & yes & yes & none\\
Isolated points dense & --- & --- & yes\\
Zero-dimensional & yes & yes & yes\\
Order-embeds into $\W$ with convex image & no & --- & yes\\
\bottomrule
\end{tabularx}
\caption{Equal ordinal spectra do not imply equal lexicographic complexity.
The assertions about $\A$ refer to the full index \eqref{lwo:eq:KS}.
\merge\ The rows below the rule are those of the summary tables of sources~06
(its Appendix~A) and~07 (its Section~2) not already above; the gap spectrum
of $\A$ is not classified.}
\label{lwo:tab:invariants}
\end{table}

\begin{table}[ht]
\centering\small
\renewcommand{\arraystretch}{1.15}
\begin{tabularx}{\textwidth}{@{}>{\raggedright\arraybackslash}p{.42\textwidth}Y@{}}
\toprule
Embedded or excluded object & Result\\
\midrule
Every linear order $L$ with $|L|\leq\kappa$ & embeds in every admissible slice $\W_\beta$ (\cref{lwo:sp:thm:universality})\\
$\B_\gamma$ and $\R^\gamma_{\lex}$, $\gamma<\kappa^+$ & embed in $\W$ (\cref{lwo:thm:full-length,lwo:cor:real-powers})\\
$\alpha$ or $\alpha^{\rev}$ & embeds exactly when $\alpha<\kappa^+$ (\cref{lwo:thm:ordinal-spectrum})\\
$S_{<\theta}$, $\theta<\kappa^+$ & embeds in $\W$ (\cref{lwo:sp:thm:surreals})\\
$S_{<\kappa^+}$ & does not embed, although its size is $2^\kappa$\\
All orders of size $2^\kappa$ & not universally embedded ($\kappa^+$ does not embed)\\
Full $\W$ into $\W_\kappa$ & excluded in ZFC (\cref{lwo:rk:cor:notP}); source~06 had it only if $2^{<\kappa}<2^\kappa$\\
$\W_\kappa$ into $\W$ with convex image & excluded (\cref{lwo:rk:cor:noconvex})\\
\bottomrule
\end{tabularx}
\caption{\merge\ Embedded and excluded objects; source~06's second summary
table (its Appendix~A), with its row on $\W$ into $\W_\kappa$ updated by the
in-batch answer and the last row added from source~07.}
\label{lwo:sp:tab:objects}
\end{table}

Under CH the ordinal cutoff is $\omega_2$: every ordinal below $\omega_2$
embeds, including $\omega_1\cdot2$ and $\omega_1^2$, but $\omega_2$ does not.
Without CH, the answer is still exactly $\kappa^+$, not an absolute cutoff at
$\omega_2$. For example, if $\mathfrak c=\aleph_2$, then $\omega_2$ embeds
and $\omega_3$ does not. These are conditional substitutions into the ZFC
theorem, not additional assumptions in its proof.

\section{Exhaustive enumeration orders}\label{lwo:sec:enumerations}

Several central facts hold for any infinite linearly ordered set $X$, not
just the real line. Write $\lambda=|X|$ and let $\W(X)$ consist of all
bijections from ordinals of cardinality $\lambda$ onto $X$, with the same
first-difference comparison.

\begin{lemma}[Prefix-freeness]\label[lemma]{lwo:lem:prefix-free}
No member of $\W(X)$ is a proper initial segment of another. Consequently
two distinct members have a first disagreement in the common part of their
domains, and first-difference comparison is a linear order.
\end{lemma}
\begin{proof}
Suppose $e:\alpha\to X$ is an initial segment of the injective map
$f:\beta\to X$, where $\alpha<\beta$. Since $e$ is onto, the value
$f(\alpha)$ already occurs at some coordinate below $\alpha$, contrary to
injectivity. If two enumerations agree throughout their common domain, one
is therefore equal to the other.

For transitivity, suppose $e<f<g$, with first disagreements $\delta$ and
$\varepsilon$. Before $\min(\delta,\varepsilon)$ all three agree. If the
coordinates are unequal, the earlier comparison gives $e<g$ there. If they
are equal, $e(\delta)<f(\delta)<g(\delta)$ gives the same conclusion.
Trichotomy follows from the linear order on $X$.
\end{proof}
\merge\ Also source~06, Lemma~2.1 and Proposition~2.3 (06:171, 06:198),
and source~07, Section~1.2 (07:132--141), by the same argument.

\begin{definition}
For an injective ordinal sequence $s$ of elements of $X$, its cylinder is
\[
                  C_s=\{e\in\W(X):s\subseteq e\}.
\]
We consider only prefixes that have an exhaustive extension. An exhaustive
prefix has a singleton cylinder.
\end{definition}

\begin{lemma}[Residual-set geometry]\label[lemma]{lwo:lem:cylinder}
Every cylinder is convex. If $s$ is not exhaustive and
$Y=X\setminus\ran s$, then concatenation gives an order isomorphism
\[
                       \W(Y)\longrightarrow C_s,
                       \qquad t\longmapsto s\cat t.
\]
For finite $Y$ of size $m$, the cylinder has exactly $m!$ points.
\end{lemma}
\begin{proof}
A point between two extensions of $s$ cannot first disagree with $s$ before
its end: that disagreement would put it on the same strict side of both
extensions. The remaining assertions follow because every exhaustive
extension must enumerate precisely $Y$ after $s$, and first differences in
the tail are unchanged by concatenation.
\end{proof}

\merge\ Sources~06 (Lemma~3.2, 06:246) and~07 (Lemma~3.1, 07:259) prove the
same lemma and add its next-value decomposition: if $Y_s=X\setminus\ran s$ is
nonempty, then
\begin{equation}\label{lwo:sp:eq:children}
 C_s\cong\sum_{a\in Y_s}^{<_X}\W(Y_s\setminus\{a\}),
\end{equation}
a lexicographic sum over $Y_s$ in its given order, whose summand indexed by
$a$ is the child cylinder $C_{s\cat\langle a\rangle}$, with $\W(\varnothing)$
a one-point order. Sorting the members of $C_s$ by their next entry proves it.
Source~06 also observes that concatenation uses ordinal addition, so no
ordinal subtraction convention is needed (06:242--244).

In particular, a long prefix may have already used all but two or three
reals. This is why the density of the usual real order does \emph{not}
imply density of $\W$.

\begin{proposition}[Cardinality]\label[proposition]{lwo:prop:cardinality}
For every infinite linearly ordered set $X$ of cardinality $\lambda$,
\[
              |\W(X)|=|\W_\lambda(X)|=2^\lambda.
\]
\end{proposition}
\begin{proof}
A well-order relation is a subset of $X^2$, giving the upper bound
$2^{|X^2|}=2^\lambda$. For the lower bound, partition $X$ into
$\lambda$ two-element sets $\{a_\xi,b_\xi\}$, with $a_\xi<b_\xi$.
For $x\in\two^\lambda$, enumerate the pair at coordinate $\xi$ as
$(a_\xi,b_\xi)$ if $x(\xi)=0$, and as $(b_\xi,a_\xi)$ otherwise.
The concatenation has length $2\cdot\lambda=\lambda$ and enumerates all of
$X$. Different bit sequences give different enumerations. In fact, this is
an order embedding $\B_\lambda\emb\W_\lambda(X)$.
\end{proof}
\merge\ Also source~06, Lemma~4.1 and Corollary~4.2 (06:331, 06:363), and
source~07, Lemma~4.1 and Corollary~4.6 (07:329, 07:417), by the same pair-swap
code. Source~06 notes that the multiplication in $2\cdot\nu=\nu$ is ordinal
multiplication, the two positions of a block being the less significant
coordinates (06:358--361). Two definitions that are \emph{not} used are ruled
out by source~06 (06:220--230): comparing characteristic functions of
relations at the ``least real where they differ'' fails, since a set of real
disagreement coordinates need not have a least member; and the order does not
descend to isomorphism types, since two enumerations of any two admissible
types can be made to compare either way by their first values.

\section{Binary cubes and the cylinder obstruction}\label{lwo:sec:cubes}

\begin{lemma}[Cut coding]\label[lemma]{lwo:lem:cut-code}
Every linear order $L$ of cardinality at most an infinite cardinal $\lambda$
embeds into $\B_\lambda$.
\end{lemma}
\begin{proof}
Choose a well-ordered list $(t_\xi)_{\xi<\mu}$ of the elements of $L$, where
$\mu\leq\lambda$. For $x\in L$, put
\[
 c_x(\xi)=\begin{cases}1&t_\xi\leq_Lx,\\0&t_\xi>_Lx,\end{cases}
 \qquad(\xi<\mu),
\]
and put $c_x(\xi)=0$ for $\mu\leq\xi<\lambda$. If $x<y$, then
$c_x\leq c_y$ coordinatewise, and the coordinate listing $y$ distinguishes
them. Their first unequal bits are consequently $0$ and $1$.
\end{proof}
\merge\ Also source~06, Lemma~4.3 (06:373), and source~07, Lemma~4.2
(07:339), coding $x$ by the characteristic function of $\{t:t<x\}$ instead of
$\{t:t\le x\}$; the argument is the same. Source~06 attributes this coding to
the standard literature (Moore's representation of Aronszajn lines,
\cite[Section~2]{Moore}) and stresses that it uses a well-ordering of the
coordinate set, not one compatible with $L$ (06:385--387).

\begin{lemma}[A cut preserves a full child]\label[lemma]{lwo:lem:child}
Let $C_s$ be a full prefix cylinder in $\B_\theta$, with
$\dom s<\theta$. If $D$ is an initial segment of $C_s$, then either
$C_{s\cat0}\subseteq D$ or $C_{s\cat1}\subseteq C_s\setminus D$.
The same assertion holds for cylinders $C_s\times F$ in
$\B_\theta\times_{\lex}F$, where $F$ is any nonempty linear order.
\end{lemma}
\begin{proof}
If the first inclusion fails, choose a point of the lower child outside
$D$. Every point of the upper child lies above it, hence is also outside
the initial segment. The additional terminal factor does not change the
ordering of the two children.
\end{proof}

\merge\ The next theorem is Corollary \texttt{cor:strict} of the ProveIt
report \emph{Every ordinal as a point-separating game value}
(\texttt{ordinal\_separation.tex}, lines 633--641), where it is proved through
the point-separating game: an embedding $\B_\alpha\emb\B_\beta$ with
$\beta<\alpha$ would give a winning strategy of length $\beta$, contrary to
that report's exact game value. Source~02's cylinder-fusion proof below is a
second route, which avoids the game.

\begin{theorem}[Strict hierarchy of binary cubes]\label{lwo:thm:cube-strict}
For all ordinals $\alpha,\beta$,
\[
                         \B_\alpha\emb\B_\beta
                         \quad\Longleftrightarrow\quad\alpha\leq\beta.
\]
\end{theorem}
\begin{proof}
Padding with zero bits proves the forward-length direction. Suppose instead
that $f:\B_\alpha\emb\B_\beta$ and $\beta<\alpha$. Construct a source
prefix $s_\xi$ of length $\xi$, for $\xi\leq\beta$, so that $f$ is constant
on its first $\xi$ target coordinates throughout $C_{s_\xi}$.

At a successor step, all previous target bits are fixed. The inverse image
of target bit $0$ at coordinate $\xi$, restricted to the current cylinder,
is an initial segment. By \cref{lwo:lem:child}, one source child lies entirely
on one side of this cut. Retain that child. At a limit stage take the union
of the retained prefixes. Its full cylinder is the intersection of the
earlier cylinders and remains nonempty. This is a fusion argument, not an
appeal to finite-intersection compactness.

After $\beta$ steps a source cylinder of length $\beta<\alpha$ still
contains at least two points, while their images agree on all $\beta$
coordinates. This contradicts injectivity.
\end{proof}

\begin{definition}[Binary lexicographic length]\label[definition]{lwo:def:binary-length}
For a set-sized linear order $L$, define
\[
               \bl(L)=\min\{\theta:L\emb\B_\theta\}.
\]
This minimum exists in ZFC by \cref{lwo:lem:cut-code}. A singleton and the empty order have length
zero. Equimorphic orders have the same binary length.
\end{definition}

\begin{lemma}[Real and binary coordinates]\label[lemma]{lwo:lem:real-code}
$\R\equi\B_\omega$. More generally,
\[
                 \R^\gamma_{\lex}\equi\B_{\omega\cdot\gamma}
\]
for every ordinal $\gamma$.
\end{lemma}
\begin{proof}
List the rationals as $(q_n)_{n<\omega}$ and send $r$ to
$j(r)(n)=1$ exactly when $q_n<r$. This is an order embedding by the same
coordinatewise argument as in \cref{lwo:lem:cut-code}. In the other direction,
send a bit sequence $x$ to
\[
                     T(x)=\sum_{n<\omega}\frac{2x(n)}{3^{n+1}}.
\]
At the first unequal digit the contribution dominates the sum of all later
possible differences, so this map is strictly increasing. Applying these
maps coordinatewise and concatenating blocks of length $\omega$ proves the
last statement. The ordinal length is $\omega\cdot\gamma$, not cardinal
multiplication and not $\gamma\cdot\omega$.
\end{proof}

For the uncountable initial ordinal $\kappa=2^{\aleph_0}$,
\begin{equation}\label{lwo:eq:omega-kappa}
                             \omega\cdot\kappa=\kappa.
\end{equation}
Indeed, each $\omega\cdot\xi$ for $\xi<\kappa$ has cardinality less than
$\kappa$, and their supremum is at least $\kappa$ and at most $\kappa$.
No regularity of $\kappa$ is assumed.

\merge\ The codes $j$ and $T$ of \cref{lwo:lem:real-code} are also source~07's
(07:436--458) and source~06's (06:433--445), and \eqref{lwo:eq:omega-kappa} is
part of source~07's Lemma~3.2 (07:288), which also records $\kappa^\kappa=2^\kappa$,
$\sup_{\lambda<\kappa}\kappa^\lambda=2^{<\kappa}$ and $\cf\kappa>\omega$.

\begin{corollary}[Strictness of real powers; Kuhlmann]\label[corollary]{lwo:rk:thm:strict}
For all ordinals $\alpha,\beta$,
\[
 \R^\alpha_{\lex}\emb\R^\beta_{\lex}\quad\Longleftrightarrow\quad\alpha\leq\beta.
\]
\end{corollary}
\begin{proof}[First proof \textup{(from source~02's lemmas)}]
By \cref{lwo:lem:real-code}, $\R^\alpha_{\lex}\equi\B_{\omega\cdot\alpha}$
and $\R^\beta_{\lex}\equi\B_{\omega\cdot\beta}$. By
\cref{lwo:thm:cube-strict}, an embedding exists exactly when
$\omega\cdot\alpha\le\omega\cdot\beta$, that is, when $\alpha\le\beta$, since
left multiplication by $\omega$ is strictly increasing on ordinals.
\end{proof}
\begin{proof}[Second route \textup{(source~07, Theorem~5.2, 07:477)}]
Extension by zero gives the forward construction when $\alpha\leq\beta$.
For the converse, suppose $\alpha>\beta$ and $F:\R^\alpha_{\lex}\emb\R^\beta_{\lex}$.
Recursively choose a source prefix $s_\delta$ of length $\delta$ and a target
prefix $t_\delta$ of the same length, for $\delta\leq\beta$, such that every
image of a source extension of $s_\delta$ extends $t_\delta$.

At step $\delta<\beta$, subdivide the source cylinder according to its next
real value $r$. Images of these subcylinders are ordered by $r$. Their
$\delta$-th target-coordinate sets are consequently ordered subsets of $\R$.
At most countably many such sets contain two distinct values: for each such
set choose a rational strictly between two of its values, and these rationals
must be distinct. Choose an $r$ outside the countable exceptional set.
The coordinate is constant on this subcylinder, so it determines
$t_{\delta+1}$.

At limits take unions. A source prefix of length at most $\beta$ always has
extensions, since $\beta<\alpha$. At the final step all extensions of
$s_\beta$ have exactly the same image, namely $t_\beta$. There are at least
two such source extensions. This contradicts injectivity.
\end{proof}
Source~07 identifies this ordinal-exponent theorem as classical: it is the
ordinal case of Kuhlmann's embedding result \cite[Corollary~2.4]{Kuhlmann}, and
source~07's coordinate fusion is its own proof of the case needed. The
point-separating report has only the binary case.

\section{The exact ordinal spectrum}\label{lwo:sec:ordinal}

The upper bound is the essential point. Cardinality alone would permit
well-ordered suborders far larger than $\kappa$.

\begin{lemma}[Eventual constancy]\label[lemma]{lwo:lem:eventual}
Let $\lambda$ be an infinite cardinal. A nondecreasing sequence of length
$\lambda^+$ with values in a linear order of cardinality at most $\lambda$
is eventually constant. The same holds for a nonincreasing sequence.
\end{lemma}
\begin{proof}
There are at most $\lambda$ distinct values. Take the first occurrence of
each value. Their supremum, increased by one if necessary, is below
$\lambda^+$: the supremum of at most $\lambda$ ordinals below $\lambda^+$
has cardinality at most $\lambda$. After this bound no new value can first
occur. Monotonicity prevents a change to a value that already occurred
before the current value. Hence the sequence is constant on a tail.
\end{proof}
\merge\ Also source~06, Lemma~5.1 (06:454), for a regular $\lambda$ and a
target of cardinality below $\lambda$ (``some fiber is unbounded, and a
monotone unbounded fiber contains a tail''), and source~07, Lemma~4.3 (07:355).

\begin{theorem}[Exhaustion forbids a successor-cardinal chain]\label{lwo:thm:no-long}
For every infinite linearly ordered set $X$ of cardinality $\lambda$,
$\W(X)$ contains neither $\lambda^+$ nor its reverse.
\end{theorem}
\begin{proof}
Suppose $(e_\eta)_{\eta<\lambda^+}$ is strictly increasing. Recursively, for
each $\xi<\lambda^+$, we shall obtain an injective prefix $s_\xi$ of length
$\xi$ shared by every $e_\eta$ on a tail of the index set.

The empty prefix starts the recursion. At a limit $\xi<\lambda^+$, the
supremum of the preceding tail bounds is still below $\lambda^+$, since
$|\xi|\leq\lambda$. The union of the prefixes is injective and is shared
by the resulting common tail.

Suppose the shared prefix $s_\xi$ is exhaustive. Prefix-freeness forces all
members of this tail to equal $s_\xi$, contrary to strict increase. Thus it
is not exhaustive. Every enumeration on the tail has a value at coordinate
$\xi$, and these values are nondecreasing because the earlier coordinates
are fixed. By \cref{lwo:lem:eventual}, they are constant on a further tail.
Their common value is not in $\ran s_\xi$, by injectivity. Adjoin it to get
$s_{\xi+1}$.

If the recursion never stops, its union is an injection
$\lambda^+\to X$, a contradiction. If it stops by exhaustion, we already
have a contradiction. Replacing nondecreasing by nonincreasing proves the
reverse-order assertion.
\end{proof}
\merge\ Also source~06, Theorem~5.2 (06:465), and source~07, Theorem~4.4
(07:367), by the same stabilization of a common prefix. Source~06 remarks
(its Remark~5.3) that the regularity used at limit stages is a ZFC fact about
the successor cardinal $\nu^+$, not an assumption about $\nu$: a union of at
most $\nu$ ordinals below $\nu^+$ has cardinality at most $\nu\cdot\nu=\nu$;
the continuum $\kappa$ itself need not be regular. Source~07 notes that no
separability of the alphabet is needed.

\begin{remark}
It would be incorrect to argue that the sequence of enumeration lengths is
bounded below $\lambda^+$. A long family can have unbounded lengths. The
proof instead stabilizes a common prefix and uses the fact that an exhaustive
prefix cannot be extended. This is exactly the extra step needed for the
full order, rather than a fixed-length product.
\end{remark}

\begin{theorem}[Universal small orders and exact ordinals]\label{lwo:thm:ordinal-spectrum}
For every infinite linearly ordered set $X$ of cardinality $\lambda$:
\begin{enumerate}[label=(\alph*)]
\item Every linear order of cardinality at most $\lambda$ embeds into
$\W_\lambda(X)$ and hence into $\W(X)$.
\item For every ordinal $\beta$,
\[
       \beta\emb\W(X)\quad\Longleftrightarrow\quad\beta<\lambda^+.
\]
The same is true for reverse ordinals and for $\W_\lambda(X)$.
\end{enumerate}
\end{theorem}
\begin{proof}
Compose cut coding with the pair-swap embedding in
\cref{lwo:prop:cardinality}. An ordinal has cardinality at most $\lambda$
exactly when it is below $\lambda^+$. Larger ordinals contain
$\lambda^+$ as an initial segment and are excluded by \cref{lwo:thm:no-long}.
Reverse orders are covered by the same small-order universality and the
reverse half of that theorem.
\end{proof}

For the real line, this proves the precise ordinal answer in
\cref{lwo:thm:main-spectrum}. Notice the distinction between the external
lexicographic order on well-orderings and the well-order carried by any one
of its elements. The external order contains long ascending and descending
chains even though every element individually codes a well-order.

\merge\ Also source~06, Corollary~5.4 (06:501), which states the spectrum for
every slice $\W_\beta(X)$ of admissible length and notes that there is no
largest embedded ordinal, and source~07, Corollary~4.5 (07:394). Source~07
adds (07:408--415) that $\kappa+1$, $\kappa\cdot2$ and $\kappa\cdot\kappa$
all embed while $\kappa^+$ does not, and that the map sending a well-order to
its order type is onto $[\kappa,\kappa^+)$ but not monotone, so choosing one
well-order of each type cannot give an increasing $\kappa^+$-sequence.
Source~06 adds (its Example~4.5) that, without CH, $\W$ already contains
$\omega_1$, its reverse, every linear order on $\omega_1$, every ordinal of
cardinality $\con$ and every continuum-sized ordered field or linear order,
as orders, without preserving field operations; and (06:530--531) that $\W$ is
not universal for orders of size $2^\kappa$, since $\kappa^+\le2^\kappa$.

\section{Fixed length and the actual ultrafilter index}\label{lwo:sec:fixed}

\begin{theorem}[Fixed-length equivalences]\label{lwo:thm:fixed-equivalence}
For every ultrafilter $U$ on $\N$,
\[
                 \W_\kappa\equi\B_\kappa\equi\A_U\equi\A.
\]
Consequently, for any set-sized linear order $L$,
\begin{equation}\label{lwo:eq:fixed-classification}
 L\emb\W_\kappa
 \quad\Longleftrightarrow\quad L\emb\A
 \quad\Longleftrightarrow\quad\bl(L)\leq\kappa.
\end{equation}
\end{theorem}
\begin{proof}
The pair-swap construction gives $\B_\kappa\emb\W_\kappa$. In the other
direction, $\W_\kappa$ is a suborder of $\R^\kappa_{\lex}$, which embeds
into $\B_{\omega\cdot\kappa}=\B_\kappa$ by \cref{lwo:lem:real-code}.

Every ultrafilter $U$ has cardinality $\kappa$. To see this, partition
$\mathcal P(\N)$ into complementary pairs $\{S,\N\setminus S\}$.
There are $\kappa$ such pairs, and an ultrafilter contains exactly one
member of each pair. Choose a surjection $u:\kappa\to U$ and two distinct
members $u_0<u_1$ of $U$, in binary lexicographic order. For
$x\in\two^\kappa$, set
\begin{equation}\label{lwo:eq:ultrafilter-code}
 a_x(2\cdot\xi)=u_{x(\xi)},\qquad
 a_x(2\cdot\xi+1)=u(\xi)\qquad(\xi<\kappa).
\end{equation}
Since $2\cdot\kappa=\kappa$, these are maps of the required length. Their
ranges are exactly $U$, and the first changed coding bit gives the first
changed value. Thus $\B_\kappa\emb\A_U$.

Finally, $\A_U\subseteq\A\subseteq(\mathcal P(\N))^\kappa$.
Flattening the characteristic sequences of subsets of $\N$ turns the last
order into $\B_{\omega\cdot\kappa}=\B_\kappa$. Composing these embeddings
proves all the equivalences and \eqref{lwo:eq:fixed-classification}.
\end{proof}

The fact that even one fixed fiber $\A_U$ has this universality is an
order-theoretic statement; it does not eliminate the role of all ultrafilters
in the definable model construction. In particular, the embeddings in
\eqref{lwo:eq:ultrafilter-code} use a chosen enumeration of $U$.

\merge\ Also source~06, Proposition~4.6 (06:425), and source~07,
\eqref{lwo:eq:omega-kappa} with its display (5.3) (07:456) and Theorem~11.2
(07:1129), where the ultrafilter coding is the same and the cardinality of an
ultrafilter is its Lemma~11.1. Source~07 states the equimorphism with real
coordinates as well: $\W_\kappa\equi\A\equi\B_\kappa\equi\R^\kappa_{\lex}$.
Both stress that equimorphism is not isomorphism: $\B_\kappa$ has endpoints,
whereas $\W_\kappa$ has none (06:447--448), and the characteristic-function
order on $\mathcal P(\N)$ has endpoints and adjacent pairs, unlike the usual
order of $\R$, so a bijection between these alphabets is not an order
isomorphism (source~07, Remark~1.2).

\begin{corollary}\label[corollary]{lwo:cor:fixed-spectra}
The binary lengths of $\W_\kappa$, $\A$, and $\A_U$ are all $\kappa$.
Their embedded ordinals and reverse ordinals are exactly those below
$\kappa^+$. Moreover,
\[
 \B_\theta\emb\W_\kappa\ \Longleftrightarrow\ \theta\leq\kappa,
 \qquad
 \R^\gamma_{\lex}\emb\W_\kappa\ \Longleftrightarrow\ \gamma\leq\kappa,
\]
with the same criteria for $\A$ and $\A_U$.
\end{corollary}
\begin{proof}
Use \cref{lwo:thm:cube-strict,lwo:thm:fixed-equivalence}. For real powers,
$\omega\cdot\gamma\leq\kappa$ is equivalent to $\gamma\leq\kappa$:
\eqref{lwo:eq:omega-kappa} proves one direction, while
$\gamma\geq\kappa+1$ implies
$\omega\cdot\gamma\geq\kappa+\omega>\kappa$. The ordinal spectrum follows
from equimorphism with $\W_\kappa$ and \cref{lwo:thm:ordinal-spectrum}.
\end{proof}

\subsection{An exact finite-tail refinement}
The last entry of an exhaustive enumeration is forced by all preceding
entries. Several final entries, however, retain a genuine finite amount of
lexicographic information. This gives a more precise hierarchy than the
cardinality or ordinal spectrum alone detects.

For a positive integer $m$, write
\[
                  D_{\kappa,m}=\B_\kappa\times_{\lex}\{0,\ldots,m-1\}.
\]
The binary sequence is the first, more significant coordinate.

\begin{theorem}[Finite-tail equivalence]\label{lwo:thm:finite-tail}
For every nonnegative integer $n$,
\begin{equation}\label{lwo:eq:finite-tail}
                         \W_{\kappa+n}\equi D_{\kappa,n!}.
\end{equation}
In particular,
\[
 \W_{\kappa+1}\equi\W_\kappa,
 \qquad
 \W_{\kappa+2}\equi\B_{\kappa+1}.
\]
\end{theorem}
\begin{proof}
Choose a fixed $n$-element set $F\subseteq\R$. Partition $\R\setminus F$
into $\kappa$ ordered pairs, and use the pair-swap construction for the
first $\kappa$ entries. Append the $r$-th permutation of $F$ in lexicographic
order, for $0\leq r<n!$. This embeds $D_{\kappa,n!}$ into
$\W_{\kappa+n}$. The case $n=0$ uses the unique empty permutation.

Conversely, code the first $\kappa$ entries of an enumeration by a member of
$\B_\kappa$, using rational-cut codes and \eqref{lwo:eq:omega-kappa}. When two
such prefixes agree, the remaining set of $n$ reals is the same. Record the
rank, between $0$ and $n!-1$, of the final permutation of those reals. The
prefix code followed by this rank is an order embedding into
$D_{\kappa,n!}$.
\end{proof}

\begin{theorem}[Finite-fiber capacity]\label{lwo:thm:finite-capacity}
For every positive finite $m$,
\begin{align}
 \bl(D_{\kappa,m})&=\kappa+\lceil\log_2m\rceil,\label{lwo:eq:finite-length}\\
 \B_\theta\emb D_{\kappa,m}
       &\quad\Longleftrightarrow\quad
             \theta\leq\kappa+\lfloor\log_2m\rfloor.
             \label{lwo:eq:finite-cubes}
\end{align}
Hence the same formulas hold for $\W_{\kappa+n}$ with $m=n!$.
\end{theorem}
\begin{proof}
A finite increasing binary code for $\{0,\ldots,m-1\}$ proves the upper
bound in \eqref{lwo:eq:finite-length}. Suppose
$D_{\kappa,m}\emb\B_{\kappa+j}$ for a finite $j$. Read the first $\kappa$
target bits and apply \cref{lwo:lem:child} successively to source cylinders
$C_s\times\{0,\ldots,m-1\}$. After $\kappa$ steps, all $m$ points over one
source bit sequence have the same first $\kappa$ target bits. They must fit
into a fiber of size $2^j$. Thus $m\leq2^j$. Targets of length below
$\kappa$ are already excluded by their failure to contain $\B_\kappa$.
This proves \eqref{lwo:eq:finite-length}.

For \eqref{lwo:eq:finite-cubes}, $2^j\leq m$ gives the direct embedding
$\B_{\kappa+j}=\B_\kappa\times_{\lex}\two^j\emb D_{\kappa,m}$.
Conversely, given an embedding in this direction, fuse source cylinders
while reading the first $\kappa$ bits of the target's $\B_\kappa$
coordinate. The final source cylinder contains $2^j$ points, but its target
fiber has only $m$ points. Hence $2^j\leq m$. Any ordinal length greater
than $\kappa+\lfloor\log_2m\rfloor$ contains the next forbidden finite
extension as an initial coordinate block, so all larger lengths are also
excluded.
\end{proof}

\begin{example}
For $n=3$ there are six final permutations. Thus
\[
 \bl(\W_{\kappa+3})=\kappa+3,
 \qquad
 \B_\theta\emb\W_{\kappa+3}
       \Longleftrightarrow\theta\leq\kappa+2.
\]
The least coding length need not itself be a cube that embeds back. This
small finite-fiber example foreshadows the much larger discrepancy for
$\W$: its coding length is $\kappa^+$, but $\B_{\kappa^+}$ does not embed.
\end{example}

\section{Unbounded lengths and full lexicographic complexity}\label{lwo:sec:length}

\begin{lemma}[Arbitrarily long small-cardinal cubes]\label[lemma]{lwo:lem:all-small-cubes}
Let $X$ be an infinite linearly ordered set of cardinality $\lambda$. For
every ordinal $\theta<\lambda^+$,
\[
                              \B_\theta\emb\W(X).
\]
\end{lemma}
\begin{proof}
Choose pairwise disjoint pairs $\{a_\xi,b_\xi\}$ for $\xi<\theta$, ordered
so that $a_\xi<b_\xi$, leaving a residual set of cardinality $\lambda$.
This is possible because $|\theta|\leq\lambda$ and $\lambda$ is infinite.
For each bit sequence, list each pair in its selected orientation, in the
ordinal order of $\theta$, and then append one fixed enumeration of the
residual set. The total length is $2\cdot\theta+\lambda$, an ordinal of
cardinality $\lambda$. First changed bits give first changed entries.
\end{proof}
\merge\ Also the second half of source~06's Lemma~4.1 (06:351) and, with
real coordinates, source~07's Proposition~5.1 (07:460).

\begin{theorem}[The full coding length]\label{lwo:thm:full-length}
For every infinite linearly ordered set $X$ of cardinality $\lambda$,
\begin{equation}\label{lwo:eq:full-length}
 \bl(\W(X))=\lambda^+,
 \qquad
 \B_\theta\emb\W(X)\ \Longleftrightarrow\ \theta<\lambda^+.
\end{equation}
In particular, $\W$ embeds into $\B_{\kappa^+}$, but the reverse embedding
does not exist.
\end{theorem}
\begin{proof}
Embed $X$ into $\B_\lambda$ by \cref{lwo:lem:cut-code}. For an enumeration of
length $\alpha<\lambda^+$, concatenate its coordinate codes into a binary
sequence of length $\lambda\cdot\alpha$, and pad with zeros to length
$\lambda^+$. Since $|\lambda\cdot\alpha|\leq\lambda$, this is a legitimate
shorter initial block. Prefix-freeness ensures that two different
exhaustive enumerations disagree before either ends, so padding introduces
no ambiguity. This proves the upper bound $\bl(\W(X))\leq\lambda^+$.

If $\W(X)$ embedded into $\B_\delta$ for $\delta<\lambda^+$, then
\cref{lwo:lem:all-small-cubes} would give
$\B_{\delta+1}\emb\B_\delta$, contrary to \cref{lwo:thm:cube-strict}. This
proves equality.

The positive cube embeddings are \cref{lwo:lem:all-small-cubes}. To exclude
$\B_{\lambda^+}$, note that this cube contains a copy of $\lambda^+$:
map $\eta<\lambda^+$ to the sequence with value $1$ below $\eta$ and $0$
from $\eta$ onward. This would contradict \cref{lwo:thm:no-long}. Larger cubes
contain this forbidden cube.
\end{proof}
\merge\ Source~07 proves the real-coordinate counterpart,
$\repr(\W)=\kappa^+$ and $\repr(\W_\kappa)=\repr(\A)=\kappa$
(\cref{lwo:rk:thm:ranks}; its Theorem~5.3, 07:505), from \cref{lwo:rk:thm:strict} by the same padding
argument; the two invariants are compared in \cref{lwo:tab:notation}.

For the real case, the upper coding construction can use countable
rational-cut blocks rather than blocks of length $\kappa$. Thus an
enumeration of length $\alpha$ is first coded at length
$\omega\cdot\alpha$. The ultimate common padding length is still
$\kappa^+$.

\begin{corollary}[No bounded-length universal suborder]\label[corollary]{lwo:cor:unbounded}
Fix $\tau<\kappa^+$. Then
\[
                      \W\not\emb\bigcup_{\alpha\leq\tau}\W_\alpha.
\]
Consequently, the image of every order embedding $\W\emb\W$ has an
unbounded set of enumeration lengths in $\kappa^+$.
\end{corollary}
\begin{proof}
The bounded union embeds into $\B_{\omega\cdot(\tau+1)}$ by coding and
padding. Its target length is below $\kappa^+$, whereas $\bl(\W)=\kappa^+$.
\end{proof}

\begin{corollary}[No fixed power below $\kappa^+$; sources~02 and~07]
\label[corollary]{lwo:rk:cor:notP}
$\W$ embeds neither into $\W_\kappa$, nor into $\A$, nor into any real power
$\R^\gamma_{\lex}$ or binary cube $\B_\gamma$ with $\gamma<\kappa^+$. This
holds in ZFC, without any hypothesis on $2^{<\kappa}$.
\end{corollary}
\begin{proof}[First proof \textup{(source~02)}]
$\W_\kappa$ is the stratum $\alpha=\kappa$, so \cref{lwo:cor:unbounded} with
$\tau=\kappa$ excludes it; $\A\equi\W_\kappa$ by
\cref{lwo:thm:fixed-equivalence}. An embedding into $\R^\gamma_{\lex}$ would
give $\W\emb\B_{\omega\cdot\gamma}$ by \cref{lwo:lem:real-code}, with
$\omega\cdot\gamma<\kappa^+$, contrary to $\bl(\W)=\kappa^+$
(\cref{lwo:thm:full-length}).
\end{proof}
\begin{proof}[Second route \textup{(source~07, Corollary~5.4, 07:525)}]
By real representation rank: $\repr(\W)=\kappa^+$, whereas
$\repr(\W_\kappa)=\repr(\A)=\kappa$ and $\repr(\R^\gamma_{\lex})=\gamma$ by
\cref{lwo:rk:thm:strict}.
\end{proof}
\merge\ This answers, unconditionally, source~06's Research question~14.1
(06:1332), which asked whether $\W$ embeds into $\W_\kappa$ or into one fixed
power $\R^\gamma_{\lex}$ with $\gamma<\kappa^+$ ``in any model of ZFC''. Source~06
had the negative answer for $\W_\kappa$ only under $2^{<\kappa}<2^\kappa$; its
argument is kept as a different route in the remark below. Its question about
the longer fixed slices $\W_\beta$ is answered as well: $\W$ embeds into no
bounded union of strata (\cref{lwo:cor:unbounded}).

\begin{remark}[The cellularity route; source~06, Corollary~9.5]
\label[remark]{lwo:sp:rem:conditional}
If $2^{<\kappa}<2^\kappa$, the separation $\W\not\emb\W_\kappa$ also follows from
cellularity. The full order has $2^\kappa$ pairwise disjoint six-element convex
blocks (\cref{lwo:prop:full-topology}). An order embedding into $\W_\kappa$ would
send, from each block, three selected points to an ordered triple. The
intervals between the first and third images would be nonempty and pairwise
disjoint. This contradicts $\cell(\W_\kappa)=2^{<\kappa}$
(\cref{lwo:prop:fixed-topology}). The proof uses intervals between images,
not continuity of the embedding. When $2^{<\kappa}=2^\kappa$ this
cardinal-invariant argument says nothing; source~06 left that case open, and
\cref{lwo:rk:cor:notP} closes it.
\end{remark}

\begin{corollary}[Real-coordinate powers]\label[corollary]{lwo:cor:real-powers}
For every ordinal $\gamma$,
\[
                        \R^\gamma_{\lex}\emb\W
                        \quad\Longleftrightarrow\quad\gamma<\kappa^+.
\]
\end{corollary}
\begin{proof}
For $\gamma<\kappa^+$, the ordinal $\omega\cdot\gamma$ also lies below
$\kappa^+$, so use \cref{lwo:lem:real-code,lwo:lem:all-small-cubes}. For
$\gamma\geq\kappa^+$, restricting all coordinates to two reals gives a
copy of $\B_{\kappa^+}$ and therefore of $\kappa^+$.
\end{proof}
\merge\ Also source~07, Proposition~5.1 and the last assertion of
Corollary~5.4 (07:460, 07:525), and, for the positive part, source~06,
Proposition~4.6 (06:425).

These results give an exact classification for binary and real-coordinate
lexicographic powers. They do not classify every linear order of cardinality
$2^\kappa$ that has no $\kappa^+$-chain in either direction. In particular,
absence of the two long monotone chains is not asserted to be a sufficient
embedding criterion for $\W$.


\section{An exact classification of every fixed ordinal type}\label{lwo:sec:all-strata}

The finite-tail result extends to \emph{every} possible ordinal length, not
only to $\kappa+n$. The key is to exploit the shrinking alphabet after a
long prefix has already enumerated most of the ground set. Coding every
remaining real by a fresh block of $\omega$ bits would miss this improvement.

\begin{lemma}[Countable alphabets]\label[lemma]{lwo:lem:countable-alphabet}
If $X$ is a countable linear order, then $X^\omega_{\lex}$ embeds into
$\B_\omega$. Consequently, for every nonzero limit ordinal $\lambda$,
$X^\lambda_{\lex}$ embeds into $\B_\lambda$.
\end{lemma}
\begin{proof}
There are countably many nonempty finite words $t$ over $X$. For a word of
length $m$, the set
\[
                 D_t=\{x\in X^\omega:x\restr m\leq_{\lex}t\}
\]
is an initial segment. Enumerate these sets and code a point by bit $0$
when it belongs to the corresponding $D_t$, and bit $1$ otherwise. If
$x<y$ first differ at coordinate $m-1$, then $t=x\restr m$ gives a cut
containing $x$ but not $y$. Every cut code is monotone, so the resulting
binary code is an order embedding.

Every limit ordinal has the form $\lambda=\omega\cdot\eta$. Apply the
countable-alphabet code separately to its successive $\omega$-blocks. The
concatenated binary length is $\omega\cdot\eta=\lambda$.
\end{proof}

\begin{theorem}[All fixed-type strata]\label{lwo:thm:all-strata}
Let $X\subseteq\R$ be infinite, with the inherited order, and let
$|\alpha|=|X|$. Write
\[
                  \alpha=\lambda+n,
       \qquad \lambda\text{ a nonzero limit ordinal},\quad n<\omega.
\]
Then
\begin{equation}\label{lwo:eq:all-strata}
                   \W_\alpha(X)\equi
                         \B_\lambda\times_{\lex}\{0,\ldots,n!-1\}.
\end{equation}
In particular, this is an exact equimorphism classification of all strata
$\W_\alpha$ of the real well-ordering space.
\end{theorem}
\begin{proof}
The lower embedding has a uniform construction. Fix an $n$-element set
$F\subseteq X$ and partition $X\setminus F$ into pairs indexed by
$\lambda$. Orient these pairs according to a member of $\B_\lambda$ and
append one of the $n!$ permutations of $F$. Since $\lambda$ is a limit
ordinal, $2\cdot\lambda=\lambda$, so the result has exactly the required
length $\lambda+n$.

For the upper embedding, induct on the infinite cardinal $\mu=|X|$,
proving the assertion simultaneously for all inherited real suborders of
that size and all appropriate $\alpha$. If $\mu=\aleph_0$, code the first
$\lambda$ entries by \cref{lwo:lem:countable-alphabet}, and rank the permutation
of the remaining $n$ elements. This gives the required target.

Suppose $\mu$ is uncountable and the assertion is known for smaller infinite
cardinals. Use ordinal division to write
\[
                         \alpha=\mu\cdot\eta+\rho,
                  \qquad 0\leq\rho<\mu,\quad\eta>0,
\]
and put $\delta=\mu\cdot\eta$. Since $\omega\cdot\mu=\mu$,
$\omega\cdot\delta=\delta$. Rational-cut coding therefore embeds the
first $\delta$ real entries into $\B_\delta$ without increasing their
ordinal length.

For a fixed such prefix $p$, the residual set
$Y=X\setminus\ran p$ has cardinality $|\rho|$. If $\rho$ is finite, its
tail is just one of $\rho!$ permutations, which finishes the construction.
If $\rho$ is infinite, write $\rho=\sigma+n$, with $\sigma$ a limit
ordinal. By the induction hypothesis, choose an embedding
\[
             \W_\rho(Y)\emb\B_\sigma\times_{\lex}\{0,\ldots,n!-1\}.
\]
Choose such an embedding for each possible residual set $Y$; this is a
set-sized use of Choice. Append its binary part after the prefix code and
retain its final finite rank. The target is
\[
          \B_{\delta+\sigma}\times_{\lex}\{0,\ldots,n!-1\}
                =\B_\lambda\times_{\lex}\{0,\ldots,n!-1\}.
\]
Two different prefixes are compared inside the first $\delta$ bits. Equal
prefixes have the same residual set and are compared by the chosen
embedding of their tails. Thus the assembled map preserves and reflects
order. The induction is complete.
\end{proof}

\begin{corollary}[Exact length and cube spectrum in every stratum]
\label[corollary]{lwo:cor:all-strata-length}
For $\alpha=\lambda+n$ as above,
\begin{align*}
 \bl(\W_\alpha(X))&=\lambda+\lceil\log_2(n!)\rceil,\\
 \B_\theta\emb\W_\alpha(X)
       &\quad\Longleftrightarrow\quad
          \theta\leq\lambda+\lfloor\log_2(n!)\rfloor.
\end{align*}
For $X=\R$, every stratum still has cardinality $2^\kappa$ and ordinal
spectrum precisely $\{\beta:\beta<\kappa^+\}$.
\end{corollary}
\begin{proof}
The finite-fiber fusion proof of \cref{lwo:thm:finite-capacity} works with any
ordinal $\lambda$ in place of $\kappa$; it only uses full binary cylinders
through $\lambda$ steps. Apply it to \cref{lwo:thm:all-strata}.
For the real case, $|\lambda|=\kappa$ and $\lambda\geq\kappa$, so the
stratum contains $\B_\kappa$ and is a suborder of $\W$. The cardinality and
ordinal conclusions follow from the earlier upper and lower bounds.
\end{proof}

\begin{corollary}[The exact comparison between strata]\label[corollary]{lwo:cor:strata-comparison}
Let $\alpha=\lambda+n$ and $\beta=\mu+m$ have cardinality $\kappa$, with
$\lambda,\mu$ nonzero limit ordinals and $m,n<\omega$. Then
\[
 \W_\alpha\emb\W_\beta
 \quad\Longleftrightarrow\quad
 \lambda<\mu\ \text{ or }\ [\lambda=\mu\text{ and }n!\leq m!].
\]
Thus the fixed strata are linearly ordered by embeddability. Distinct
strata are equimorphic exactly for the pairs $\W_\lambda,\W_{\lambda+1}$
with $\lambda$ a limit ordinal.
\end{corollary}
\begin{proof}
For the same limit length, fusing source cylinders while reading the
$\lambda$ target bits forces an entire source fiber of size $n!$ into a
target fiber of size $m!$. Hence $n!\leq m!$ is necessary, and it is
plainly sufficient. If $\lambda<\mu$, then
$\lambda+\lceil\log_2(n!)\rceil<\mu$, since $\mu$ is a limit ordinal.
Code the source into that shorter cube and then into $\B_\mu$ in the
target. If $\mu<\lambda$, the target codes into a cube of length
$\mu+\lceil\log_2(m!)\rceil<\lambda$, whereas the source contains
$\B_\lambda$. The strict cube hierarchy excludes an embedding.
\end{proof}

\begin{example}
For the following infinite-tail types, the fixed stratum is equimorphic
with the binary cube of exactly the same length:
\[
 \W_{\kappa+\omega}\equi\B_{\kappa+\omega},\qquad
 \W_{\kappa+\omega^2}\equi\B_{\kappa+\omega^2},\qquad
 \W_{\kappa\cdot2}\equi\B_{\kappa\cdot2}.
\]
More generally, $\W_\lambda\equi\W_{\lambda+1}\equi\B_\lambda$
for every limit $\lambda$ of cardinality $\kappa$. The formula is about
embeddability, not an order isomorphism with a cube that has endpoints.
\end{example}

\begin{proposition}[Density and cellularity across the strata]
\label[proposition]{lwo:prop:all-strata-density}
For every $\kappa\leq\alpha<\kappa^+$,
\[
 \dens(\W_\alpha)=\cell(\W_\alpha)=
 \begin{cases}
 2^{<\kappa},&\alpha=\kappa,\ \kappa+1,\ \kappa+2,\\
 2^\kappa,&\alpha\geq\kappa+3.
 \end{cases}
\]
\end{proposition}
\begin{proof}
The first three cases are proved in \cref{lwo:prop:finite-topology}. For the
others, write $\alpha=\kappa+\tau$, where $\tau\geq3$. Choose a set
$Y\subseteq\R$ of cardinality $|\tau|$ whose complement still has
cardinality $\kappa$. There are $2^\kappa$ distinct exhaustive prefixes of
length $\kappa$ listing $\R\setminus Y$. Each corresponding cylinder in
$\W_\alpha$ is isomorphic to $\W_\tau(Y)$ and has at least three points.
Inside each choose a nonempty open interval between three such points.
These intervals are pairwise disjoint. Cardinality gives the upper bound.
\end{proof}

\section{Self-absorption and closure under lexicographic constructions}
\label{lwo:sec:absorption}

The availability of arbitrary exhaustive lengths lets one concatenate entire
well-orderings, rather than squeezing them into a fixed coordinate budget.

\begin{lemma}[Many disjoint real copies]\label[lemma]{lwo:lem:disjoint-reals}
There are $\kappa$ pairwise disjoint subsets of $(0,1)$, each admitting an
order isomorphism from $\R$ onto that subset.
\end{lemma}
\begin{proof}
Use the increasing injection $j:\R\to\two^\omega$ from
\cref{lwo:lem:real-code}. For $t\in\two^\omega$, interleave the fixed bits of
$t$ with the variable bits of $j(r)$:
\[
 z_{t,r}(2n)=t(n),\qquad z_{t,r}(2n+1)=j(r)(n).
\]
For fixed $t$, first disagreements are those of $j(r)$, so
$r\mapsto T(z_{t,r})$ is increasing. Different $t$ give disjoint images
because ternary coding is injective. An affine rescaling puts all images
strictly inside $(0,1)$. There are $2^{\aleph_0}=\kappa$ parameters $t$.
\end{proof}

\begin{theorem}[Sharp power absorption]\label{lwo:thm:absorption}
For every ordinal $\gamma$,
\begin{equation}\label{lwo:eq:absorption}
                       \W^\gamma_{\lex}\emb\W
                       \quad\Longleftrightarrow\quad\gamma<\kappa^+.
\end{equation}
For $0<\gamma<\kappa^+$, $\W^\gamma_{\lex}\equi\W$.
\end{theorem}
\begin{proof}
Assume $0<\gamma<\kappa^+$. Choose disjoint images
$X_\xi=h_\xi[\R]\subseteq(0,1)$ for $\xi<\gamma$, with each $h_\xi$
increasing. Let
\[
                  Z=\R\setminus\bigcup_{\xi<\gamma}X_\xi,
\]
and choose a fixed exhaustive enumeration $z$ of $Z$. The set $Z$ has
cardinality $\kappa$, because it contains $\R\setminus(0,1)$.
For a tuple $(e_\xi)_{\xi<\gamma}\in\W^\gamma$, define
\begin{equation}\label{lwo:eq:concatenation}
 E((e_\xi)_{\xi<\gamma})
       =\mathop{\cat}_{\xi<\gamma}(h_\xi\circ e_\xi)\cat z.
\end{equation}
All images being disjoint, this is an injective enumeration of all reals.
Its ordinal length is the sum of at most $\kappa$ ordinals, each of
cardinality at most $\kappa$, followed by the length of $z$. The total has
cardinality $\kappa$, so it is an allowed member of $\W$.

For two distinct tuples, let $\xi$ be their first different component.
Earlier components are equal as exhaustive enumerations, including their
lengths; hence the concatenated prefixes are identical. Within component
$\xi$, prefix-freeness supplies a disagreement before either component
ends. Since $h_\xi$ is increasing, that disagreement has exactly the
required sign. Thus $E$ is an order embedding.

The empty product is a singleton and embeds trivially. For
$\gamma\geq\kappa^+$, choose two points of $\W$ and use them as binary
coordinate values. The product then contains $\B_{\kappa^+}$ and hence
$\kappa^+$, which $\W$ cannot contain. Finally, for $\gamma>0$, varying one
coordinate and fixing all other coordinates embeds $\W$ into the product.
\end{proof}
\merge\ Also source~07, Lemma~6.1 and Theorem~6.2 (07:548, 07:563), by
the same interleaved copies of $\R$ and the same concatenation. Source~07
stresses that the copies must be disjoint \emph{ordered copies} of $\R$, not
disjoint intervals, of which only countably many exist (07:543--546), and that
the ranges need not be convex. Absorption is mutual embeddability, not
isomorphism (\cref{lwo:rk:rem:omega-power}).

\begin{corollary}[Closure of the embedding class]\label[corollary]{lwo:cor:closure}
Let $\mathscr E(\W)=\{L:L\emb\W\}$. This class is closed under reversal,
under ordinal-indexed lexicographic products of at most $\kappa$ factors,
and under lexicographic sums whose index order and every summand belong to
$\mathscr E(\W)$.
\end{corollary}
\begin{proof}
Negation of all real entries reverses $\W$, so the class is closed under
reversal. Embed each factor into $\W$ and apply \cref{lwo:thm:absorption} for
products. For a sum $\sum_{i\in I}L_i$, choose embeddings $f:I\to\W$ and
$g_i:L_i\to\W$. Then
\[
                      (i,x)\longmapsto(f(i),g_i(x))
\]
embeds the sum into $\W\times_{\lex}\W$, and the latter embeds into $\W$.
The family of chosen embeddings is a set-sized choice in ZFC.
\end{proof}
\merge\ Also source~07, Corollary~6.3 (07:599).

\begin{corollary}[Failure of absorption at fixed length]\label[corollary]{lwo:cor:no-fixed-square}
For every ordinal $\gamma$,
\[
                 (\W_\kappa)^\gamma_{\lex}\equi\B_{\kappa\cdot\gamma},
                 \qquad
                 \A^\gamma_{\lex}\equi\B_{\kappa\cdot\gamma}.
\]
Either product embeds back into its original factor exactly when
$\gamma\leq1$.
\end{corollary}
\begin{proof}
Equimorphisms can be applied coordinatewise, and consecutive blocks of
$\kappa$ binary coordinates have ordinal length $\kappa\cdot\gamma$.
For $\gamma\geq2$, this length exceeds $\kappa$; use
\cref{lwo:thm:cube-strict}. The cases $0$ and $1$ are immediate.
\end{proof}
\merge\ Source~07 states the same equimorphism with real coordinates,
$(\W_\kappa)^\beta_{\lex}\equi\R^{\kappa\cdot\beta}_{\lex}$ and
$\A^\beta_{\lex}\equi\R^{\kappa\cdot\beta}_{\lex}$ for $\beta>0$ (its
Proposition~6.4 and Theorem~11.2, 07:613, 07:1129). This is the same
statement: $\R^{\kappa\cdot\beta}_{\lex}\equi\B_{\omega\cdot\kappa\cdot\beta}$
and $\omega\cdot\kappa\cdot\beta=\kappa\cdot\beta$ by
\eqref{lwo:eq:omega-kappa}. Source~07 adds the rank
$\repr((\W_\kappa)^\beta_{\lex})=\repr(\A^\beta_{\lex})=\kappa\cdot\beta$.

\subsection{Many convex copies inside the full order}
The preceding embeddings need not have convex images. There is also a
particularly strong convex version for a large family of copies.

\begin{proposition}\label[proposition]{lwo:prop:convex-copies}
$\W$ contains $2^\kappa$ pairwise disjoint open convex suborders, each
isomorphic to $\W$. They can be chosen so that their index order is
$\B_\kappa$.
\end{proposition}
\begin{proof}
Let $I=(0,1)$. Partition $\R\setminus I$ into $\kappa$ ordered pairs.
For $x\in\two^\kappa$, let $s_x$ be the pair-swap enumeration of
$\R\setminus I$. The cylinders $C_{s_x}$ are pairwise disjoint and ordered
according to $x$. By \cref{lwo:lem:cylinder}, each is isomorphic to $\W(I)$,
and an increasing bijection $\R\to I$ identifies $\W(I)$ with $\W$.
The order $\W(I)$ has no endpoints, since $I$ has none. Thus every point of
a cylinder is bracketed by two points in that cylinder, making the convex
cylinder open in $\W$.
\end{proof}

\section{Adjacent pairs and finite condensation}\label{lwo:sec:adjacency}

\begin{lemma}[Extrema of exhaustive enumeration orders]\label[lemma]{lwo:lem:extrema}
For any nonempty linearly ordered set $Y$, $\W(Y)$ has a least element
exactly when the given order on $Y$ is a well-order. That element is the
increasing enumeration of $Y$. It has a greatest element exactly when the
given order on $Y$ is a reverse well-order, in which case it is the
decreasing enumeration.
\end{lemma}
\begin{proof}
In a lexicographically least exhaustive enumeration, the value at each
coordinate must be the least as-yet-unused element. Otherwise a smaller
unused choice, followed by any exhaustive enumeration of the residual set,
would give a smaller member of $\W(Y)$. Thus a least enumeration must be
increasing and exhibit $Y$ as a well-order. Conversely, at the first
changed coordinate of any other enumeration, the increasing enumeration
uses the least available value, so it is lexicographically smaller.
Reversing the ground order proves the maximum assertion.
\end{proof}
\merge\ Also source~06, Lemma~3.3 (06:263), and the first half of
source~07, Lemma~8.1 (07:714).

\begin{theorem}[Complete adjacency criterion]\label{lwo:thm:adjacency}
Let $e<f$ in $\W(X)$, where $X$ is any infinite linearly ordered set.
Let $s$ be their maximal common prefix, of length $\delta$, and write
\[
                Y=X\setminus\ran s,\qquad a=e(\delta),\quad b=f(\delta).
\]
Then $e$ and $f$ are adjacent if and only if all of the following hold:
\begin{enumerate}[label=(\alph*)]
\item $Y$ is finite;
\item $a<b$ are consecutive elements of $Y$ in its given order;
\item after $a$, the enumeration $e$ lists $Y\setminus\{a\}$ in decreasing
order;
\item after $b$, the enumeration $f$ lists $Y\setminus\{b\}$ in increasing
order.
\end{enumerate}
In particular, an adjacent pair has equal enumeration lengths and differs
only in a finite final segment.
\end{theorem}
\begin{proof}
If $a<c<b$ for some $c\in Y$, then an exhaustive extension of $s\cat c$
lies strictly between $e$ and $f$. If $e$ is not the greatest member of its
$a$-cylinder, a larger member of that cylinder also lies between them.
Likewise, $f$ must be the least member of its $b$-cylinder. By
\cref{lwo:lem:extrema}, $Y\setminus\{a\}$ must be reverse well-ordered and
$Y\setminus\{b\}$ must be well-ordered.

The intersection $Y\setminus\{a,b\}$ is consequently both well-ordered and
reverse well-ordered. It is finite: an infinite well-order contains an
increasing $\omega$-sequence, whose range has no greatest member and cannot
be reverse well-ordered. Thus $Y$ is finite. The two extremal tails are then
exactly the decreasing and increasing finite lists in the statement.

Conversely, any point between $e$ and $f$ would have to extend $s$, by
convexity. Its next entry could not lie strictly between $a$ and $b$.
Choosing $a$ cannot exceed the greatest tail $e$, and choosing $b$ cannot
precede the least tail $f$. No intermediate point exists.
\end{proof}
\merge\ Also source~06, Theorem~6.2 (06:548), and source~07,
Theorem~8.2 (07:747), by the same argument. Source~07 states the criterion
without condition~(a) and derives finiteness of $Y$; its conditions (ii) and
(iii) are (c) and (d) here, through \cref{lwo:lem:extrema}.

\begin{example}[A jump and an isolated point]
Choose any exhaustive prefix $s$ enumerating
$\R\setminus\{a,b\}$, with $a<b$. Then
\[
                         s\cat(a,b)<s\cat(b,a)
\]
is an adjacent pair in the full order $\W$, not just in a fixed-length
suborder.

If $a<b<c$ and $s$ enumerates $\R\setminus\{a,b,c\}$, its cylinder is the
six-point chain
\[
\begin{split}
 s\cat(a,b,c)&<s\cat(a,c,b)<s\cat(b,a,c)\\
             &<s\cat(b,c,a)<s\cat(c,a,b)<s\cat(c,b,a).
\end{split}
\]
The four interior points are isolated in the order topology of $\W$.
\end{example}
\merge\ This is also source~06's Example~6.1 and its Appendix~B. Source~06
reads the middle jump $s\cat(b,c,a)<s\cat(c,a,b)$ through
\cref{lwo:thm:adjacency}: the old suffix after $b$ is decreasing ($c,a$), the
new suffix after $c$ is increasing ($a,b$), and $b,c$ are adjacent in the
unused alphabet. It warns that the example must not be extrapolated to all
points: an enumeration of length $\kappa$ has no finite final tail, lies in a
singleton finite-condensation class, and has uncountable character on both
sides; nor should a permutation generator be read as an algorithm that
enumerates all well-orderings of $\R$.

\begin{definition}
The finite-condensation relation on a linear order is
\[
        x\sim_{\mathrm{fin}}y\quad\Longleftrightarrow\quad
        [\min(x,y),\max(x,y)]\text{ is finite}.
\]
It is a convex equivalence relation.
\end{definition}

\begin{theorem}[Exact finite blocks]\label{lwo:thm:blocks}
Let $e\in\W$, and write its length uniquely as
$\alpha=\lambda+n$, where $\lambda$ is a nonzero limit ordinal and
$n<\omega$. If $n=0$, its finite-condensation class is a singleton. If
$n>0$, its class consists exactly of all $n!$ permutations of the final
$n$ reals after the fixed prefix $e\restr\lambda$.

These classes are the maximal finite convex blocks. A point is isolated
exactly when it is an interior point of one of the blocks with $n\geq3$.
\end{theorem}
\begin{proof}
The displayed permutations form a finite convex cylinder. If a point is
adjacent to $e$, \cref{lwo:thm:adjacency} says that the first disagreement
leaves only finitely many entries. Hence the two enumerations have the same
length and agree before its final finite remainder $n$. No adjacency can
cross to a different displayed block. Every finite interval is connected
by finitely many adjacency steps, so the displayed block is the full
finite-condensation class.

In a linear order with no endpoints, a point is isolated exactly when it
has both an immediate predecessor and an immediate successor. The first
and last points of these maximal blocks have no adjacent point beyond the
block; an interior point has both required neighbors. A factorial block has
an interior point precisely when $n!\geq3$, equivalently $n\geq3$.
\end{proof}
\merge\ Also source~06, Proposition~6.3 and Corollary~6.4 (06:590,
06:607), and source~07, Theorem~8.3 (07:770). Source~06 counts the isolated
points of a block: exactly its $n!-2$ nonextremal permutations. Source~07 makes
explicit why an adjacency chain cannot leave the block: for the length
$\lambda+n$, a common prefix with finitely many unused values has length at
least $\lambda$, since adding finitely many positions to an ordinal below the
limit $\lambda$ cannot reach $\lambda$. It observes that the finite blocks
explain why extrapolating density from $\W_\kappa$ to $\W$ is incorrect: once a
prefix has used all but finitely many reals, there is no reservoir of
intermediate values left (07:813--815).

The quotient by finite condensation is a dense linear order: if two
different classes were adjacent in the quotient, their union would be a
finite interval and they would be the same class. We do not need to assume
that this quotient is isomorphic to any familiar generalized real line.
Its precise isomorphism type remains a separate question.

\merge\ Sources~06 and~07 prove that $\Wq$ is equimorphic with $\W$
(\cref{lwo:sp:thm:quotient}; source~06, Theorem~11.1, 06:1133; source~07,
Proposition~10.5, 07:1072),
answering the question asked in source~02's \cref{lwo:sec:research}. Its
isomorphism type is still open.

\section{Order topology, size of intervals, and cofinalities}
\label{lwo:sec:topology}

\subsection{Global invariants of the full order}

\begin{proposition}\label[proposition]{lwo:prop:full-topology}
The order $\W$ has no endpoints, is isomorphic to its reverse, and satisfies
\[
 \cf(\W)=\ci(\W)=\omega,
 \qquad
 \dens(\W)=\cell(\W)=2^\kappa.
\]
It has exactly $2^\kappa$ isolated points.
\end{proposition}
\begin{proof}
Every real can be the first entry of an exhaustive enumeration. Choosing
first entries tending to $+\infty$ gives a countable cofinal sequence;
choosing first entries tending to $-\infty$ gives a countable coinitial
sequence. Neither endpoint exists. The map $e\mapsto -e$, with the same
ordinal domain, is an order-reversing bijection.

Fix three reals and use all pair-swap enumerations of their complement as
prefixes of length $\kappa$. For each prefix choose one of the four
interior permutations in its six-point block. This gives $2^\kappa$
isolated points by \cref{lwo:thm:blocks}. Each isolated singleton must be met
by a dense set, and the singletons are pairwise disjoint nonempty open
sets. The upper bounds follow from $|\W|=2^\kappa$.
\end{proof}
\merge\ Also source~06, Corollary~3.6 and Theorem~6.6 (06:320, 06:650),
and source~07, Lemma~8.1 and Corollary~8.4 (07:714, 07:798). Both add the
weight: $\weight(\W)=2^\kappa$, since the interval base has at most
$|\W|^2=2^\kappa$ members and $\cell\le\dens\le\weight$. Source~06 also
shows that the isolated points are dense (\cref{lwo:sp:thm:denseisolated});
sources~06 and~07 obtain $\cf\W=\ci\W=\omega$ from the stronger fact that every
residual order $\W(Y)$, $Y\subseteq\R$, has cofinality and coinitiality in
$\{1,\omega\}$ (\cref{lwo:sp:lem:countends}).

There is a second, independent reason for the large cellularity:
\cref{lwo:prop:convex-copies} gives $2^\kappa$ disjoint open copies of $\W$.
Thus the phenomenon is not solely a by-product of isolated finite tails.
For example, selecting one point from each finite-condensation class still
leaves $2^\kappa$ disjoint nontrivial interval regions coming from those
large cylinders.

It would nevertheless be false to say that \emph{every} interval in $\W$
is large. Some intervals are empty, some are finite, and a cylinder leaving
countably infinitely many unused reals has cardinality only $2^{\aleph_0}$.
The residual-set lemma, not the ambient cardinality, controls such examples.

\subsection{The dense fixed-length order}

\begin{proposition}\label[proposition]{lwo:prop:fixed-topology}
The order $\W_\kappa$ is dense, has no endpoints, and is self-dual. Its
global cofinality and coinitiality are $\omega$, and
\begin{equation}\label{lwo:eq:fixed-density}
                \dens(\W_\kappa)=\cell(\W_\kappa)=2^{<\kappa}.
\end{equation}
At every point $e$, the cofinality of the left side and the coinitiality of
the right side are both $\cf(\kappa)$.
\end{proposition}
\begin{proof}
Let $e<f$ first disagree at $\delta<\kappa$. The common prefix has used
fewer than $\kappa$ reals, whereas the interval
$(e(\delta),f(\delta))$ has cardinality $\kappa$. Choose an unused real
strictly between the two entries. Extending the resulting prefix to a
bijection of length $\kappa$ gives an intermediate point. This proves
density. Endpoints, global cofinalities, and self-duality are treated as for
$\W$.

There are at most
\[
       \sum_{\xi<\kappa}\kappa^{|\xi|}=2^{<\kappa}
\]
prefixes of length below $\kappa$. Here finite nonzero exponents contribute
$\kappa=2^{\aleph_0}$, and for infinite $\mu<\kappa$,
$\kappa^\mu=(2^{\aleph_0})^\mu=2^\mu$. Choose one exhaustive extension
for each such prefix. The resulting set is dense, since the intermediate
prefix in the previous paragraph determines a whole cylinder inside the
given interval.

For every infinite cardinal $\mu<\kappa$, orient $\mu$ disjoint pairs to
obtain $2^\mu$ incompatible prefixes of length $2\cdot\mu<\kappa$.
Their cylinders are nonempty, convex, and have no endpoints: one can change
a later coordinate upward or downward while fewer than $\kappa$ values have
been used. The cylinders are therefore open. This proves the lower
cellularity bound after taking the supremum over $\mu<\kappa$.
Cellularity never exceeds density, proving \eqref{lwo:eq:fixed-density}.

Finally, for every prefix length $\xi<\kappa$, there are points on both
sides of $e$ agreeing with it below $\xi$: change a still later coordinate
to an unused smaller or larger real and complete the enumeration. Choose
such points at a cofinal set of $\cf(\kappa)$ prefix lengths. They are
cofinal on the left and coinitial on the right. Conversely, fewer than
$\cf(\kappa)$ points on one side have a bounded set of first-disagreement
coordinates with $e$. A perturbation agreeing longer than that bound lies
closer to $e$ than all of them. This proves both local lower bounds.
\end{proof}
\merge\ Also source~06, Theorems~9.2 and~9.3 (06:973, 06:1001), and
source~07, Theorem~7.1 (07:637), by the same short-prefix arguments. Both add
$\weight(\W_\kappa)=2^{<\kappa}$, since the short-prefix cylinders, at most
$2^{<\kappa}$ of them, are clopen and form a base
(\cref{lwo:sp:thm:minimaldense,lwo:sp:prop:minimalweight}).

\begin{remark}[Regularity is not assumed]
The formula uses $\cf(\kappa)$, not automatically $\kappa$. However,
$\cf(2^{\aleph_0})>\omega$. Here is a direct diagonal proof. If
$\two^\omega=\bigcup_{n<\omega}X_n$ with each $|X_n|<2^{\aleph_0}$,
partition $\omega$ into infinite sets $I_n$. Choose a pattern on $I_n$ that
is not the restriction of any member of $X_n$. The union of these patterns
belongs to none of the $X_n$, a contradiction. A countable cofinal sequence
of smaller cardinals would produce just such a cover.
Consequently $\W_\kappa$ is not first countable at any point.
\end{remark}
\merge\ Sources~06 (06:1029--1038) and~07 (Lemma~3.2, 07:305--310) prove
$\cf\kappa>\omega$ through K\"onig's inequality: if cardinals $\kappa_n<\kappa$
were cofinal in $\kappa$, then
$\kappa=\sum_n\kappa_n<\prod_n\kappa_n^+\le\kappa^{\aleph_0}=\kappa$, a
contradiction; the strict inequality is the same diagonal argument.

\subsection{A topology refinement for finite tails}

\begin{proposition}\label[proposition]{lwo:prop:finite-topology}
For finite $n\geq0$, the following hold.
\begin{enumerate}[label=(\alph*)]
\item For $n=0,1$, $\W_{\kappa+n}$ is dense and has no isolated points.
\item For $n=2$, it has $2^\kappa$ adjacent pairs and no isolated points.
Its density and cellularity are still $2^{<\kappa}$.
\item For $n\geq3$, it has $2^\kappa$ isolated points and hence density and
cellularity $2^\kappa$.
\end{enumerate}
For $n=0,1$, density and cellularity are $2^{<\kappa}$ as well. The weight
of $\W_{\kappa+2}$ is $2^\kappa$, even when its density is smaller.
\end{proposition}
\begin{proof}
If two points first disagree below $\kappa$, the unused-real argument in
\cref{lwo:prop:fixed-topology} gives an intermediate point of the same length.
If their first $\kappa$ entries agree, they lie in one finite permutation
block. The block sizes are $1,1,2,6,\ldots$, proving the assertions about
adjacency and isolation. Varying a pair-swap prefix while fixing the final
finite set gives $2^\kappa$ blocks.

For $n\leq2$, a nonempty open interval cannot be finite: a finite open set
in a $T_1$ space would contain an isolated point. Such an interval therefore
contains two points with different prefixes below $\kappa$, since each
full $\kappa$-prefix has at most two extensions. As before, a shorter
cylinder fits between those points. Choosing one extension of every prefix
of length below $\kappa$ gives the upper density bound, and incompatible
short prefixes give the lower cellularity bound.

For the weight assertion, every adjacent pair $a<b$ requires a basic open
set containing $b$ and contained in $(a,\infty)$. Different adjacent pairs
require different such basic sets: if $a<b\leq c<d$, one set cannot both
contain $b$ and be contained in $(c,\infty)$. Thus any base has size at
least $2^\kappa$. The set of all endpoint intervals supplies the reverse
bound.
\end{proof}

\section{Finer properties of the Kanovei--Shelah index}\label{lwo:sec:KS-topology}

Equimorphism with $\B_\kappa$ gives the exact embeddable orders, but does not
determine global cofinality or topology. Those require separate arguments.

\begin{lemma}\label[lemma]{lwo:lem:A-cylinders}
Every nonempty cylinder in $\A$ determined by fewer than $\kappa$
coordinates has no endpoints. In particular, $\A$ is a dense linear order
without endpoints.
\end{lemma}
\begin{proof}
Fix $a\in\A_U$ and a prefix of length $\xi<\kappa$. Every ultrafilter
contains $\N$ and has $\kappa$ members. Since the prefix lists fewer than
$\kappa$ values, arbitrarily far out in the tail there is a value other
than $\N$. Replacing such a value by $\N$, and then enumerating all of
$U$ on the remaining coordinates, gives a larger point with the prescribed
prefix.

For a smaller point, if $U$ has a least member, the tail cannot consist
only of this member; choose a later larger value and replace it by the
least one. If $U$ has no least member, replace any chosen later value by
a smaller member of $U$. In either case complete the remaining coordinates
by an enumeration of $U$. Thus there are both larger and smaller points in
the cylinder.

If $a<b$ first differ at $\delta$, apply the preceding construction to
$a$ with its prefix through $\delta$ fixed. A larger extension in that
cylinder still lies below $b$.
\end{proof}

\begin{proposition}\label[proposition]{lwo:prop:A-topology}
The full index $\A$ satisfies
\[
 |\A|=2^\kappa,\qquad
 \dens(\A)=\cell(\A)=2^{<\kappa},\qquad
 \cf(\A)=\cf(\kappa),\qquad\ci(\A)=\omega.
\]
At every point, the left and right local cofinalities are both
$\cf(\kappa)$. In particular, $\A$ is not isomorphic to its reverse and
is not isomorphic to $\W_\kappa$.
\end{proposition}
\begin{proof}
Cardinality follows from the embeddings in \cref{lwo:thm:fixed-equivalence}.
For density, there are at most $2^{<\kappa}$ prefixes of length below
$\kappa$. By \cref{lwo:lem:A-cylinders}, every nonempty interval contains a
cylinder with such a prefix; choose one extension for each extendible
prefix. Conversely, fix $U$ and two of its members. All binary choices for
a prefix of an infinite length $\mu<\kappa$ extend to members of $\A_U$.
Their full $\A$-cylinders are disjoint nonempty open sets. This proves the
cellularity and density formulas.

Let $\nu(a)$ be the first coordinate at which $a$ is not $\N$. It exists
because $\ran a$ is an ultrafilter with $\kappa$ members. For a cofinal set
of ordinals $\delta<\kappa$, construct a map beginning with $\delta$ copies
of $\N$, then a fixed non-$\N$ member of a fixed ultrafilter, then an
exhaustive tail. These maps are cofinal in $\A$. A family of fewer than
$\cf(\kappa)$ maps has its values of $\nu(a)$ bounded below $\kappa$;
a sufficiently long initial string of $\N$ lies above the whole family.
Hence $\cf(\A)=\cf(\kappa)$.

For each natural number $n$, take a member of $\A$ whose first value is
$\{n\}$ and whose range is the principal ultrafilter at $n$. The singletons
$\{n\}$ descend coinitially among the nonempty subsets of $\N$ in binary
lexicographic order. Therefore these maps are coinitial in $\A$, which has
no least element. This proves $\ci(\A)=\omega$.

Local cofinalities follow from the same bounded-disagreement argument used
for $\W_\kappa$, now using \cref{lwo:lem:A-cylinders} for perturbations after
arbitrary prescribed prefixes. Finally, $\cf(\kappa)>\omega$, so the two
global cofinalities are unequal. They rule out self-duality and also rule
out an isomorphism with $\W_\kappa$.
\end{proof}
\merge\ Also source~07, Theorem~11.3 (07:1163), with the same
cofinality and coinitiality witnesses; source~07 adds the weight, a clopen base
of short prefixes, and coinitiality witnesses for the free-ultrafilter
variant of $\A$ (\cref{lwo:rk:prop:Atopology}).

For comparison, a fixed principal-ultrafilter fiber $\A_U$ has coinitiality
$\cf(\kappa)$: it has a least available coordinate value, which can be
repeated for arbitrarily long initial segments. A nonprincipal fiber has
coinitiality $\omega$, because the cofinite sets form a decreasing
coinitial sequence in its coordinate order. All fibers nevertheless have
the same embeddability class by \cref{lwo:thm:fixed-equivalence}.

\section{Explicit countable gaps and Dedekind completion}\label{lwo:sec:completion}

\subsection{A gap caused by a repeated real}

\begin{proposition}\label[proposition]{lwo:prop:real-gap}
Both $\W$ and $\W_\kappa$ have an unfilled cut of cofinality type
$(\omega,\omega)$.
\end{proposition}
\begin{proof}
Divide either order into the following initial and final segments:
\begin{align*}
 L&=\{e:e(0)<0\text{ or }[e(0)=0\text{ and }e(1)<0]\},\\
 R&=\{e:e(0)>0\text{ or }[e(0)=0\text{ and }e(1)>0]\}.
\end{align*}
Every exhaustive enumeration has at least two entries, and injectivity
excludes $e(0)=e(1)=0$. Thus this is a partition and $L<R$.
For $n\geq1$, complete $(0,-1/n)$ and $(0,1/n)$ to exhaustive
enumerations of length $\kappa$. The first family is cofinal in $L$;
the second is coinitial in $R$. Neither side has an endpoint, and a point
between all the bounds would have to start $(0,0)$, which is forbidden.
\end{proof}
\merge\ Also source~07, Example~7.3 (07:703), with the same bounds;
source~07 concludes that neither order has the countable interpolation
property of an $\eta_1$-order.

This also illustrates why large local character at every point does not
imply countable saturation as a dense order. The obstruction is an omitted
cut, rather than a point with small local character.

\subsection{A gap caused by incompatible ultrafilter values}

\begin{proposition}\label[proposition]{lwo:prop:A-gap}
The order $\A$ also has an unfilled $(\omega,\omega)$ cut.
\end{proposition}
\begin{proof}
Let $u$ be the even natural numbers and $v$ the odd natural numbers. Put
\[
 v^-_n=(v\setminus\{2n+1\})\cup\{2n+2\},
 \qquad
 v^+_n=v\cup\{2n\}.
\]
In binary lexicographic order, $v^-_n$ increases toward $v$ from below,
and $v^+_n$ decreases toward $v$ from above. Each pair $u,v^-_n$ belongs to
the principal ultrafilter at $2n+2$; each pair $u,v^+_n$ belongs to the
principal ultrafilter at $2n$. Therefore both two-term prefixes
$(u,v^-_n)$ and $(u,v^+_n)$ extend to members of $\A$.

A map between all these bounds must have first value $u$ and second value
$v$. Indeed, if a binary sequence differs from $v$, its first differing bit
is eventually shared by neither bounding side in the required direction.
But no ultrafilter contains the disjoint sets $u$ and $v$. Thus the cut is
not realized. The initial segment of maps with first value below $u$, or
first value $u$ and second value below $v$, has the lower sequence cofinal;
the complementary final segment has the upper sequence coinitial.
\end{proof}
\merge\ Source~07 proves the same statement with different bounds, and
adds a gap of character $(\cf\kappa,\cf\kappa)$ (\cref{lwo:rk:prop:Agaps};
its Proposition~11.4, 07:1218).

These statements concern the index orders. They do not contradict the
countable saturation of the elementary extension constructed from the index
in the Kanovei--Shelah theorem.

\subsection{Completion does not create longer ordinals}

We use the Dedekind completion that fills unrealized two-sided cuts and
keeps existing points and jumps. Adding endpoints is optional and does not
change the following cardinal or ordinal conclusions.

\begin{lemma}[Cuts have small cofinal witnesses]\label[lemma]{lwo:lem:small-cut}
Every nonempty initial segment of $\W$ has a cofinal subset of cardinality
at most $\kappa$. Every nonempty final segment has a coinitial subset of
cardinality at most $\kappa$.
\end{lemma}
\begin{proof}
If an initial segment had no such cofinal subset, at every stage
$\xi<\kappa^+$ one could choose a point larger than all previously chosen
points: the preceding set has cardinality at most $\kappa$ and is not
cofinal. This constructs the forbidden increasing sequence of length
$\kappa^+$. Reverse the order for the other assertion.
\end{proof}
\merge\ Also source~06, Corollary~5.5 (06:520), for every nonempty
suborder of $\W(X)$, and source~07, Corollary~4.6 (07:417).

\begin{theorem}[Completion invariants]\label{lwo:thm:completion}
Let $\widehat\W$ be a Dedekind completion of $\W$. Then
\[
 |\widehat\W|=2^\kappa,
 \qquad
 \beta\emb\widehat\W\ \Longleftrightarrow\ \beta<\kappa^+,
 \qquad
 \bl(\widehat\W)=\kappa^+.
\]
The ordinal equivalence also holds for reverse ordinals. Every two-sided
cut of $\W$ has each of its two cofinalities at most $\kappa$.
\end{theorem}
\begin{proof}
By \cref{lwo:lem:small-cut}, every lower half of a cut is determined by a
cofinal subset of size at most $\kappa$. Thus the number of cuts is at most
\[
                 |\W|^\kappa=(2^\kappa)^\kappa=2^\kappa.
\]
The lower cardinal bound follows from the copy of $\W$ already present.

Suppose the completion contained an increasing sequence of length
$\kappa^+$. Between the first and third members of each successive
three-point block, choose a point of $\W$. Such a point always exists:
if the middle point is old, use it; if it is an added cut, the cut property
supplies an old point between the outer endpoints. The resulting points
form a strictly increasing sequence of length $\kappa^+$ in $\W$, a
contradiction. The same argument applies in reverse. The positive ordinal
embeddings were already present in $\W$.

For binary length, $\W\emb\widehat\W$ gives the lower bound. The cube
$\B_{\kappa^+}$ is Dedekind complete, as verified below. An embedding
$f:\W\to\B_{\kappa^+}$ extends over each unrealized cut $(L,R)$ by
assigning it $\sup f[L]$. This value is strictly above every $f(l)$ because
$L$ has no greatest element, and strictly below every $f(r)$ because $R$
has no least element. Distinct cuts remain distinct, by an old point
between them. Thus the extension is an order embedding. If endpoints are
to be added, first put $f[\W]$ inside a proper binary cylinder by a fixed
finite prefix and then use the ambient endpoints. Its length remains
$\kappa^+$ because finite $+\,\kappa^+=\kappa^+$.
\end{proof}
\merge\ The cardinality also follows from source~07's count of initial
segments (Corollary~4.6, 07:417) and from source~06's count of proper gaps
(\cref{lwo:sp:thm:gapcount}; Theorem~7.6, 06:824), which shows that there
are exactly $2^\kappa$ of them.
Source~06 also proves the ordinal assertion by the same interval argument
(06:851--856).

\begin{lemma}[Completeness of binary lexicographic cubes]\label[lemma]{lwo:lem:cube-complete}
Every $\B_\theta$ is Dedekind complete and has both endpoints.
\end{lemma}
\begin{proof}
The constant sequences give endpoints. For a nonempty subset $S$, construct
its supremum by retaining a nonempty set of candidates with a common
prefix. At a successor step choose bit $1$ if some current candidate has
bit $1$, and retain those candidates; otherwise choose bit $0$ and retain
all current candidates. At a limit, intersect the candidate sets.

If the first empty intersection occurs at a limit $\lambda\leq\theta$,
complete the accumulated prefix by zeros. Every member of $S$ was discarded
at an earlier first disagreement where it had $0$ and the chosen prefix had
$1$, so the completed sequence is an upper bound. Any strictly smaller
sequence first disagrees with this upper bound before $\lambda$; at that
stage there was a candidate with the chosen $1$, which is larger than the
proposed smaller bound. Thus the completed sequence is the least upper
bound. If no empty intersection occurs through $\theta$, the final sequence
belongs to $S$ and is its maximum. This covers all cases.
\end{proof}

\begin{remark}
The completion theorem does not prove $\widehat\W\emb\W$. Both orders have
the same cardinality, ordinal spectrum, and binary length, but these
invariants do not in general decide embeddability. Whether the completion
can be absorbed is one of the natural next questions.
\end{remark}

\section{Universality in every slice}\label{lwo:sp:sec:slices}

\merge\ This section and the next five are source~06's, with source~07's
proofs of the same statements recorded as routes. The cardinality,
universality and ordinal spectrum of \cref{lwo:sec:enumerations,lwo:sec:ordinal}
hold for every infinite linear alphabet $X$; source~06 sharpens universality
to every single slice $\W_\beta(X)$.

\begin{theorem}[Universality in every fixed-type slice; source~06, Theorem~4.4]
\label{lwo:sp:thm:universality}
Let $X$ have infinite cardinality $\nu$, and let $\beta$ be any ordinal of
cardinality $\nu$. Every linear order $L$ of cardinality at most $\nu$ embeds
in $\W_\beta(X)$. Consequently $|\W_\beta(X)|=2^\nu$.
\end{theorem}
\begin{proof}
Write $\beta=\nu+\tau$, where $\tau$ is the order type of the interval of
positions after the initial $\nu$ positions. Partition $X=S\sqcup T$ so
that $|S|=\nu$ and $|T|=|\tau|$, taking $T=\varnothing$ if $\tau=0$.
Pair-code all bits of $\B_\nu$ in $S$ over the first $\nu$ positions
(\cref{lwo:prop:cardinality}), then append a fixed bijection $\tau\to T$.
The resulting embeddings all have domain $\nu+\tau=\beta$. Compose with
\cref{lwo:lem:cut-code}. The same construction embeds all $2^\nu$ bit
strings, proving the cardinality claim together with the relation-code upper
bound.
\end{proof}

For $X=\R$ this contains \cref{lwo:cor:all-strata-length}'s statement that
every real stratum has cardinality $2^\kappa$ and ordinal spectrum
$\{\beta:\beta<\kappa^+\}$, by a direct construction instead of the
all-strata classification.

\begin{remark}[Larger embedded orders; source~06, 06:415--423]
\label[remark]{lwo:sp:rem:separating}
Cardinality at most $\nu$ is a sufficient condition, not a necessary one:
$\B_\nu$ itself has size $2^\nu$ and embeds. Another useful sufficient
condition is a separating family $(I_\xi)_{\xi<\nu}$ of initial segments:
assume that whenever $x<y$, some $I_\xi$ contains $x$ but not $y$. Then the
nonmembership vectors $(\mathbf 1_{x\notin I_\xi})_{\xi<\nu}$ are
coordinatewise increasing and strictly distinguish points, giving an embedding
into $\B_\nu$. Separately, the collection of all initial segments of a linear
order of size $\nu$, ordered by inclusion, embeds by the characteristic
functions of those initial segments.
\end{remark}

\section{Countable ends and dense isolation}\label{lwo:sp:sec:isolation}

\begin{lemma}[Real subsets have countable monotone length; source~06, Lemma~3.4]
\label[lemma]{lwo:sp:lem:realmono}
Every strictly increasing or strictly decreasing ordinal-indexed sequence of
real numbers has countable domain. A subset of $\R$ without a greatest element
has countable cofinality; the dual statement holds for coinitiality.
\end{lemma}
\begin{proof}
For a strictly increasing sequence choose a rational strictly between each
successive pair of terms. These rational intervals are disjoint, so there are
only countably many successor pairs. An uncountable ordinal has uncountably
many successor positions, a contradiction. The decreasing case is identical.
For the last assertion, a subset unbounded above has a countable cofinal
sequence chosen past the integers. A nonempty bounded subset with no maximum
has supremum $r$, and points chosen above $r-1/n$ form a countable cofinal
set. A finite cofinal set would imply a maximum.
\end{proof}

\begin{lemma}[Countable ends of every residual permutation order; source~06, Lemma~3.5]
\label[lemma]{lwo:sp:lem:countends}
For every nonempty $Y\subseteq\R$, each of $\cf(\W(Y))$ and $\ci(\W(Y))$ is
either $1$ or $\omega$. More precisely, the former is $1$ exactly when $Y$ is
reverse well-ordered, and the latter is $1$ exactly when $Y$ is well-ordered.

If $\W(Y)$ has no greatest element, it has a cofinal family of cylinders with
infinite residual alphabet. The dual assertion holds at the lower end.
\end{lemma}
\begin{proof}
Repeatedly remove the greatest still remaining element of $Y$, taking unions
at limits. This greedy decreasing sequence cannot continue for $\omega_1$
steps, by \cref{lwo:sp:lem:realmono}. It therefore stops at a countable
ordinal $\tau$, either because $Y$ has been exhausted or because the
remaining nonempty set $B$ has no maximum. In the first case the decreasing
enumeration is greatest. In the second, let $s$ be the removed prefix. Every
enumeration not extending $s$ is smaller than every member of $C_s$: at its
first disagreement it failed to choose the current maximum. Thus $C_s$ is a
final segment. Choose a strictly increasing cofinal sequence $b_n$ in $B$. The
cylinders $C_{s\cat\langle b_n\rangle}$ are cofinal in $\W(Y)$, and each has
infinite residual alphabet $B\setminus\{b_n\}$. Choosing one completion in
each cylinder gives countable cofinality. There is no greatest element by
\cref{lwo:lem:extrema}. Reverse all comparisons for the lower end.
\end{proof}
\merge\ Source~07 proves the same statement by the same greedy stripping
(second half of its Lemma~8.1, 07:731--742), including the empty set.
Applied to $Y=\R$, it gives a second route to the global invariants
$\cf\W=\ci\W=\omega$ of \cref{lwo:prop:full-topology} (source~06,
Corollary~3.6). Source~06 adds that every increasing bijection
$h:\R\to\R$ induces an automorphism $e\mapsto h\circ e$ of $\W$, and that
these natural symmetries and negation need not exhaust the automorphism group
(06:325--327).

\begin{lemma}[Infinite-disagreement intervals contain large cylinders; source~06, Lemma~6.5]
\label[lemma]{lwo:sp:lem:intervalcone}
Let $e<f$ in $\W$ have maximal common prefix $s$, residual set $Y_s$, and
first entries $a=e(\delta)$, $b=f(\delta)$ as in \cref{lwo:thm:adjacency}. If
$Y_s$ is infinite then $(e,f)$ contains a cylinder with infinite residual
alphabet. In particular, such an interval contains a six-element finite
completion cylinder, with isolated internal points.
\end{lemma}
\begin{proof}
If some $c\in Y_s$ lies strictly between $a$ and $b$, the cylinder with next
value $c$ is contained in $(e,f)$ and still has infinite residual alphabet.
Otherwise consider the two end child cylinders.

The orders $\W(Y_s\setminus\{a\})$ and $\W(Y_s\setminus\{b\})$ cannot have,
respectively, a greatest and a least element: by the proof of
\cref{lwo:thm:adjacency}, that would force $Y_s$ finite. Hence either the
$a$-child has no maximum or the $b$-child has no minimum. In the first case,
\cref{lwo:sp:lem:countends} supplies a cylinder with infinite residual
alphabet wholly above $e$ in the $a$-child. In the second it supplies such a
cylinder wholly below $f$ in the $b$-child. These cylinders are within $(e,f)$.

In any cylinder with infinite residual set $B$, extend its prefix by listing
all but three elements of $B$. The remaining completion cylinder has $3!=6$
elements and is still contained in $(e,f)$. Its four internal points have
immediate neighbors and are isolated globally.
\end{proof}

\begin{theorem}[Dense isolation and cardinal invariants; source~06, Theorem~6.6]
\label{lwo:sp:thm:denseisolated}
The set $I$ of isolated points of $\W$ is dense and has cardinality
$2^\kappa$. Moreover,
\[
 \dens(\W)=\cell(\W)=\weight(\W)=2^\kappa.
\]
\end{theorem}
\begin{proof}
Consider any nonempty open interval $(e,f)$. If its first-disagreement
residual $Y_s$ is infinite, use \cref{lwo:sp:lem:intervalcone}. If $Y_s$ is
finite, $[e,f]$ is a finite chain, and every point strictly between its
endpoints has immediate neighbors. Every nonempty open set contains a
nonempty open interval because $\W$ has no endpoints. Thus $I$ is dense.

Fix three reals and take all enumerations of their complement of length
$\kappa$. There are $2^\kappa$ such prefixes, by \cref{lwo:prop:cardinality}.
Their six-element completion cylinders are disjoint and each contributes four
isolated points. Hence $|I|=2^\kappa$. Every dense set must contain every
open singleton, so $\dens(\W)\geq2^\kappa$; these singletons also witness the
same lower bound for cellularity. The usual interval base has size at most
$|\W|^2=2^\kappa$, proving the weight upper bound and all equalities.
\end{proof}
\merge\ The counts are those of \cref{lwo:prop:full-topology} (source~02)
and of source~07's Corollary~8.4 (07:798), proved there by the same reserved
three reals; the density of $I$ is source~06's alone.

\begin{proposition}[The perfect remainder and zero-dimensionality; source~06, Proposition~6.7]
\label[proposition]{lwo:sp:prop:perfect}
The first derived set $\W'=\W\setminus I$ is closed, nowhere dense, has
cardinality $2^\kappa$, and has no isolated points in its subspace topology.
Thus $\W''=\W'$. The whole space $\W$ is zero-dimensional.
\end{proposition}
\begin{proof}
The set $I$ is open and dense, so its complement is closed and nowhere dense.
All enumerations of limit length $\kappa$ are nonisolated, giving
$|\W'|=2^\kappa$.

There is a nonisolated point between any two different finite-condensation
classes. Indeed, the closed interval between representatives of distinct
classes is infinite, so its first-disagreement residual is infinite by
\cref{lwo:thm:adjacency} and convexity of finite cylinders. Apply
\cref{lwo:sp:lem:intervalcone}, and complete the resulting prefix by an
enumeration of its infinite residual set of initial-cardinal length. The
resulting total domain is a limit ordinal, so the point has singleton finite
block and is nonisolated. Inside a nonsingleton finite block, $\W'$ retains
exactly the first and last points. Neither can become isolated in $\W'$, since
its outward side has no adjacent different class. Singleton classes have the
same property on both sides. This proves $\W''=\W'$.

For zero-dimensionality, the cut at a jump yields clopen initial and final
rays. Near a nonisolated point, any side lacking an immediate neighbor has
isolated points arbitrarily close, by density of $I$; their jumps supply
such clopen cuts. A side with an immediate neighbor already has a jump cut.
Combining a lower and an upper clopen cut gives a clopen convex neighborhood
inside any prescribed neighborhood. Isolated points themselves have clopen
singleton neighborhoods.
\end{proof}
The derived-set statement is topological, not an assertion that the order is
scattered: $\W$ contains a copy of $\Q$, and much larger dense orders.

\section{The point-character spectrum of the full order}\label{lwo:sp:sec:characters}

\merge\ This answers, for $\W$, the point-character half of source~02's
research topic ``Exact spectra of cuts and point characters''
(\cref{lwo:sec:research}). Sources~06 and~07 prove it by the same
disagreement-position argument; the statement and proof below are source~06's,
and source~07's differing realizations are recorded after the proof.

For $e:\alpha\to\R$ in $\W$, define
\begin{align*}
 D^-(e)&=\{\xi<\alpha:\exists a\notin\ran(e\restr\xi),\ a<e(\xi)\},\\
 D^+(e)&=\{\xi<\alpha:\exists a\notin\ran(e\restr\xi),\ a>e(\xi)\}.
\end{align*}
These sets record positions at which an enumeration can branch below or
above $e$. Write
\[
 \chm(e)=\cf\{f:f<e\},\qquad
 \chp(e)=\ci\{f:e<f\}.
\]
Both sides are nonempty. The value $1$ means an immediate neighbor.

\begin{theorem}[Character from available disagreement positions; source~06, Theorem~8.1]
\label{lwo:sp:thm:characterformula}
If $D^-(e)$ has no largest element and $\lambda^-=\sup\{\xi+1:\xi\in D^-(e)\}$,
then
\[
 \chm(e)=\cf\lambda^-.
\]
If $D^-(e)$ has a largest element, then $\chm(e)\in\{1,\omega\}$. The identical
statements with $+$ and reversed order determine $\chp(e)$.
\end{theorem}
\begin{proof}
Every $f<e$ has first disagreement at a position in $D^-(e)$, and every such
position permits a lower completion. Later first disagreements lie strictly
closer to $e$. If the position set has no largest element, choosing one lower
completion at each position of a cofinal subset gives a cofinal family of size
$\cf\lambda^-$, and fewer than $\cf\lambda^-$ points have bounded
first-disagreement positions and miss a closer lower branch. This gives the
formula.

If the last lower-branch position is $\delta$, the sector of all points first
disagreeing at $\delta$ is cofinal below $e$. It is a lexicographic sum over
the real set of available next values below $e(\delta)$, with residual
permutation orders as summands. A next-value set with no maximum has
cofinality $\omega$; with a maximum, use the last residual summand and
\cref{lwo:sp:lem:countends}. Hence the sector has cofinality $1$ or $\omega$.
The upper-side proof is dual.
\end{proof}
\merge\ Source~07's Proposition~9.1 (07:832) is the same statement, written
$\chm(e)=\cf D^-(e)$, with the same proof. It adds when an immediate
predecessor occurs: the set $B$ of available next values below $e(\delta)$ at
the last position $\delta$ must have a greatest value, and the final fiber
over it must have a greatest enumeration.

\begin{proposition}[No unequal uncountable point characters; source~06, Proposition~8.2]
\label[proposition]{lwo:sp:prop:equalhigh}
If both $\chm(e)$ and $\chp(e)$ are uncountable, then they are equal.
\end{proposition}
\begin{proof}
Both disagreement-position sets have no largest element. Let their suprema as
in \cref{lwo:sp:thm:characterformula} be $\lambda^-$ and $\lambda^+$. Suppose
$\lambda^-<\lambda^+$. At every position at least $\lambda^-$, $e$ chooses the
least unused real. Thus this whole final part of $e$ is strictly increasing
and has countable length, by \cref{lwo:sp:lem:realmono}. Every strictly later
limit position in that tail, in particular $\lambda^+$, has countable
cofinality. This contradicts uncountability of $\chp(e)$. The reverse
inequality is handled by a decreasing tail. Therefore the suprema agree, and
so do their cofinalities.
\end{proof}

\begin{theorem}[Exact point-character spectrum; sources~06 and~07]
\label{lwo:sp:thm:point-spectrum}
Let $K=\{1\}\cup\Reg(\le\kappa)$, where $\Reg(\le\kappa)$ is the set of infinite
regular cardinals at most $\kappa$. The realized pairs $(\chm(e),\chp(e))$,
$e\in\W$, are exactly
\begin{equation}\label{lwo:sp:eq:pointpairs}
 \{(u,v)\in K^2:u=v\text{ or }\min(u,v)\leq\omega\}.
\end{equation}
In particular, every infinite regular cardinal at most $\kappa$ occurs as a
one-sided point character, and a point can have pair $(\rho,\rho)$ for every
uncountable regular $\rho\leq\kappa$, whereas a proper gap cannot have two
uncountable cofinalities (\cref{lwo:sp:thm:gaprestriction}).
\end{theorem}
\begin{proof}
The restriction to $K$ follows from \cref{lwo:lem:small-cut} (source~06's
Corollary~5.5): every nontrivial infinite cofinality of a linear order is
regular, and the sides of a point are bounded by $\kappa$.
\Cref{lwo:sp:prop:equalhigh} gives the remaining exclusion. We now realize
every allowed pair.

For infinite regular $\rho\leq\kappa$ set $\theta=\kappa\cdot\rho$, an ordinal
of cardinality $\kappa$ and cofinality $\rho$. To realize $(\rho,\rho)$,
partition $\R$ into four-element sets indexed by $\theta$. Within each set
label the values $a<b<c<d$ and use the block $\langle b,a,c,d\rangle$. The
lower branches at $b$ and upper branches at $c$ are cofinal in the resulting
limit domain $4\cdot\theta=\theta$. The character formula gives $(\rho,\rho)$.

To realize $(\rho,1)$, reserve $\{0,1\}$, enumerate the complement using
cofinally many decreasing two-element blocks over length $\theta$, and append
$\langle0,1\rangle$. The lower branch positions are cofinal in $\theta$ and
stop there. The last two-element block gives an immediate successor.
Reversing real values gives $(1,\rho)$.

To realize $(\rho,\omega)$, reserve $\N$, enumerate $\R\setminus\N$ in
decreasing pairs over length $\theta$, and append $0,1,2,\ldots$ in its usual
order. Lower branches have supremum $\theta$; upper branches are cofinal in
the final $\omega$-tail. Negation gives $(\omega,\rho)$. These constructions
include the cases $\rho=\omega$. Finally an internal point of a six-element
final block realizes $(1,1)$.
\end{proof}
\begin{proof}[Second route \textup{(source~07, Theorem~9.2, 07:861)}]
The exclusions are argued as above. For the realizations fix an infinite
regular $\mu\leq\kappa$ and put $\lambda=\kappa\cdot\mu$. Any $\kappa$-sized
subset of $\R$ can be partitioned into triples indexed by $\lambda$. List each
sorted triple as \emph{middle, least, greatest}. The flattened enumeration has
length $3\cdot\lambda=\lambda$, and at the first position of each triple both
a smaller and a larger value are still unused, so both divergence sets are
cofinal in $\lambda$, whose cofinality is $\mu$. Apply this to all of $\R$ for
$(\mu,\mu)$; to $\R\setminus\N$, followed by $\N$ in increasing order, for
$(\mu,\omega)$, and negate for $(\omega,\mu)$. For $(\mu,1)$ reserve two reals
$a<b$, use the triple construction on their complement, and append $(a,b)$;
appending $(b,a)$ instead gives $(1,\mu)$. An interior point of a six-element
block gives $(1,1)$.
\end{proof}
Source~07 notes that if $\kappa$ is singular, $\kappa$ itself is not a
character. A point's neighborhood-base cardinality is the maximum of its two
one-sided values; when both are $1$, it is an isolated point with a
one-element local base (source~06). The two-sided spectrum is substantially
more informative than this single-cardinal invariant: it detects one-sided
endpoints of finite blocks and the asymmetry of countable tails (source~07).

\section{Gaps of the full order and the size of its completion}\label{lwo:sp:sec:gaps}

\merge\ This answers, for $\W$, the cut half of source~02's research topic
``Exact spectra of cuts and point characters''. The proof is source~06's cut
spine; source~07's prefix tracking (its Lemma~10.1 and Theorem~10.2,
07:921--992) follows the same straddling cylinders, and its different
realization of the gaps is printed as a second route.

\begin{definition}[Source~06, Definition~7.1]\label[definition]{lwo:sp:def:gap}
A \emph{proper Dedekind gap} is a partition $\W=L\sqcup U$ into nonempty sets
with $L<U$, such that $L$ has no greatest element and $U$ no least element. Its
cofinality pair, or \emph{character}, is $(\cf L,\ci U)$; $\GapChar(M)$ is the
set of characters of the gaps of a linear order $M$. A jump is not a proper
gap. The Dedekind completion here adjoins one point for each proper gap and
keeps all original points and jumps; endpoints may separately be added if
desired.
\end{definition}

\begin{lemma}[Cut-spine alternative; source~06, Lemma~7.2]\label[lemma]{lwo:sp:lem:cutspine}
Every proper gap is described in one of the following two ways:
\begin{enumerate}[label=(\roman*)]
\item At an injective prefix $p$, a nonempty proper initial segment $B$ of
$Y_p$ divides the next-value child cylinders: those with next value in $B$
lie on the lower side, and the other child cylinders lie on the upper side.
\item At a limit-length injective prefix $p$, the whole cylinder $C_p$ lies on
one side and is cofinal or coinitial on that side. The gap is the lower or
upper boundary of $C_p$.
\end{enumerate}
The length of the describing prefix is below $\kappa^+$.
\end{lemma}
\begin{proof}
Say that a cylinder \emph{straddles} the gap if it meets both $L$ and $U$.
Start with the empty prefix, whose cylinder is all of $\W$. In any straddling
cylinder, at most one child cylinder straddles: the children are linearly
arranged by their next real value. If none straddles, the children partition
into a nonempty initial family contained in $L$ and a nonempty final family
contained in $U$. This is (i). If one straddles, extend the prefix by its next
value.

At a limit stage take the union of the prefixes. It is an injection and has a
completion. If its cylinder still straddles, continue. If the cylinder lies
entirely in $L$, it is cofinal in $L$. To see this, suppose a lower-side point
lay above every member of the cylinder. Its first disagreement with the union
prefix would occur at an earlier stage. The straddling cylinder immediately
after that coordinate would be wholly below this lower-side point, but would
also contain an upper-side point, a contradiction. The upper-side case is dual.
This gives (ii).

The process cannot continue for $\kappa^+$ stages, since it would then have
chosen $\kappa^+$ distinct reals. Nor can a singleton exhausted cylinder
straddle. Thus it stops at one of the stated configurations.
\end{proof}
\merge\ Source~07's Lemma~10.1 (07:921) is the straddling step and the limit
step of this proof; it adds that the first difference used at a limit exists
because a complete enumeration cannot properly end inside an injective prefix
that still extends further.

\begin{theorem}[Gap restriction; source~06, Theorem~7.3]\label{lwo:sp:thm:gaprestriction}
Every proper gap of $\W$ has countable cofinality on at least one side.
Both sides have infinite regular cofinality at most $\kappa$.
\end{theorem}
\begin{proof}
In case (i) of \cref{lwo:sp:lem:cutspine}, the lower-side part inside $C_p$
is cofinal in $L$ because $C_p$ straddles the gap. It is a lexicographic sum
over $B$ of residual permutation orders, by \eqref{lwo:sp:eq:children}. If $B$
has no maximum, countably many next values are cofinal by
\cref{lwo:sp:lem:realmono}. If $B$ has a maximum, the last child cylinder has
cofinality at most $\omega$ by \cref{lwo:sp:lem:countends}. Therefore
$\cf L\leq\omega$. The dual argument gives $\ci U\leq\omega$. Neither side has
an endpoint, so both equal $\omega$.

In case (ii), the side on which the whole cylinder is cofinal or coinitial has
cofinality at most $\omega$ by \cref{lwo:sp:lem:countends}. If that value were
$1$, the side itself would have an endpoint, contrary to properness. Thus it
is $\omega$.

The other side is bounded by $\kappa$ using \cref{lwo:lem:small-cut}. The
cofinality of an endpoint-free linear order is an infinite regular cardinal:
a cofinal set can be replaced by an increasing cofinal sequence of minimal
cardinality, and a singular indexing cardinal would admit a smaller cofinal
subsequence. Apply the same fact to the reverse upper side.
\end{proof}
Thus there is no proper $(\omega_1,\omega_1)$-gap, or any proper gap with both
cofinalities uncountable. This does not assert countable saturation:
$(\omega,\omega)$-gaps exist (\cref{lwo:prop:real-gap}), and jumps exist.

\begin{lemma}[Boundary cofinality of a partial enumeration; source~06, Lemma~7.4]
\label[lemma]{lwo:sp:lem:boundarycof}
Let $p:\theta\to\R$ be injective, with $\theta$ limit, and let
\[
 D^-(p)=\{\xi<\theta:\text{some real unused before }\xi
                                  \text{ is smaller than }p(\xi)\}.
\]
If $D^-(p)$ is cofinal in $\theta$, then the set of all full enumerations
strictly below the entire cylinder $C_p$ has cofinality $\cf\theta$. The
reversed statement holds for the upper boundary and $D^+(p)$.
\end{lemma}
\begin{proof}
Every enumeration below the cylinder first disagrees with $p$ at a coordinate
in $D^-(p)$. Points that first disagree at a later coordinate are closer to
the cylinder and greater than those disagreeing earlier. Choosing one
completion at each coordinate in a cofinal subset of $D^-(p)$ gives a cofinal
family of size $\cf\theta$. Conversely, fewer than $\cf\theta$
first-disagreement coordinates are bounded in $\theta$ and miss a closer
available lower branch.
\end{proof}

\begin{theorem}[Exact gap spectrum of $\W$; sources~06 and~07]\label{lwo:sp:thm:gapspectrum}
\[
 \GapChar(\W)=\{(\omega,\rho),(\rho,\omega):\rho\in\Reg(\le\kappa)\}.
\]
In particular a gap cannot have two uncountable sides.
\end{theorem}
\begin{proof}
\Cref{lwo:sp:thm:gaprestriction} gives the inclusion $\subseteq$. For each
infinite regular $\rho\leq\kappa$ put $\theta=\kappa\cdot\rho$, an ordinal of
cardinality $\kappa$ and cofinality $\rho$. Partition $S=(-\infty,0]$ into
pairs indexed by $\theta$, which is possible since $|\theta|=\kappa$. List
each pair in decreasing order, and concatenate the blocks. The resulting
prefix $p$ enumerates $S$ and has length $2\cdot\theta=\theta$. Its
lower-branch set $D^-(p)$ is cofinal, since the first element of each
decreasing pair leaves a smaller partner unused. The residual alphabet is
$T=(0,\infty)$.

At the lower boundary of $C_p$, the lower side has cofinality $\rho$ by
\cref{lwo:sp:lem:boundarycof}; it has no greatest element. The upper side has
$C_p\cong\W(T)$ as a coinitial suborder, and this order has no minimum and
has coinitiality $\omega$ by \cref{lwo:sp:lem:countends}. Hence this is a
proper $(\rho,\omega)$-gap. Negation reverses the pair.
\end{proof}
\begin{proof}[Second route for the realization \textup{(source~07, Theorem~10.2, 07:976--991)}]
Fix an infinite regular $\mu\leq\kappa$ and put $\lambda=\kappa\cdot\mu$.
Reserve $S=(-1,1)$, and enumerate $\R\setminus S$ in $\lambda$ positions by the
middle--least--greatest triple construction of the second route of
\cref{lwo:sp:thm:point-spectrum}. Denote this prefix by $p$ and its cylinder by
$C$. It is isomorphic to $\W(S)$ and has no greatest point, with cofinality
$\omega$. Form the cut immediately above $C$:
\[
 L=\{x:\exists c\in C\ (x\leq c)\},\qquad U=\W\setminus L.
\]
A point in $U$ first differs positively from $p$ at a coordinate below
$\lambda$. Such positive differences can be made at cofinally many triple
starts. Perturbations at a cofinal set of these starts form a coinitial family
in $U$ of size $\mu$. Fewer than $\mu$ differences are bounded and can all be
improved by a later one, so $\ci U=\mu$ and $U$ has no least point. Since $C$ is
cofinal in $L$, $\cf L=\omega$. Negation gives $(\mu,\omega)$.
\end{proof}

\begin{theorem}[Completion cardinality; source~06, Theorem~7.6]\label{lwo:sp:thm:gapcount}
There are exactly $2^\kappa$ proper gaps of $\W$. Its Dedekind completion, with
or without two added endpoints, has cardinality $2^\kappa$.
\end{theorem}
\begin{proof}
There are at most $2^\kappa$ injective prefixes. For each length
$\theta<\kappa^+$ there are at most $\kappa^{|\theta|}\leq\kappa^\kappa=2^\kappa$
functions, and there are $\kappa^+\leq2^\kappa$ lengths. A subset of $\R$ has
at most $\kappa$ initial segments: an initial segment is determined by its
supremum in the extended real line, together with whether a possible boundary
point is included. The descriptions in \cref{lwo:sp:lem:cutspine} therefore
give at most $2^\kappa$ proper gaps.

For the reverse inequality, let $p$ range over all length-$\kappa$
enumerations of $S=(-\infty,0]$. There are $2^\kappa$ choices. The
lower-branch set $D^-(p)$ is cofinal in $\kappa$: otherwise a tail of $p$
would always choose its least unused value, hence would be an increasing
sequence of $\kappa$ reals, contrary to \cref{lwo:sp:lem:realmono}. By
\cref{lwo:sp:lem:boundarycof}, the lower boundary of $C_p$ has no greatest
lower point. Its upper side has the minimum-free cylinder $\W((0,\infty))$
coinitial. Thus every such boundary is a proper gap. Distinct prefixes of the
same length have disjoint ordered cylinders and distinct lower boundaries.
This produces $2^\kappa$ gaps.

Adding one point for each proper gap to the original $2^\kappa$ points leaves
the cardinality unchanged. Two endpoints make no difference.
\end{proof}
The cardinality is also part of \cref{lwo:thm:completion} (source~02) and of
source~07's Corollary~4.6. Source~06 notes, as source~02 does, that the
completion still has no increasing or decreasing sequence of length $\kappa^+$:
such a sequence would give $\kappa^+$ separated nonempty intervals, from which
original points could be chosen in order, since every nonempty open interval
of the completion meets the original order.

Under CH the gap spectrum is
$\GapChar(\W)=\{(\omega,\omega),(\omega,\omega_1),(\omega_1,\omega)\}$
(source~07, 07:1057--1063). These are gap characters, not point characters;
conflating the two would reverse some of the conclusions.

\section{Real powers and representation rank}\label{lwo:rk:sec:rank}

\merge\ This section and \cref{lwo:rk:sec:KS,lwo:rk:sec:greedy} are
source~07's. Source~07 measures lexicographic complexity by real coordinates,
following the representability terminology of Giarlotta \cite{Giarlotta}:
\[
 \repr(L)=\min\{\alpha:L\emb\R^\alpha_{\lex}\}.
\]
For an ordinal exponent every real-valued function is allowed; there is no
finite-support restriction. The binary invariant $\bl$ of
\cref{lwo:def:binary-length} satisfies $\repr(L)\le\bl(L)\le\omega\cdot\repr(L)$
(\cref{lwo:tab:notation}).

\begin{theorem}[Representation ranks; source~07, Theorem~5.3]\label{lwo:rk:thm:ranks}
\[
 \repr(\W_\kappa)=\repr(\A)=\kappa,\qquad \repr(\W)=\kappa^+.
\]
In fact $\W$ embeds in $\B_{\kappa^+}$ as well as in $\R^{\kappa^+}_{\lex}$.
\end{theorem}
\begin{proof}
$\W_\kappa\equi\A\equi\R^\kappa_{\lex}$ (\cref{lwo:thm:fixed-equivalence} and
the note after it) and \cref{lwo:rk:thm:strict} give the first identity. If
$\W$ embedded in $\R^\gamma_{\lex}$ for some $\gamma<\kappa^+$,
\cref{lwo:cor:real-powers} would imply
$\R^{\gamma+1}_{\lex}\emb\R^\gamma_{\lex}$, a contradiction.

For the upper bound, extend each complete enumeration $e:\alpha\to\R$ by zero
to a function on $\kappa^+$. Distinct complete enumerations already differ
before either ends, so padding changes no comparison. For a binary target,
first replace every real coordinate by its $j$-code (\cref{lwo:lem:real-code})
and then pad by zeros to length $\kappa^+$. Each code length is
$\omega\cdot\alpha<\kappa^+$, and the same prefix-free argument applies.
\end{proof}
\merge\ The binary half of the upper bound, and $\bl(\W)=\kappa^+$, are
source~02's \cref{lwo:thm:full-length}.

Thus $\W$ contains every real power strictly below the critical exponent
(\cref{lwo:cor:real-powers}), but not the real power at that exponent. Its
representation rank is the critical exponent nonetheless. Ordinal spectrum and
representation rank measure genuinely different aspects of an order
(source~07, 07:536--539). For fixed-length powers source~07 adds
$\repr((\W_\kappa)^\beta_{\lex})=\repr(\A^\beta_{\lex})=\kappa\cdot\beta$ for
every nonzero ordinal $\beta$ (its Proposition~6.4 and Theorem~11.2): by the
note after \cref{lwo:cor:no-fixed-square} both powers are equimorphic with
$\R^{\kappa\cdot\beta}_{\lex}$, and \cref{lwo:rk:thm:strict} applies.

\begin{remark}[Absorption is not isomorphism; source~07, 07:628--633]
\label[remark]{lwo:rk:rem:omega-power}
Mutual embeddability must not be upgraded to isomorphism.
$\W^\omega_{\lex}$ is dense: between two tuples, increase a coordinate of
the smaller tuple strictly after their first difference, keeping all earlier
coordinates fixed. This is possible because $\W$ has no greatest element. The
full order $\W$ is not dense, by its finite blocks (\cref{lwo:thm:blocks}).
Thus $\W^\omega_{\lex}\equi\W$ (\cref{lwo:thm:absorption}) but
$\W^\omega_{\lex}\not\cong\W$.
\end{remark}

\section{More on the fixed-length order}\label{lwo:sp:sec:minimal}

\merge\ Sources~06 and~07 sharpen \cref{lwo:prop:fixed-topology}. The full
order allows prefixes that have already used $\kappa$ reals but have not
exhausted $\R$. In $\W_\kappa$, every proper prefix has length less than
$\kappa$, hence uses fewer than $\kappa$ values. This elementary difference
controls its topology (source~06, 06:956--959).

\begin{lemma}[Short residuals remain locally large; source~06, Lemma~9.1]
\label[lemma]{lwo:sp:lem:shortresidual}
If $p:\theta\to\R$ is injective and $\theta<\kappa$, then every nonempty open
real interval contains $\kappa$ elements of $Y_p$. Every such prefix has
extensions in $\W_\kappa$.
\end{lemma}
\begin{proof}
An open real interval has cardinality $\kappa$. Removing fewer than $\kappa$
points leaves cardinality $\kappa$, without any regularity assumption on
$\kappa$. Enumerate $Y_p$ with domain $\kappa$ and concatenate; since
$\theta+\kappa=\kappa$, the resulting full domain is $\kappa$.
\end{proof}
Source~07 states the same fact (07:281--286) and stresses that it is what
makes short-prefix cylinders of $\W_\kappa$ behave differently from
long-prefix cylinders of $\W$.

\begin{theorem}[Dense local universality and a clopen base; source~06, Theorem~9.2]
\label{lwo:sp:thm:minimaldense}
$\W_\kappa$ is dense without endpoints. Every nonempty open interval in
$\W_\kappa$ contains a copy of $\B_\kappa$, and hence contains every linear
order of size at most $\kappa$; in particular every such interval has
cardinality $2^\kappa$. The short-prefix cylinders
\[
 C_p\cap\W_\kappa,\qquad \dom p<\kappa,
\]
are clopen and form a base for its order topology.
\end{theorem}
\begin{proof}
Given $e<f$, their first disagreement occurs at $\delta<\kappa$. By
\cref{lwo:sp:lem:shortresidual}, an unused real lies strictly between their
values there. The corresponding next-value cylinder is inside $(e,f)$. Its
residual set has cardinality $\kappa$, so pair coding followed by the fixed
prefix embeds the binary power, still with total length $\kappa$. Varying the
first real also shows that there are no endpoints.

A short-prefix cylinder has no endpoints: after its prescribed prefix, one can
choose unused next values both below and above any given next value. The
cylinder is therefore open. If a point lies outside it, its first disagreement
with the prefix occurs at a coordinate below the prefix length; a longer
short-prefix neighborhood of that outside point remains outside. Thus the
complement is open. Finally, inside a prescribed interval around $e$, fix its
prefix beyond its first disagreements with both bounds. The resulting
cylinder lies inside the interval and proves the base assertion.
\end{proof}
\merge\ Source~07 proves the clopen base in its Theorem~7.1 and the copy of
$\B_\kappa$ in every interval in its Corollary~7.2 (07:637, 07:680), by the
same between-values construction.

\begin{proposition}[Weight and convergence; sources~06 and~07]\label[proposition]{lwo:sp:prop:minimalweight}
$\weight(\W_\kappa)=2^{<\kappa}$, so that
$\dens(\W_\kappa)=\cell(\W_\kappa)=\weight(\W_\kappa)=2^{<\kappa}$.
No sequence of distinct points of $\W_\kappa$ converges in its order topology.
\end{proposition}
\begin{proof}
The number of short prefixes is at most $\kappa^{<\kappa}=2^{<\kappa}$: for
infinite $\lambda<\kappa$, $(2^{\aleph_0})^\lambda=2^\lambda$, and the finite
lengths do not exceed the same supremum. The clopen base of
\cref{lwo:sp:thm:minimaldense} gives the weight upper bound; the lower bound is
$\cell(\W_\kappa)=2^{<\kappa}$ (\cref{lwo:prop:fixed-topology}) and
$\cell\le\dens\le\weight$.

Given a point $e$ and countably many other points, their first disagreement
coordinates with $e$ are bounded below $\kappa$, since $\cf\kappa>\omega$. A
longer short-prefix neighborhood of $e$ excludes them all. Hence a sequence of
distinct points cannot converge to $e$ (source~06, Theorem~9.3).
\end{proof}

\begin{corollary}[A $P$-space; source~07, Corollary~7.2]\label[corollary]{lwo:rk:cor:Pspace}
Intersections of fewer than $\cf\kappa$ open subsets of $\W_\kappa$ are open.
In particular $\W_\kappa$ is a zero-dimensional $P$-space, is not first
countable, and every convergent sequence is eventually constant.
\end{corollary}
\begin{proof}
At a point of an intersection of fewer than $\cf\kappa$ open sets, choose a
prefix cylinder inside each set. Their lengths have supremum below $\kappa$,
so their union is still a short prefix and its cylinder lies in the
intersection. Recall $\cf\kappa>\omega$. A $T_1$ $P$-space has no
non-eventually-constant convergent sequence: otherwise remove a countable set
of distinct nonlimit sequence values by intersecting their open complements.
This gives a neighborhood of the alleged limit that misses a subsequence. The
character computation of \cref{lwo:prop:fixed-topology} rules out first
countability, and a clopen base gives zero-dimensionality.
\end{proof}

\begin{remark}[Do not confuse global and local cofinality; source~06, Remark~9.4]
\label[remark]{lwo:sp:rem:globallocal}
Although every point of $\W_\kappa$ has uncountable two-sided character, the
whole order has countable cofinality and coinitiality: choose enumerations
whose first values tend to $+\infty$ or $-\infty$. These are different
invariants.
\end{remark}

\subsection{The gaps of the fixed-length order}
\merge\ This answers, for $\W_\kappa$, the cut half of source~02's research
topic on spectra.

\begin{theorem}[All gaps of the minimal-type stratum; source~07, Theorem~10.3]
\label{lwo:rk:thm:Pgaps}
\[
 \GapChar(\W_\kappa)=
 \{(\omega,\mu),(\mu,\omega):\mu\in\Reg(<\kappa)\}
 \ \cup\ \{(\cf\kappa,\cf\kappa)\},
\]
where $\Reg(<\kappa)$ is the set of infinite regular cardinals below $\kappa$.
\end{theorem}
\begin{proof}
Track a gap $(L,U)$ by short prefixes exactly as in \cref{lwo:sp:lem:cutspine},
but only up to length $\kappa$. Every prefix of length below $\kappa$ has a
completion in $\W_\kappa$. A split between next-value fibers again gives
$(\omega,\omega)$: short-prefix cylinders, as well as their child cylinders,
have countable cofinality and coinitiality by varying the next unused real
through a cofinal or coinitial sequence.

Suppose instead that at a limit length $\delta<\kappa$ the union-prefix
cylinder $C$ first lies wholly in $L$. It is cofinal in $L$, so
$\cf L=\omega$. The first differences of points in $U$ from this union prefix
are cofinal in $\delta$, because every earlier cylinder was straddling. Later
such differences are closer to $C$ than earlier ones. Choosing a cofinal set
of indices gives $\ci U\leq\cf\delta$, and a family of fewer than $\cf\delta$
points has bounded differences and is not coinitial. Thus the character is
$(\omega,\cf\delta)$. This second cardinal is regular and below $\kappa$. The
case $C\subseteq U$ is dual.

It remains possible that the process reaches length $\kappa$, producing an
injective sequence $p:\kappa\to\R$. If it is onto, it is a point of
$\W_\kappa$. The straddling property at every shorter stage then makes that
point an endpoint of one side, contrary to the gap assumption. If it is not
onto, there is no member of $\W_\kappa$ extending it. Every member differs
from $p$ somewhere, and comparison with this virtual branch divides
$\W_\kappa$ into the two gap sides. Downward and upward perturbations are
possible at every coordinate after fixing any shorter prefix. The
bounded-difference argument of \cref{lwo:prop:fixed-topology} gives character
$(\cf\kappa,\cf\kappa)$.

For realization of $(\omega,\mu)$ with $\mu<\kappa$ regular infinite, take any
injective prefix of length $\mu$ and cut immediately above its
$\W_\kappa$-cylinder. This cylinder has cofinality $\omega$; outside positive
perturbations approach it with coinitiality $\mu$. Negation gives the reversed
pair. Finally choose a bijection $p:\kappa\to\R\setminus\{0\}$. Its virtual
branch gives the $(\cf\kappa,\cf\kappa)$ gap just described.
\end{proof}
The last construction illustrates what the full order adds: $p\cat(0)$ is a
genuine point of $\W$ lying in that previously unfilled cut. This fills the gap
by allowing a longer \emph{injective} enumeration, not by introducing
repetitions. Longer prefixes also create the mixed gaps of
\cref{lwo:sp:thm:gapspectrum} whose uncountable side can be any regular
cardinal at most $\kappa$ (source~07, 07:1040--1044). Under CH,
$\GapChar(\W_\kappa)=\{(\omega,\omega),(\omega_1,\omega_1)\}$.

\begin{corollary}[No convex copy; source~07, Corollary~10.4]\label[corollary]{lwo:rk:cor:noconvex}
Although $\W_\kappa\subseteq\W$, there is no order embedding of $\W_\kappa$ into
$\W$ with convex image.
\end{corollary}
\begin{proof}
$\cf\kappa>\omega$. The $(\cf\kappa,\cf\kappa)$ gap of $\W_\kappa$ would induce
a gap with the same character in $\W$ if its image were convex: any point
between its two sides would have to be in the image.
\Cref{lwo:sp:thm:gapspectrum} forbids such a gap.
\end{proof}

\subsection{How longer enumerations fill gaps of the minimal slice}
\merge\ Source~06's Section~10 relates the full order and its minimal-type
slice beyond shared embedding spectra. It uses the order $\Jinj$ of all
injections $\kappa\to\R$; source~06 calls it $J$, and the letter is changed
here because source~07's different $J$ appears in \cref{lwo:rk:sec:greedy}.

\begin{definition}[Source~06, Definition~10.1]\label[definition]{lwo:sp:def:Jinj}
Let $\Jinj$ be the lexicographically ordered set of all injections
$p:\kappa\to\R$, without requiring surjectivity. Thus $\W_\kappa\subseteq\Jinj$.
\end{definition}

\begin{theorem}[Canonical lexicographic-sum decomposition; source~06, Theorem~10.2]
\label{lwo:sp:thm:inflation}
The first-$\kappa$-coordinates map gives an order isomorphism
\begin{equation}\label{lwo:sp:eq:inflation}
 \W\cong\sum_{p\in\Jinj}\W(\R\setminus\ran p),
\end{equation}
with $\W(\varnothing)$ interpreted as a singleton. Further, $\W_\kappa$ is
order-dense in $\Jinj$, and $\Jinj$ and $\W_\kappa$ are mutually
order-embeddable.
\end{theorem}
\begin{proof}
Every full enumeration has length at least $\kappa$. Its restriction to the
first $\kappa$ positions is in $\Jinj$, and deleting those positions
identifies the corresponding fiber with the residual permutation order.
Distinct prefixes compare at a coordinate below $\kappa$, so the fibers are
arranged exactly as the lexicographic sum in \eqref{lwo:sp:eq:inflation}.

Given $p<q$ in $\Jinj$, choose an unused real strictly between their values at
the first disagreement and complete the resulting short prefix within
$\W_\kappa$. This gives a point of $\W_\kappa$ between them. Finally,
$\W_\kappa\subseteq\Jinj\subseteq\R^\kappa_{\lex}$, and the last order embeds
into $\W_\kappa$ (\cref{lwo:thm:fixed-equivalence,lwo:lem:real-code}).
\end{proof}

\begin{proposition}[The missing exhaustive branch; source~06, Proposition~10.3]
\label[proposition]{lwo:sp:prop:balancedP}
Every $p\in\Jinj\setminus\W_\kappa$ determines a proper gap of $\W_\kappa$ with
cofinality pair $(\cf\kappa,\cf\kappa)$. Inside $\W$, the set of all points
strictly between the two sides of this gap is exactly $C_p$.
\end{proposition}
\begin{proof}
No member of $\W_\kappa$ agrees with $p$ on all $\kappa$ positions, since $p$ is
not onto. Thus comparison with $p$, as functions of the same length,
partitions $\W_\kappa$ into a lower and upper part. The available lower and
upper branch positions of $p$ are cofinal in $\kappa$: a bound on either would
make a $\kappa$-long tail of $p$ monotone in $\R$. Every short branch can be
completed in $\W_\kappa$. The disagreement-coordinate argument therefore
gives cofinality $\cf\kappa$ on both sides, with no endpoints.

Every member of $C_p$ lies above all lower-side members of $\W_\kappa$ and
below all upper-side members. Conversely, a full enumeration not extending $p$
first disagrees with it at some $\xi<\kappa$. A permitted disagreement on the
same side at a later coordinate gives a member of $\W_\kappa$ strictly between
that enumeration and $C_p$. Hence the enumeration does not lie between
\emph{all} lower-side and \emph{all} upper-side members.
\end{proof}
This identifies an actual mechanism behind the absence of uncountable balanced
gaps in the full order: some such gaps in the minimal slice are replaced by
entire residual permutation orders, not merely by one point. The
decomposition is exact; source~06 does not claim that every gap of
$\W_\kappa$ comes from an injective missing branch, and horizontal
countable-cofinality gaps remain a separate phenomenon (06:1122--1127).
\merge\ \Cref{lwo:rk:thm:Pgaps} (source~07) shows that the missing branches
give the only uncountable balanced gaps of $\W_\kappa$; its other gaps have
one countable side.

\section{The dense quotient and its connected completion}\label{lwo:sp:sec:quotient}

\merge\ This answers source~02's question whether $\Wq$ is equimorphic with
$\W$ (\cref{lwo:sec:research}). Part~(a) is source~06's Theorem~11.1; part~(b)
is source~07's Proposition~10.5 beyond~(a).

\begin{theorem}[Removing finite permutation blocks; sources~06 and~07]\label{lwo:sp:thm:quotient}
\leavevmode
\begin{enumerate}[label=(\alph*)]
\item The order $\Wq$ is dense without endpoints, has cardinality $2^\kappa$,
and has exactly the same ordinal embedding spectrum as $\W$. In fact $\Wq$ and
$\W$ are mutually order-embeddable. Its Dedekind completion is a linear
continuum of cardinality $2^\kappa$, with no endpoints and with no copy of
$\kappa^+$ or $(\kappa^+)^{\rev}$.
\item $\dens(\Wq)=\cell(\Wq)=\weight(\Wq)=2^\kappa$, and
$\GapChar(\Wq)=\GapChar(\W)$.
\end{enumerate}
\end{theorem}
\begin{proof}[Proof of \textup{(a)} \textup{(source~06)}]
The classes are nonempty and finite, so their number is $2^\kappa$. There are
no first or last classes, since an endpoint class would have an endpoint in
$\W$. If two distinct classes had no class strictly between them, their finite
union would be convex, contradicting maximality of the classes. Thus $\Wq$ is
dense. Choosing one representative from each class embeds $\Wq$ in $\W$.

For the reverse embedding choose an increasing bijection $h:\R\to(0,1)$ and a
fixed enumeration $t:\kappa\to\R\setminus(0,1)$. For $e:\alpha\to\R$, let
$E(e)=(h\circ e)\cat t$. This is a full enumeration, and its domain
$\alpha+\kappa$ is a limit ordinal of cardinality $\kappa$. Distinct $e,f$
disagree before either finishes, by \cref{lwo:lem:prefix-free}; hence applying
$h$ preserves their first disagreement, and appending $t$ does not alter it.
Thus $E$ is an order embedding into $\W$. Every $E(e)$ has a singleton
finite-condensation class because its domain is a limit ordinal. Composing
with the quotient map is consequently an order embedding $\W\emb\Wq$.

A proper gap of $\Wq$ lifts to a proper gap of $\W$ by taking unions of its
classes. Neither lifted side can gain an endpoint, since that endpoint's class
would be an endpoint on the corresponding side of $\Wq$. Therefore $\Wq$ has at
most $2^\kappa$ proper gaps. Its completion has cardinality $2^\kappa$, since
$\Wq$ already does. Being a complete dense order, it is a linear continuum; the
usual order-topological connectedness proof applies. For clarity, if a
separation of a nontrivial interval existed, the supremum of the part on one
side would violate openness of one of the separating sets. Finally, a
$\kappa^+$-long monotone sequence in this completion would yield one in its
dense suborder $\Wq$ by choosing intervening points, and hence one in $\W$,
which is impossible.
\end{proof}
\begin{proof}[Proof of \textup{(b)}, and a second route to the equimorphism \textup{(source~07)}]
Source~07 embeds $\Wq$ into $\W$ by the least point of each block, and $\W$
into $\Wq$ by $e\mapsto(h\circ e)\cat q\cat n$, with $h:\R\to(-2,-1)$
increasing, $q$ an enumeration of $\R\setminus((-2,-1)\cup\N)$ and $n$ the
increasing enumeration of $\N$; every image has a limit domain, so lies in a
singleton block.

For the density lower bound, fix $S=[0,1]$. There are $2^\kappa$ distinct
prefixes of type $\kappa$ enumerating $\R\setminus S$. Their disjoint convex
cylinders each contain at least three distinct singleton finite blocks,
obtained from limit-length completions. After projection each contains a
nonempty open interval. These intervals witness cellularity $2^\kappa$;
cardinality gives the weight upper bound as before.

A gap cannot cut through a finite block, because a nontrivial split of a
finite block would supply endpoints. Conversely a gap of the quotient lifts to
a gap of $\W$. Infinite cofinalities on either side are unchanged by replacing
points by finite blocks. This proves the spectrum assertion.
\end{proof}
Adding two endpoints to this completion gives a compact connected ordered space
of the same cardinality. This connected object is different from the
completion of $\W$ itself: the latter retains jumps. Finite condensation is an
essential step in the connected construction (source~06).

The quotient and $\W^\omega_{\lex}$ (\cref{lwo:rk:rem:omega-power}) are both
dense and mutually embeddable with $\W$, but they are not isomorphic to each
other. Every point of $\W^\omega_{\lex}$ has character $(\omega,\omega)$, by
perturbing later and later factors. The quotient retains, for example, the
$(\omega_1,\omega_1)$ character of a singleton limit-domain block constructed
in \cref{lwo:sp:thm:point-spectrum} (source~07, 07:1103--1108).

\section{Bounded surreal sign orders}\label{lwo:sp:sec:surreal}

\merge\ This section is source~06's. The repository's sign-sequence
foundations suggest a useful comparison, but the termination convention
differs. A surreal sign sequence is compared using
\[
 -\ <\ \text{termination}\ <\ +
\]
at the first place where the sequences cease to agree. In contrast,
exhaustive enumerations are prefix-free, so termination never participates in
their comparison. For an ordinal $\theta$, define
\[
 S_{<\theta}=\bigcup_{\alpha<\theta}\{-,+\}^\alpha,
\]
with this sign order. One may view it as the order of surreal numbers whose
sign-sequence lengths, or birthdays, are below $\theta$; no field structure is
used here.

\begin{theorem}[Every bounded sign order embeds, but their full union does not; source~06, Theorem~12.1]
\label{lwo:sp:thm:surreals}
For every $\theta<\kappa^+$, $S_{<\theta}\emb\W$. However,
\[
 S_{<\kappa^+}\not\emb\W,
 \qquad |S_{<\kappa^+}|=|\W|=2^\kappa.
\]
\end{theorem}
\begin{proof}
Pad a sign sequence $s$ of length $\alpha<\theta$ to length $\theta$ by using
the alphabet $\{0,1,2\}$: write $0$ for $-$, $2$ for $+$, and $1$ at all
coordinates from $\alpha$ onward. At the first disagreement this implements
exactly the sign/termination convention, so it is an order embedding into
$(\{0,1,2\}^\theta,<_{\lex})$. Replace the symbols $0,1,2$ by the binary
blocks $00,01,10$. This embeds into $\B_{2\cdot\theta}$, which embeds into
$\W$ because $|2\cdot\theta|\leq\kappa$ (\cref{lwo:lem:all-small-cubes}).

For each $\alpha<\kappa^+$, the all-plus sequence of length $\alpha$ is
strictly below the all-plus sequence of any greater length. Thus the full
union contains $\kappa^+$, excluded by \cref{lwo:thm:no-long}. Its cardinality
is $2^\kappa$: the length-$\kappa$ level gives the lower bound, and there are
$\kappa^+$ levels, each of size at most $2^\kappa$.
\end{proof}
This is a concrete warning against a tempting union argument. Every bounded
stage embeds, yet the set-sized union of all stages below $\kappa^+$ does not,
even though its cardinality is no larger than the target's. The issue is the
order structure of the union, not cardinality or a proper-class paradox. The
construction above supplies no coherent family of embeddings whose union would
be an embedding of the full sign order, and the theorem shows that such
coherence is impossible.

The sign-simplicity module inspected by source~06,
\nolinkurl{Algebra/SurrealNumbers/Surreal/Foundations/SignSequenceSimplicity.lean},
proves uniqueness of the minimum-birthday element of a nonempty convex set of
sign sequences and separates that statement from existence of cut separators
\cite{signsimplicity}. That distinction is relevant here: prefix and simplicity
arguments alone do not establish cut filling. The cut-filling and non-filling
statements of this report instead come from exhaustive completion and the cut
spine. \merge\ The module is context only: none of its declarations states a
result of this report.

\noindent[Added 2 October 2026, batch~80: for an infinite cardinal $\theta$
the surreal-well-orders report takes $S_{<\theta}$ itself as the alphabet and
computes the binary length of its minimal slice; see the dated note after the
question ``Other ground orders'' in \cref{lwo:sec:research}.]

\section{The Kanovei--Shelah index: further topology, gaps and fibres}\label{lwo:rk:sec:KS}

\merge\ Source~07's Section~11 studies the index $\A$ of \eqref{lwo:eq:KS} with
the binary characteristic-function order on $\mathcal P(\N)$. Its embedding
type (source~07's Lemma~11.1 and Theorem~11.2) is
\cref{lwo:thm:fixed-equivalence} and \cref{lwo:rk:thm:ranks}, and its
cofinality, coinitiality, density, cellularity and point characters are
\cref{lwo:prop:A-topology}. The formal definition includes principal
ultrafilters; the embedding and topological results remain true when the
family is restricted to free ultrafilters, with the changes indicated below.
Thus a difference in raw enumeration rules does not change the embedding
class at length $\kappa$, but it does change the order-isomorphism type.

\begin{proposition}[Base, weight and the free variant; source~07, Theorem~11.3]
\label[proposition]{lwo:rk:prop:Atopology}
The order topology of $\A$ has a clopen base of admissible prefixes of length
below $\kappa$, and $\weight(\A)=2^{<\kappa}$. The free-ultrafilter variant of
$\A$ also has coinitiality $\omega$. Intersections of fewer than $\cf\kappa$
open subsets of $\A$ are open, and every convergent sequence in $\A$ is
eventually constant.
\end{proposition}
\begin{proof}
Fix $a\in\A$ with range $U$. Arbitrarily late coordinates have value strictly
below $\N$: otherwise the range would have cardinality below $\kappa$. At any
such coordinate replace the value by $\N$, and complete the remaining
coordinates so that the range is still $U$. This gives arbitrarily late upward
perturbations. There are also arbitrarily late downward perturbations. If $U$
has no least member in its binary order, any coordinate value has a smaller
member of $U$. If $U$ has a least member, the tail cannot be constantly equal
to it, again because $|U|=\kappa$. Choose a late nonleast value and replace it
by a smaller member of $U$. The unused coordinates have order type $\kappa$,
so a full surjection onto $U$ can always be appended after a short altered
prefix. The same argument as in \cref{lwo:sp:thm:minimaldense} gives clopen
prefix cylinders and a base; there are at most $2^{<\kappa}$ short prefixes,
and $\weight\ge\cell=2^{<\kappa}$ by \cref{lwo:prop:A-topology}.

For coinitiality of the free variant, replace the singletons $\{n\}$ of the
proof of \cref{lwo:prop:A-topology} by the cofinite sets $[n,\infty)\cap\N$
and complete to any fixed free ultrafilter. The $P$-space assertions follow as
in \cref{lwo:rk:cor:Pspace}. These topological facts do not supply
order-theoretic saturation.
\end{proof}

\begin{proposition}[Two explicit gaps in the ultrafilter index; source~07, Proposition~11.4]
\label[proposition]{lwo:rk:prop:Agaps}
The order $\A$ has gaps of characters $(\omega,\omega)$ and
$(\cf\kappa,\cf\kappa)$. In particular it is not an $\eta_1$-order.
\end{proposition}
\begin{proof}[Proof \textup{(the $(\omega,\omega)$ gap is a second route to \cref{lwo:prop:A-gap})}]
Fix a free ultrafilter $U$. Let $a_m$ begin with $\N\setminus\{m\}$, and let
$b_n$ begin with $\N,[n,\infty)\cap\N$, completing both to surjections onto
$U$. The $a_m$ increase, the $b_n$ decrease, and $a_m<b_n$ for every $m,n$. A
point between them all would have to have first value $\N$ and second value
$\varnothing$. The latter is in no ultrafilter. More explicitly, a first value
below $\N$ lies below some $\N\setminus\{m\}$; if the first value is $\N$, a
nonempty second value lies above some $[n,\infty)\cap\N$. Thus the two
countable families really determine complementary gap sides, both with
countable character.

For the other gap, put $v=\N\setminus\{0\}$ and consider the constant virtual
branch $p(\xi)=v$ for $\xi<\kappa$. Its singleton range is not an ultrafilter,
so every member of $\A$ differs from it somewhere. Every short constant prefix
has extensions on both sides: at the next coordinate use
$\N\setminus\{0,1\}<v$ or $\N>v$ and then enumerate $U$. Arbitrarily late
deviations give cofinal and coinitial families of size $\cf\kappa$; fewer
deviations have bounded indices and are not sufficient. Comparing with $p$
therefore gives a $(\cf\kappa,\cf\kappa)$ gap.
\end{proof}

\begin{remark}[The labelling fibres are not convex; source~07, 07:1243--1247]
\label[remark]{lwo:rk:rem:fibres}
The fibers of the labelling map $a\mapsto\ran(a)$ are not convex. For
example, let $U_0,U_1$ be principal at $0,1$ respectively. Two enumerations
of $U_0$ can start with $\{0\}$ and $\N$, while an enumeration of $U_1$ starts
with $\{0,1\}$ and lies between them. Thus identifying all repeated copies of
each ultrafilter is not an ordered quotient of $\A$ by convex classes.
\end{remark}

\section{A well-ordering-based ultrafilter index}\label{lwo:rk:sec:greedy}

\merge\ This section is source~07's Section~12. It connects the literal
well-ordering interpretation to ultrafilters directly, as a separate
construction, not as a claim about the notation or exact index of
\cite{KS}. Source~07 writes $J$ for the order below; it is $\JU$ here, because
source~06's $J$ is $\Jinj$ (\cref{lwo:sp:def:Jinj}).

Let
\[
 \JU=\{e:\kappa\to\mathcal P(\N):e\text{ is bijective}\},
\]
with the characteristic-function alphabet and the first-difference order.
Every $e\in\JU$ canonically determines a free ultrafilter $U_e$. Begin with
the cofinite filter. At stage $\xi<\kappa$, adjoin $e(\xi)$ and close under
the filter operations if this remains a proper filter; otherwise leave the
filter unchanged. At limit stages take unions.

\begin{theorem}[Source~07, Theorem~12.1]\label{lwo:rk:thm:greedyultra}
The resulting $U_e$ is a free ultrafilter, is definable from $e$, and every
free ultrafilter on $\N$ equals $U_e$ for some $e\in\JU$. Moreover
$\JU\equi\R^\kappa_{\lex}$.
\end{theorem}
\begin{proof}
If adjoining a set $X$ to a filter $F$ makes it improper, some $Y\in F$ has
$Y\cap X=\varnothing$, so $\N\setminus X\in F$ already. Thus every set or its
complement is present after that set has been considered. Increasing unions at
limits are proper filters because every finite intersection already occurs at
some earlier stage. At the end the filter decides every set and is an
ultrafilter. It contains every cofinite set, so cannot be principal. All
choices during this recursion are determined by $e$, giving definability from
that parameter.

Fix a free ultrafilter $U$. Order the complementary pairs in type $\kappa$,
and in each pair list the $U$-member first and its complement second. The
resulting enumeration has domain $2\cdot\kappa=\kappa$. Inductively the current
filter stays inside $U$. Its next $U$-member can always be adjoined properly,
and that member prevents adjoining its complement at the following step.
Consequently the final filter is $U$.

Pair switches on the alphabet $\mathcal P(\N)$ give $\B_\kappa\emb\JU$, just
as in \cref{lwo:prop:cardinality}; the ternary coordinate map $T$ of
\cref{lwo:lem:real-code} gives $\JU\emb\R^\kappa_{\lex}$, and
$\R^\kappa_{\lex}\equi\B_\kappa$ proves the last assertion.
\end{proof}
This supplies a definable linearly ordered family of labels covering every
free ultrafilter, without choosing a preferred global well-order. It is a
literal well-order-enumeration counterpart of the repeated-index idea.

In the published finite-support construction, a finite set of labels is sorted
before applying its ultrafilter quantifiers; no well-order of the entire label
set is required. A further $\omega_1$-stage iteration supplies countable
saturation \cite[\S\S2--4]{KS}. The gaps proved here belong to index orders,
not to that final hyperreal field. An embedding or even an isomorphism of
unlabelled index orders need not preserve their ultrafilter labels or the
resulting ultrapower construction.

\section{Foundational scope and interpretation}\label{lwo:sec:scope}

\subsection{A definable order is not a selected well-order}
The set of all well-order relations on $\R$ is definable from $\R$, and
first-difference comparison can be expressed without choosing one member
of that set as privileged. This observation does not produce a definable
well-ordering of $\R$. It defines a linear order \emph{on the collection}
of all its well-orderings. Likewise, forming all the presentations in
\eqref{lwo:eq:KS} does not select one ultrafilter or one presentation.

Our proofs nevertheless use Choice. It supplies well-orderings of the
real line, partitions into pairs, selected enumerations of residual sets,
and simultaneous choices of embeddings. In ZF alone, $\R$ need not be
well-orderable, so the literal collection $\W$ may be empty. Even when
that collection is nonempty, the arguments using families of choices
should not automatically be read as proofs in a weaker theory.

The set-sized nature of the construction is also essential. Every possible
length is below $\kappa^+$, and every well-ordering is a subset of
$\R^2$. There is no proper class of elements hidden in $\W$, despite its
unbounded range of lengths below $\kappa^+$. In contrast, the phrase
``all ordinals'' denotes a proper class and does not describe a suborder
of this set.

\subsection{Which aspects are genuinely independent?}
The proofs separate four kinds of information. The ordinal spectrum records
how far a single monotone well-ordered chain can run. Binary length measures
the depth needed to distinguish all points through ordered binary cuts.
Adjacency records what happens when only finitely many ground elements
remain. Finally, topology and cut cofinalities describe local approach and
missing limits. These quantities interact, but none should be silently
substituted for another.

For example, $\W_\kappa$ and $\A$ have identical embedding classes yet
unequal global cofinalities. The full order and its completion have the
same ordinal spectrum and binary length, but differ in completeness.
The strata $\W_{\kappa+2}$ and $\W_{\kappa+3}$ have different binary
lengths and different finite blocks, though their cardinalities and ordinal
spectra coincide. Under some cardinal-arithmetic assumptions even density
fails to distinguish those last two orders; their isolated points still do.

All embeddings in this article are purely order-theoretic. An order
embedding need not be continuous for the two order topologies. In
particular, the fact that a cube or a complete order embeds into a dense
order does not transfer its order-topological compactness to the ambient
space. No topological or model-theoretic consequence is inferred from an
embedding without a separate argument.

\subsection{Choice and cardinal arithmetic in sources~06 and~07}\label{lwo:sec:choice}
\merge\ Source~06 (its Section~13.1) and source~07 (its Section~13.1) make the
same points independently, and add the following. The order comparison is
defined from canonical ordinal enumerations and the given real order; it does
not choose a preferred well-ordering of $\R$. ZFC is used to ensure that these
enumerations exist, to choose pair partitions and background completions,
and to use the regularity of the successor cardinal in the stabilization
proof. Once a background well-order of an alphabet is fixed, every one of its
subsets has a specified completion by restriction, which avoids a fresh choice
at every prefix (source~06). The no-long-chain argument applies to every
infinite well-orderable alphabet; it does not rely on the completeness or
density of $\R$ (source~07).

Neither source proves that full AC is necessary for every theorem, and
replacing ZFC by a particular weaker choice hypothesis would require a
separate audit of the cardinal and cofinality arguments. In ZF the defining
formula for the family of well-order relations on $\R$ still makes sense, but
its nonemptiness says that $\R$ is well-orderable. A definable \emph{family} is
not a definable \emph{chosen member}; the absence of a least point in the
lexicographic order explains why taking ``the lexicographically least
well-order'' cannot perform such a selection. No assertion here gives a Borel
real code for a well-order of all reals, and no least choice principle
sufficient for every theorem has been identified (source~07).

No continuum hypothesis, generalized continuum hypothesis, forcing axiom, large
cardinal, or regularity assumption on $\con$ is used. In particular
$|\W|=2^\con$ must not be replaced by $\con^+$ without an additional
cardinal-arithmetic hypothesis. Under CH the ordinal barrier is $\omega_2$ and
$\dens(\W_\kappa)=\omega_1$, while $|\W|=2^{\omega_1}$ can still require a
separate cardinal-arithmetic specification (source~06).

\section{Further research questions and topics}\label{lwo:sec:research}

The following are questions proposed by this investigation. We have not
verified that every one is open in the published literature. The solved
fixed-stratum \emph{equimorphism} problem is deliberately excluded from the
list; the next issues concern sharper structure and broader universality.

\merge\ Source~06's seven and source~07's seven research questions are merged
into source~02's ten topics; each is printed once, under the topic it shares,
citing every source that asks it. Neither source~06 nor source~07 advertises
its questions as catalogued open problems: source~06 ties each to a proved
limitation or construction of its own text, and source~07 states them so that
a solution does not repeat a result already proved. Source~06's Research
question~14.1 (``Unconditional comparison with one fixed lexicographic
power'') is answered in ZFC by \cref{lwo:rk:cor:notP,lwo:cor:unbounded} and is
not repeated. Source~06 had suggested ``an ordinal-valued
lexicographic-dimension invariant rather than cardinality alone'' as the
approach, warning that any argument must respect the known limitations of
converses based solely on absence of long monotone sequences
\cite{Ramakrishnan}; the two answers use exactly such invariants, $\bl$ and
$\repr$. The topic on spectra of cuts and point characters is re-scoped
below, and the quotient part of the second topic is answered. Source~07 names
the square-isomorphism question (\cref{lwo:rk:q:square}) as the most direct
next mathematical target and the ordinal-spectrum theorem as the most direct
verification target.

\subsection{An intrinsic embedding criterion for the full order}
Find a necessary and sufficient condition on a set-sized linear order $L$
for $L\emb\W$. We have an exact answer when $\bl(L)\leq\kappa$ and
exact answers for the binary and real-coordinate powers, but not for every
order of cardinality $2^\kappa$ and binary length $\kappa^+$.
Absence of $\kappa^+$ and its reverse is necessary; this article does not
prove it sufficient.

A concrete route is to study trees of convex partitions whose branches
consume distinct symbols from an alphabet of size $\kappa$. The exhaustion
argument shows what a forbidden branch would do. The remaining problem is
to characterize when a compatible coding can be organized across all
branches, not merely along each branch separately.

\merge\ Sources~06 and~07 ask the same question.

\begin{question}[An intrinsic embedding criterion for large orders; source~06, Research question~14.2]
\label[question]{lwo:sp:q:criterion}
Which linear orders of cardinality between $\kappa$ and $2^\kappa$ embed
into $\W$?
\end{question}
Source~06 has several sufficient classes: all orders of size at most $\kappa$,
all binary and real powers of fixed length below $\kappa^+$, and orders with
small separating families (\cref{lwo:sp:rem:separating}). Necessary conditions
include the absence of $\kappa^+$ and its reverse. A useful characterization
should explain the additional role of an exhaustion-compatible tree
representation. Neither necessity nor sufficiency beyond what is proved should
be assumed.

\begin{question}[An exact criterion for embedding in $\W$; source~07, Research question~14.1]
\label[question]{lwo:rk:q:criterion}
Characterize the linear orders $L$ that embed in $\W$. The necessary
conditions
\[
 |L|\leq2^\kappa,\quad \kappa^+\not\emb L,\quad (\kappa^+)^{\rev}\not\emb L,
 \quad \repr(L)\leq\kappa^+
\]
are not shown to be sufficient. Every order with representation rank
\emph{strictly below} $\kappa^+$ does embed (\cref{lwo:cor:real-powers}), so the
unresolved part of this criterion is the boundary rank $\kappa^+$. Seek either
a sufficient intrinsic prefix-tree condition or a counterexample satisfying
these necessary bounds.
\end{question}

\subsection{Can the full order absorb its completion?}
Decide whether
\[
                         \widehat\W\emb\W.
\]
The affirmative direction would make the two orders equimorphic, since the
other embedding is automatic. The invariants in \cref{lwo:thm:completion} do
not obstruct it, and the power-absorption theorem suggests considerable
room for coding cuts. A proof must nevertheless assign codes to different
cuts coherently. Independently choosing a cofinal sequence for each cut
need not preserve their lexicographic comparison.

A related problem asks whether the dense finite-condensation quotient
$\W/{\sim_{\mathrm{fin}}}$ is equimorphic with $\W$, and whether its
completion has a more recognizable presentation.

\merge\ \emph{Answered in part.} The quotient is equimorphic with $\W$
(\cref{lwo:sp:thm:quotient}, sources~06 and~07), and its completion is a
linear continuum of cardinality $2^\kappa$. A recognizable presentation of the
completion, and whether $\widehat\W\emb\W$, remain open; sources~06 and~07 ask
the following.

\begin{question}[Internal structure of the dense quotient and its completion; source~06, Research question~14.5]
\label[question]{lwo:sp:q:quotient}
Can the dense finite-condensation quotient $\Wq$, or its connected completion,
be characterized by a natural uniqueness theorem?
\end{question}
Cardinality, the ordinal spectrum, and the gap spectrum alone need not
characterize a linear order. The canonical injection-indexed decomposition
\eqref{lwo:sp:eq:inflation} suggests stronger local invariants. One can also
ask which compact connected ordered spaces arise from replacing the residual
fibers by their own dense quotients before completion.

\begin{question}[A canonical description of the completions; source~07, Research question~14.4]
\label[question]{lwo:rk:q:completions}
Give a normal form for every cut using the prefix-tracking proof, including
its first stopping stage and remaining alphabet. The cardinalities and gap
characters of the completions are known, but an explicit isomorphism-level
description is not. Compare the completion of $\W_\kappa$ with expansions of
its non-surjective terminal branches into full residual-enumeration cylinders.
The latter process must also account for the new mixed gaps of $\W$.
\end{question}

\subsection{Isomorphism types beyond equimorphism}
Determine the order-isomorphism types of individual strata and their
natural quotients. In particular, equimorphism of $\W_\lambda$ and
$\W_{\lambda+1}$ does not settle whether they are isomorphic. Their
presentations impose different exhaustion conditions, and an isomorphism
need not preserve the represented ordinal length.

The comparison theorem \cref{lwo:cor:strata-comparison} supplies a complete
baseline: the task is no longer to decide which fixed strata embed, but
to identify additional invariants capable of distinguishing equimorphic
ones. Cut spectra, automorphism orbits, and local characters are plausible
candidates.

\merge\ The cut and character spectra of $\W$ and $\W_\kappa$ are now known
(\cref{lwo:sp:thm:point-spectrum,lwo:sp:thm:gapspectrum,lwo:rk:thm:Pgaps}).
Source~07 asks the analogous question for the powers of $\W$.

\begin{question}[Isomorphism rather than mutual embeddability; source~07, Research question~14.2]
\label[question]{lwo:rk:q:square}
Determine whether $\W^2_{\lex}\cong\W$. More generally classify the nonzero
successor exponents $\beta<\kappa^+$ for which $\W^\beta_{\lex}\cong\W$. The
countable limit power is already excluded by density
(\cref{lwo:rk:rem:omega-power}), and all these powers are mutually embeddable
with $\W$. Finite-block size, global cofinality, and representation rank alone
do not settle the square.
\end{question}

\subsection{Automorphisms and recovery of the prefix structure}
Every increasing bijection $h:\R\to\R$ induces an automorphism
$e\mapsto h\circ e$ of $\W$. Determine the full automorphism group and
whether some canonical fragment of the prefix tree can be reconstructed
from the external linear order alone.

Finite-condensation classes and their sizes are intrinsic and therefore
preserved. It is less clear how much of the non-finite ordinal remainder
of an enumeration can be recovered. One useful intermediate target is a
classification of the orbits of endpoints and interior points of the
finite blocks, together with the singleton-block points.

\merge\ Sources~06 and~07 ask the same question.

\begin{question}[Automorphisms and recoverable structure; source~06, Research question~14.3]
\label[question]{lwo:sp:q:automorphisms}
Which features of an enumeration are recoverable from the abstract order $\W$,
and what is its full automorphism group?
\end{question}
Finite-condensation class size and one-sided cofinalities are order
invariants. Ordinary increasing bijections of $\R$ give automorphisms, but
that construction alone does not show that the ordinal length of an
enumeration is invariant under every automorphism. A first concrete target is
the orbit structure of singleton finite blocks with equal point characters.

\begin{question}[Automorphisms and finite condensation; source~07, Research question~14.3]
\label[question]{lwo:rk:q:automorphisms}
Describe the order-automorphism group of $\W$ and of its dense finite-block
quotient. Increasing bijections of $\R$ act on enumerations, and negation
reverses the order, but these need not exhaust the possibilities. Are block
size, position inside the block, and the two external point characters enough
to determine an orbit? Determine which information about an enumeration's
ordinal domain can be recovered from the unlabelled order structure.
\end{question}

\subsection{Exact spectra of cuts and point characters}
\merge\ \emph{Re-scoped.} For $\W$ this topic is answered by sources~06
and~07: the point-character spectrum is \eqref{lwo:sp:eq:pointpairs}
(\cref{lwo:sp:thm:point-spectrum}) and the gap spectrum is
$\{(\omega,\rho),(\rho,\omega):\rho\in\Reg(\le\kappa)\}$
(\cref{lwo:sp:thm:gapspectrum}). For $\W_\kappa$ the point characters are
constant (\cref{lwo:prop:fixed-topology}) and source~07 computes the gap
spectrum (\cref{lwo:rk:thm:Pgaps}). For $\A$ the point characters are constant
(\cref{lwo:prop:A-topology}), but only two of its gap characters are known
(\cref{lwo:prop:A-gap,lwo:rk:prop:Agaps}). The topic therefore stays open for
the longer strata $\W_\beta$ with $\kappa<\beta<\kappa^+$, for the gap
spectrum of $\A$, and for the convex-suborder criteria of the second
paragraph below. Source~02's original text follows.

Classify all pairs of left and right cofinalities occurring at points and
unrealized cuts of $\W$ and its fixed strata. We have computed the local
pair $(\cf\kappa,\cf\kappa)$ throughout $\W_\kappa$, exhibited missing
$(\omega,\omega)$ cuts, and bounded every side of a full-order cut by
$\kappa$. Those facts leave many cases undecided.

A productive decomposition may separate three sources of cuts: a missing
next coordinate in an ordered residual set, a limit of increasingly long
injective prefixes, and exhaustion at the end of a finite block. This
could also lead to exact criteria for convex suborders to be complete,
separable, or first countable.

\merge\ The cut-spine proof of \cref{lwo:sp:lem:cutspine} realizes source~02's
proposed decomposition: case~(i) is the missing next coordinate in an ordered
residual set, case~(ii) the limit of increasingly long injective prefixes.

\subsection{Other ground orders}
The ordinal cutoff and full binary length were proved for every infinite
linearly ordered ground set $X$. The all-strata classification used
$X\subseteq\R$, particularly countable separating cuts and rational
coordinate codes. Determine the weakest hypotheses on $X$ under which
\[
           \W_{\lambda+n}(X)\equi\B_\lambda\times_{\lex}n!
\]
still holds. For orders with large intrinsic binary length, develop a
replacement formula that records both coordinate complexity and the
shrinking size of the residual alphabet.

This question naturally includes nonseparable real-closed fields and
set-sized suborders of surreal numbers. The distinction between the
underlying order and any algebraic structure must remain explicit.

\merge\ Sources~06 and~07 ask the same question.

\noindent[Added 2 October 2026, batch~80: answered in part for one family of
set-sized suborders of surreal numbers, at the minimal stratum $n=0$. The
surreal-well-orders report,
\nolinkurl{Algebra/SurrealNumbers/docs/foundations-and-computation/surreal-well-orders/},
merged from batch-80 manuscripts 08, 09, 11 and~12, takes $X=S_{<\theta}$ for
an infinite cardinal $\theta$, so that $\lambda=|S_{<\theta}|=2^{<\theta}$
(it writes $\mathbf{No}_{<\theta}$; the order is the sign order of
\cref{lwo:sp:sec:surreal}). Its Theorem~8.3 (label
\texttt{swo:st:thm:transition}, source~12's) proves
$\W_\lambda(S_{<\theta})\equi\B_\lambda$, which is the case $n=0$ of the
display above, except when $\theta$ is a singular strong-limit cardinal. Then
$\lambda=\theta$ and $\W_\theta(S_{<\theta})\equi\B_{\theta\cdot\theta}$, with
the ordinal product, so $\bl(\W_\theta(S_{<\theta}))=\theta\cdot\theta>\theta$
and the display fails for $n=0$. Its Example~8.5 (\texttt{swo:st:n10.5}) takes
$\theta=\beth_\omega$, a singular strong limit in ZFC alone. The question stays
open for $n>0$, for the longer strata, for cutoffs $\theta$ that are not
cardinals, for other alphabets, and for the weakest hypotheses on $X$. The same
report reprints, with pointers to the labels here and no novelty claim, its
sources' re-proofs of results that this report proves for every infinite linear
alphabet: \cref{lwo:lem:prefix-free,lwo:prop:cardinality,lwo:lem:cut-code,lwo:thm:cube-strict,lwo:lem:eventual,lwo:thm:no-long,lwo:thm:ordinal-spectrum,lwo:lem:all-small-cubes,lwo:thm:full-length,lwo:lem:extrema,lwo:thm:adjacency,lwo:thm:blocks,lwo:sp:thm:universality}.
No result here uses that report.]

\noindent[Added 3 October 2026, batch~81: answered further for the same
alphabets. The surreal-well-orders report now also merges batch-81 manuscripts
01--06 (its sources~13--18). Write $\mu=|S_{<\theta}|=2^{<\theta}$ (the
$\lambda$ of the preceding note). For every stratum $\lambda=\mu\cdot\gamma$
with $0<\gamma<\mu^+$ and every $n<\omega$, its source~15 proves
$\W_{\lambda+n}(S_{<\theta})\equi\B_{\rho\cdot\gamma}\times_{\lex}n!$, where
$\rho=\mu$ unless $\theta$ is a singular strong-limit cardinal, and
$\rho=\theta\cdot\theta$ (ordinal product) when it is (Theorem~31.2 there,
\texttt{swo:cb:new:thm:block-types}). Hence the display above holds for all
these strata when $\theta$ is not a singular strong limit. At a singular strong
limit it holds exactly when $\theta\cdot\theta\cdot\gamma=\theta\cdot\gamma$,
that is, when $\gamma$ is a left multiple of $\theta^\omega$, and fails
otherwise, for instance for $\gamma=1,2,\theta$ (that report's Remark~31.8,
\texttt{swo:vi:rem:lwo-display}); so the display is not simply false at
singular strong limits. The question stays open for strata that are not
multiples of $\mu$ (the first is $\mu+\omega$), for cutoffs $\theta$ that are
not cardinals, for other alphabets, and for the weakest hypotheses on $X$. No
result here uses that report.]

\begin{question}[Other alphabets and exact gap spectra; source~06, Research question~14.4]
\label[question]{lwo:sp:q:alphabets}
For a general infinite linear alphabet $X$, how is the gap spectrum of $\W(X)$
determined by the cuts and monotone sequences of $X$?
\end{question}
The cardinality, universality, and ordinal spectrum are already settled for
all such alphabets (\cref{lwo:thm:ordinal-spectrum,lwo:sp:thm:universality}).
The real-specific gap result uses that extremal stripping has countable length
and residual permutation orders have countable ends. Replacing those
properties by explicit cofinality bounds should yield a general cut-spine
transfer theorem. Dense nonseparable alphabets, well-ordered alphabets, and
ordered fields with large cofinality provide informative tests.

\noindent[Added 3 October 2026, batch~81: for minimal slices a transfer
theorem of this kind is proved in the surreal-well-orders report by its
source~18 (Theorem~30.10 there,
\nolinkurl{swo:sh:thm:general-permutation-gap-spectrum}). If $D$ is a dense
linear order without endpoints of size $\mu$, all of whose nonempty open
intervals have size $\mu$, whose points all have two-sided character $\nu$, and
whose cofinality and coinitiality are $\nu$, then the gap characters of the
order of bijections $\mu\to D$ are those of $D$ together with $(\nu,\nu)$,
$(\operatorname{cf}\mu,\operatorname{cf}\mu)$, and $(\rho,\nu)$, $(\nu,\rho)$
for every infinite regular $\rho<\mu$. For $D=S_{<\theta}$ its Theorem~30.16
(\texttt{swo:sh:thm:gap-multiplicities}) counts them: $2^\mu$ gaps of
character $(\operatorname{cf}\mu,\operatorname{cf}\mu)$ and $\mu^{<\mu}$ of
every other character, the analogue for that alphabet of
\cref{lwo:sp:thm:gapcount}. The question for the full order $\W(X)$ stays
open.]

\begin{question}[Other alphabets and other lengths; source~07, Research question~14.5]
\label[question]{lwo:rk:q:alphabets}
Extend the detailed classification from $\R$ to subsets of $\R$ and then to
arbitrary linear alphabets. The no-$\nu^+$ theorem already works for any
infinite alphabet of cardinality $\nu$. The endpoint, density, and rank
arguments require additional structure. Countable alphabets deserve a separate
treatment: for example $\omega\cdot\omega\ne\omega$, so the real-to-binary
block flattening used at the uncountable initial ordinal $\kappa$ cannot
simply be copied at length $\omega$.
\end{question}
\merge\ For subsets $X\subseteq\R$, source~02's \cref{lwo:thm:all-strata}
already classifies every stratum up to equimorphism, including countable $X$
through \cref{lwo:lem:countable-alphabet}.

\noindent[Added 3 October 2026, batch~81: the density part is answered for
every infinite linear alphabet. For $|X|=\mu$ and $|\alpha|=\mu$, source~16 of
the surreal-well-orders report proves that the density and cellularity of
$\W_\alpha(X)$ are $\mu^{<\mu}$ for $\alpha\in\{\mu,\mu+1,\mu+2\}$ and $2^\mu$
otherwise, and $2^\mu$ for $\W(X)$ (Theorem~27.2 there,
\texttt{swo:ds:thm:density}; the case $X=\R$ is
\cref{lwo:prop:all-strata-density}), and that every infinite stratum, and
$\W(X)$, is conditionally incomplete (Theorem~28.1,
\texttt{swo:ds:thm:incomplete}). The gap spectrum of the minimal slice is
computed for $X=S_{<\theta}$ (its source~17, Theorem~30.11) and for regular
$\eta_\mu$ alphabets of size $\mu$ (its source~16, Theorem~30.14). Endpoints and
binary rank for other alphabets, and the gap spectra of the longer strata, stay
open.]

\subsection{Ultrafilter fibers and the hyperreal presentation}
Among nonprincipal ultrafilters $U$, determine whether the order-isomorphism
type of $\A_U$ depends on $U$, even though its embedding class does not.
Principal and nonprincipal fibers are already distinguished by their
coinitialities, so an undifferentiated claim that all fibers are isomorphic
would be false.

One can also ask which transformations of the index preserve a specified
finite-support ultrapower presentation, rather than merely providing an
order embedding of indices. The index order by itself does not encode all
ultrafilter data used by the construction; equimorphism of index orders is
not a proof of isomorphism of the associated nonstandard models.

\merge\ Source~07 asks the same question.

\begin{question}[Ultrafilter labels and changes of index; source~07, Research question~14.6]
\label[question]{lwo:rk:q:labels}
Find useful conditions under which an embedding between labelled linear
indices induces an elementary embedding, or an isomorphism, between the
associated finite-support ultrapowers. The maps must respect the ultrafilter
labels and the order of the finite quantifier iterations, not just the
underlying order type. The greedy labelling of $\JU$
(\cref{lwo:rk:thm:greedyultra}) supplies a concrete test case against the
actual repeated-range index $\A$.
\end{question}

\subsection{Weak choice and uniform definability}
Determine the choice principles needed for the separate results: existence
of a nonempty exhaustive order, the successor-cardinal chain obstruction,
universality for small orders, optimal fixed-stratum coding, and power
absorption. Distinguish existence of arbitrary embeddings from embeddings
definable uniformly from their input orders.

Herzberg, Kanovei, Katz, and Lyubetsky develop a modification of the
hyperreal construction with transfer under countable choice and a further
well-ordering hypothesis for properness \cite{HKKL}. Their work provides
motivation for separating the choice strength of the index-order analysis
from that of the nonstandard model. No equivalence with their precise
axiomatic hypotheses is asserted here.

\merge\ Sources~06 and~07 ask the same question.

\begin{question}[Choice strength; source~06, Research question~14.6]
\label[question]{lwo:sp:q:choice}
Which conclusions remain valid in ZF plus the assertion that $\R$ is
well-orderable, and which require stronger choice or regularity principles?
\end{question}
The existence of a background well-order already supplies all subset
completions of $\R$. The more delicate issues are the successor-cardinal
stabilization argument, sizes of collections of prefixes, and extraction of
cofinal sequences. A dependency-by-dependency analysis is preferable to a
blanket claim that choice is either indispensable or irrelevant.

\begin{question}[Choice strength and a verified core; source~07, Research question~14.7]
\label[question]{lwo:rk:q:choice}
Determine the exact weak-choice assumptions needed for the main positive
codings and for prefix completion over $\R$. In parallel, formalize the
no-successor-sized-chain theorem and the pair-switch embedding first; these
would give a small, auditable formal core with a complete answer to the
ordinal-embedding part of the problem. Treat the gap spectra as a later,
separate verification milestone.
\end{question}

\subsection{Cardinal arithmetic and stability of the classification}
The embeddability theorems are stated in ZFC, but their numerical
interpretation depends on the continuum. Investigate which finer
isomorphism invariants depend on $2^{<\kappa}$, $2^\kappa$, or
$\cf(\kappa)$. For instance, the density distinction between the first
three real strata and later ones is strict exactly when
$2^{<\kappa}<2^\kappa$; the adjacency and isolation distinctions do not
require that inequality.

A model-comparison study should specify whether the real set, the
collection of its well-orderings, or only selected cardinal parameters are
preserved. Keeping the reals fixed does not by itself keep every subset of
$\R^2$, and hence every member of $\W$, fixed.

\subsection{A proof-assistant development for ProveIt}
A formalization can be divided into reusable modules. First formalize
prefix-free ordinal words and their first-difference order. Then establish
cut coding, the full-cylinder obstruction, and the successor-cardinal
stabilization argument. On that foundation, formalize ordinal block
concatenation and finite-fiber capacity, followed by the cardinal induction
for all fixed strata. Adjacency and finite condensation form a largely
independent final module.

The most delicate obligations are not numerical. They are the existence of
common tail bounds at limit stages, correct ordinal rather than cardinal
multiplication, and preservation of a full cylinder under transfinite
fusion. The finite verification program accompanying this article checks
none of these infinitary obligations and should not be presented as a
substitute for them.

\merge\ Sources~06 and~07 propose formalization plans of their own.
Source~06 (its Section~13.4) suggests a relation-based carrier, a structure
consisting of a well-order relation on $\R$ from which the ordinal enumeration
is derived, or alternatively an ordinal-indexed bijection carrier with an
explicit bound on its domain. It orders the work as follows: the prefix-free
lemma, first-disagreement transitivity, cylinder convexity and the pair
encoder; then initial-segment coding and the regular-cardinal stabilization
lemma; the exhaustion barrier as the main transfinite-recursion milestone;
for the real-specific theory, countability of monotone real sequences, greedy
extremal stripping, the jump criterion and disagreement-position
cofinalities; and the cut-spine recursion last, followed by its counting and
cofinality corollaries. The existing sign-sequence modules are useful
architectural examples, but copying a sign-order definition without changing
its termination semantics would be incorrect, and the repository's
distinction between a mathematical report and a checked theorem must be
preserved \cite{reports,signsimplicity}. Source~07 (its Section~13.3) first
isolates four reusable components---prefix-free ordinal sequences, order
embeddings by the first difference, regular-cardinal tail stabilization, and
completion of injective prefixes---which suffice for the exact ordinal
spectrum; a second stage adds the coordinate-fusion proof and pair and block
codings for representation ranks and product absorption; the topology comes
later: the endpoint lemma, adjacency, and finally cut tracking with explicit
successor and limit stopping cases. No unspecified theorem about
``lexicographic universality'' should substitute for the actual coding, and
no assumed continuity should enter proofs about arbitrary order embeddings.
Both are proof plans, not claims that those steps have been executed.

\begin{question}[Formalizing the cut spine; source~06, Research question~14.7]
\label[question]{lwo:sp:q:formal}
Can one produce a reusable, kernel-checked theorem for lexicographic leaf
orders of prefix trees with no branch longer than their resource alphabet?
\end{question}
The proof should separate three inputs: linear ordering of child sets, closure
of compatible prefixes under unions, and eventual exhaustion. Specializing it
to injective real sequences would yield the main gap theorem; specializing the
stabilization lemma would yield the ordinal barrier. This is a more reusable
formalization target than encoding only the particular real-number instance.
Source~07's \cref{lwo:rk:q:choice} proposes the verified ordinal core.

\section{Conclusions of sources 06 and 07}\label{lwo:sec:conclusions}
\merge\ Source~02 has no concluding section; sources~06 and~07 close as follows.

\emph{Source~06.} The lexicographic order of all well-orderings of $\R$ is
simultaneously large and sharply bounded. It has $2^\con$ elements and embeds
every continuum-sized linear order, yet it cannot contain the
successor-continuum ordinal. Its exact ordinal spectrum is therefore determined
in ZFC without any continuum hypothesis. The exhaustion requirement also
creates structure invisible in a bare lexicographic power. Finite residual
alphabets create factorial-sized jump blocks and a dense set of isolated
points. Continuing beyond the initial cardinal length fills certain
uncountable balanced gaps of the minimal slice by whole residual orders.
Nevertheless countable-sided gaps remain, and the cut-spine analysis
classifies all possible gap cofinalities. The contrast between points, gaps,
the minimal slice, and the dense quotient shows why none of these objects
should be substituted for another without an explicit argument.

\emph{Source~07.} The lexicographic order of all well-orderings of $\R$ is not
remotely a well-order, but its well-ordered suborders have a sharp ceiling: it
contains every ordinal of cardinality at most the continuum and none larger.
Within that ceiling it has unexpectedly large embedding capacity, reaching
every real lexicographic exponent below $\kappa^+$ and absorbing all
corresponding powers of itself. At the same time its local structure records
exhausted finite tails, long divergence patterns, and a rigid restriction on
unfilled cuts: one side of every gap is countable. The fixed-length stratum
and the actual ultrafilter index share the ordinal ceiling but have smaller
representation rank and different topologies. Keeping these orders separate
makes the connection with definable hyperreals more precise and prevents
cardinality, saturation, and isomorphism from being conflated with mere
mutual embeddability.

\appendix
\section{Proof dependencies and reproducible finite checks}\label{lwo:app:checks}

The central arguments have a short dependency structure:
\[
\begin{gathered}
 \text{prefix-freeness} + \text{eventual constancy}
       \ \Longrightarrow\ \text{exact ordinal cutoff},\\
 \text{full-child lemma} + \text{fusion}
       \ \Longrightarrow\ \text{strict cube hierarchy},\\
 \text{pair swaps} + \text{cut coding}
       \ \Longrightarrow\ \text{small-order universality},\\
 \text{countable codes} + \text{cardinal induction}
       \ \Longrightarrow\ \text{all fixed-stratum equivalences},\\
 \text{disjoint real copies} + \text{exhaustive concatenation}
       \ \Longrightarrow\ \text{full power absorption}.
\end{gathered}
\]
These arguments are self-contained apart from elementary ZFC facts about
well-ordering, cardinal arithmetic, and transfinite recursion, which are
explained at their points of use.

\merge\ Source~06's dependency map (its Section~13.2) adds the chains
\[
\begin{gathered}
 \text{extremal enumerations}\ \Longrightarrow\ \text{finite-tail jump classification},\\
 \text{greedy extremal stripping} + \text{cut spine}
   \ \Longrightarrow\ \text{gap restriction and completion size},\\
 \text{disagreement positions}
   \ \Longrightarrow\ \text{point characters and gap realizations},
\end{gathered}
\]
and notes that the comparison with surreal sign orders depends only on binary
coding and the ordinal barrier. Source~07's audit gives the analogous
dependencies of its representation ranks (real-to-binary coding, real-power
strictness, fixed-length padding of prefix-free codes), of its fixed-length
gaps (prefix tracking up to $\kappa$, with special care for non-surjective
branches of length exactly $\kappa$) and of its index results (surjectivity
onto an ultrafilter of size $\kappa$, arbitrarily late perturbations, explicit
missing branches); no hyperreal saturation claim is inferred from index
topology.

The archive contains \texttt{finite\_checks.py}, a Python program using only
the standard library, together with its actual output in
\texttt{finite\_checks.json}. Run it with
\begin{quote}\ttfamily
python finite\_checks.py --output finite\_checks.json
\end{quote}
It checks pair-swap order preservation, cut coding of finite labeled
orders, every candidate adjacent pair of permutations through size six,
consecutive permutation pairs through size eight, finite concatenation,
and the integer capacities used in the terminal-factor formulas. The supplied run passes \textbf{552,450} exact checks. Counts
and parameter bounds are recorded by the program, rather than inferred
from an informal estimate.

\merge\ In this report the program is shipped as
\nolinkurl{code/02-strata-finite_checks.py} and its recorded output as
\nolinkurl{data/02-strata-finite_checks.json}; the text of the program still
names its delivered output \texttt{finite\_checks.json}. Its default output
path is \texttt{finite\_checks.json} in the current directory, so a rerun
should pass an explicit scratch path, for example
\nolinkurl{py code/02-strata-finite_checks.py --output <scratch>/finite_checks.json},
and never the shipped \texttt{data/} file. A rerun on a copy during placement
reproduced the recorded JSON, identical after removing the carriage returns
that Python writes in text mode on Windows. The finite checks of sources~06 and~07
are shipped as \nolinkurl{code/06-spectra-finite_checks.py} and
\nolinkurl{code/07-rank-finite_checks.py}, with their recorded outputs in
\texttt{data/}; they too test finite instances only.

\merge\ \emph{Source~06's checks} (its Section~13.3) use only the Python standard
library. They check the pair encoder on all binary inputs for small block
counts, verify convexity and factorial cardinalities of finite prefix
cylinders, compare the jump criterion with consecutive lexicographic
permutations, and exhaustively check the converse jump criterion for all
pairs of permutations up to a stated finite size. The recorded run passed 511
binary encodings (with 43,435 first-disagreement comparisons), 16,072 finite
prefix cylinders, 46,225 consecutive-pair checks through alphabet size eight,
and 266,272 all-pair comparisons through alphabet size six, and checked
initial-segment coding for 5,914 finite linear orders through size seven. These
checks can catch indexing, orientation, and finite-suffix mistakes. They do
not verify the transfinite stabilization proof, the cut-spine argument,
cardinal arithmetic, or any claim about an uncountable order.

\emph{Source~07's checks} test finite permutation adjacency (46,225 adjacent
pairs, 7,432 exhaustive pairs and 30,000 pairs sampled with a recorded fixed
seed), pair-switch coding (2,047 words), initial-segment cut coding (5,913
coordinate orders), concatenation of interleaved ordered copies (258 product
elements), and a nonconvex ultrafilter fibre on a three-point universe
(\cref{lwo:rk:rem:fibres}). They do not check transfinite induction, choice,
cardinal arithmetic, topology, or the gap classification.

\begin{resultbox}
\textbf{What the checks establish.} They are exact finite tests of the
stated constructions and permutation formulas. They do not verify a
transfinite theorem, establish a literature-priority claim, or constitute
a Lean or Rocq formalization. The mathematical evidence for the infinite
claims is the written proof, subject to ordinary review.
\end{resultbox}

\section{Source audit and manuscript status}\label{lwo:app:sources}

The Kanovei--Shelah definition was checked directly in the paper, including
the displayed index and its comparison rule on page~2 of the arXiv PDF.
\merge\ Source~07 locates the same definition in \cite[\S2, p.~160]{KS},
on the printed page of the published version (its audit, item~1), whereas
source~02 cites Section~1 of the paper; both quote the same index and the
same comparison rule. The section numbering was not rechecked at the merge.
The paper's bibliographic information was cross-checked against Shelah's
official paper archive. The weak-choice follow-up was checked through its
authors' arXiv record. Knoblauch's work is cited for the related research
terminology and bibliographic context, not as an uninspected source for any
specific theorem proved here.

The ProveIt source was inspected at the immutable commit
\begin{center}\small\ttfamily
6ea60e3677bd847a00baf023c3c532c83884530e.
\end{center}
Its relevant report is the binary-cube and point-separation manuscript
cited below. Only its elementary cylinder perspective is reused; the
required lemma is reproved. No existing formal verification of the new
article's theorems is claimed, and this work does not modify the repository.

The literature search was targeted rather than exhaustive. It identified
related primary sources, but it does not settle historical priority for the
classification and absorption statements. The questions in
\cref{lwo:sec:research} are proposed next steps, not an authenticated list of
published open problems. All websites and repository material were
consulted for this manuscript on 2 October 2026.

\merge\ Sources~06 and~07 record their own audits in the shipped files
\nolinkurl{06-spectra-research_audit.md} and \nolinkurl{07-rank-source_audit.md},
both pinned to the same ProveIt revision and both dated 2 October 2026. Neither
claims priority, peer review or machine verification, and both describe their
repository search as targeted rather than exhaustive. In summary:
\begin{itemize}
\item Source~06 inspected the Cardinals README, the report index
\cite{reports} and the module \cite{signsimplicity} of the repository
\cite{repo}; it read but did not compile the module, and uses it as context
only. It cites Moore \cite[Section~2]{Moore} for the characteristic-function
representation, Ramakrishnan \cite{Ramakrishnan} for the literature on
lexicographic representation (including Fleischer's paper and its
correction, whose original source~06 did not obtain), and Tao \cite{Tao} for
background on ordinals; consequently it does \emph{not} infer an
arbitrary-order embedding theorem from the absence of long monotone
sequences. No nontrivial unverified assertion from a repository report is used
as a premise.
\item Source~07 inspected the Kanovei--Shelah definition on the printed page
of the published version, Kuhlmann's Corollary~2.4 on a rendered preprint page
(the parsed text being garbled; the preprint's support convention permits all
functions for an ordinal exponent), and Giarlotta's first page for the
terminology of the representability number \cite{Giarlotta}. In the
repository it read the Cardinals README \cite{ProveItCardinals}, which separates
unformalized reports from the Lean library, and a recorded result
\nolinkurl{lexicographic_tails_result.json} of the cofinal-strata report, which
concerns a different lexicographic sum and is not used. It separates
published input (the index definition, Kuhlmann's theorem, the terminology)
from its own deductions, which are not attributed to the cited papers merely
because those papers motivated the definitions.
\end{itemize}

\section{Theorem crosswalk of the three sources}\label{lwo:app:crosswalk}
\merge\ Every numbered statement of the three sources, by topic. Source
numbers are those of the delivered PDFs (T theorem, L lemma, P proposition, C
corollary, D definition, E example, R remark, Q research question), followed
by the line of the delivered \texttt{.tex}. ``Here'' is where the statement is
printed in this report. Unnumbered statements of a source are cited by
line.

{\small
\renewcommand{\arraystretch}{1.12}
\begin{longtable}{@{}>{\raggedright\arraybackslash}p{.24\textwidth}>{\raggedright\arraybackslash}p{.165\textwidth}>{\raggedright\arraybackslash}p{.165\textwidth}>{\raggedright\arraybackslash}p{.165\textwidth}>{\raggedright\arraybackslash}p{.135\textwidth}@{}}
\caption{\merge\ Crosswalk.}\label{lwo:tab:crosswalk}\\
\toprule
Result & Source 02 & Source 06 & Source 07 & Here\\
\midrule
\endfirsthead
\toprule
Result & Source 02 & Source 06 & Source 07 & Here\\
\midrule
\endhead
\bottomrule
\endfoot
Main results at a glance & T2.1 (191) & T1.1 (102) & T2.1 (187) & \ref{lwo:thm:main-spectrum}\\
Definition; prefix freeness; linear order & L3.1 (269) & L2.1 (171), D2.2, P2.3 (198) & D1.1 (118), §1.2 & \ref{lwo:lem:prefix-free}\\
Cylinders and residual sets; next-value sum & D3.2, L3.3 (297) & D3.1, L3.2 (246) & L3.1 (259) & \ref{lwo:lem:cylinder}, \eqref{lwo:sp:eq:children}\\
Cardinality $2^\lambda$ by pair swaps & P3.4 (318) & L4.1 (331), C4.2 (363) & L4.1 (329), C4.6 (417) & \ref{lwo:prop:cardinality}\\
Cut (initial-segment) coding & L4.1 (337) & L4.3 (373) & L4.2 (339) & \ref{lwo:lem:cut-code}\\
Full child; strict cube hierarchy & L4.2, T4.3 (367) & --- & --- & \ref{lwo:thm:cube-strict} (= PSG \texttt{cor:strict})\\
Binary coding length $\bl$ & D4.4 (393) & --- & --- & \ref{lwo:def:binary-length}\\
Real and binary codes & L4.5 (402) & in P4.6 & (5.1)--(5.3) & \ref{lwo:lem:real-code}\\
Cardinal identities; $\cf\kappa>\omega$ & 02:425, R11.3 & in T9.3 & L3.2 (288) & \eqref{lwo:eq:omega-kappa}, \ref{lwo:prop:fixed-topology}\\
Strict real powers (Kuhlmann) & --- & --- & T5.2 (477) & \ref{lwo:rk:thm:strict}\\
Eventual constancy & L5.1 (437) & L5.1 (454) & L4.3 (355) & \ref{lwo:lem:eventual}\\
No $\nu^+$-chain & T5.2 (451), R5.3 & T5.2 (465), R5.3 & T4.4 (367) & \ref{lwo:thm:no-long}\\
Universality; ordinal spectrum & T5.4 (487) & C5.4 (501), E4.5 & L4.2, C4.5 (394) & \ref{lwo:thm:ordinal-spectrum}\\
Universality in every slice $\W_\beta(X)$ & --- & T4.4 (389) & --- & \ref{lwo:sp:thm:universality}\\
Separating families of initial segments & --- & 06:415--423 & --- & \ref{lwo:sp:rem:separating}\\
$\W_\kappa\equi\B_\kappa\equi\A$ & T6.1 (516), C6.2 & P4.6 (425) & (5.3), L11.1, T11.2 (1129), R1.2 & \ref{lwo:thm:fixed-equivalence}\\
Finite tails $\W_{\kappa+n}$; capacity & T6.3, T6.4, E6.5 & --- & --- & \ref{lwo:thm:finite-tail}, \ref{lwo:thm:finite-capacity}\\
Small cubes and real powers below $\kappa^+$ & L7.1 (665), C7.4 (732) & L4.1, P4.6 & P5.1 (460) & \ref{lwo:lem:all-small-cubes}, \ref{lwo:cor:real-powers}\\
Coding length $\bl(\W)=\kappa^+$ & T7.2 (682) & --- & --- & \ref{lwo:thm:full-length}\\
Representation ranks $\repr$ & --- & --- & T5.3 (505) & \ref{lwo:rk:thm:ranks}\\
$\W$ into no fixed power below $\kappa^+$ & C7.3 (719) & C9.5 (1053), Q14.1 & C5.4 (525) & \ref{lwo:cor:unbounded}, \ref{lwo:rk:cor:notP}, \ref{lwo:sp:rem:conditional}\\
All fixed strata & L8.1, T8.2 (782), C8.3--C8.4, E8.5, P8.6 & --- & --- & \ref{lwo:thm:all-strata}--\ref{lwo:prop:all-strata-density}\\
Disjoint copies; power absorption; closure & L9.1, T9.2 (946), C9.3 & --- & L6.1, T6.2 (563), C6.3 & \ref{lwo:thm:absorption}, \ref{lwo:cor:closure}\\
Fixed-length powers & C9.4 (1005) & --- & P6.4 (613) & \ref{lwo:cor:no-fixed-square}\\
$\W^\omega_{\lex}\equi\W$ but not $\cong$ & --- & --- & 07:628--633 & \ref{lwo:rk:rem:omega-power}\\
Convex copies of $\W$ & P9.5 (1026) & --- & --- & \ref{lwo:prop:convex-copies}\\
Extremal enumerations & L10.1 (1044) & L3.3 (263) & L8.1 (714) & \ref{lwo:lem:extrema}\\
Countable monotone real sequences; countable ends & --- & L3.4, L3.5 (294), C3.6 & L8.1 (714) & \ref{lwo:sp:lem:realmono}, \ref{lwo:sp:lem:countends}\\
Adjacency criterion; example & T10.2 (1062), E10.3 & E6.1, T6.2 (548), App.\ B & T8.2 (747) & \ref{lwo:thm:adjacency}\\
Finite blocks; isolated points & D10.4, T10.5 (1129) & P6.3 (590), C6.4 & T8.3 (770) & \ref{lwo:thm:blocks}\\
Dense isolated points; perfect remainder & --- & L6.5, T6.6 (650), P6.7 & C8.4 (798) & \ref{lwo:sp:lem:intervalcone}--\ref{lwo:sp:prop:perfect}\\
Global invariants of $\W$ & P11.1 (1166) & C3.6, T6.6 & C8.4 & \ref{lwo:prop:full-topology}\\
Topology of $\W_\kappa$ & P11.2 (1204), R11.3 & L9.1, T9.2, T9.3, R9.4 & T7.1 (637), C7.2 & \ref{lwo:prop:fixed-topology}, \ref{lwo:sp:lem:shortresidual}--\ref{lwo:sp:rem:globallocal}\\
Finite-tail strata topology & P11.4 (1264) & --- & --- & \ref{lwo:prop:finite-topology}\\
Topology of $\A$ & L12.1, P12.2 (1330) & --- & T11.3 (1163) & \ref{lwo:prop:A-topology}, \ref{lwo:rk:prop:Atopology}\\
$(\omega,\omega)$ gaps of $\W$, $\W_\kappa$, $\A$ & P13.1 (1384), P13.2 & --- & E7.3 (703), P11.4 & \ref{lwo:prop:real-gap}, \ref{lwo:prop:A-gap}\\
Small cuts; completion; complete cubes & L13.3, T13.4, L13.5, R13.6 & C5.5, T7.6 (824) & C4.6 & \ref{lwo:lem:small-cut}--\ref{lwo:lem:cube-complete}\\
Gap spectrum of $\W$ & topic (1635) & D7.1, L7.2, T7.3, L7.4, T7.5 (802) & L10.1, T10.2 (941) & \ref{lwo:sp:def:gap}--\ref{lwo:sp:thm:gapcount}\\
Point-character spectrum & topic (1635) & T8.1, P8.2, T8.3 (913) & P9.1, T9.2 (861) & \ref{lwo:sp:thm:characterformula}--\ref{lwo:sp:thm:point-spectrum}\\
Gaps of $\W_\kappa$; no convex copy & --- & P10.3 (1101) & T10.3 (996), C10.4 & \ref{lwo:rk:thm:Pgaps}, \ref{lwo:rk:cor:noconvex}, \ref{lwo:sp:prop:balancedP}\\
Lexicographic-sum decomposition & --- & D10.1, T10.2 (1078) & --- & \ref{lwo:sp:def:Jinj}, \ref{lwo:sp:thm:inflation}\\
Dense quotient $\Wq$ & question (1604) & T11.1 (1133) & P10.5 (1072) & \ref{lwo:sp:thm:quotient}\\
Bounded surreal sign orders & --- & T12.1 (1198) & --- & \ref{lwo:sp:thm:surreals}\\
Further gaps and fibres of $\A$ & --- & --- & P11.4 (1218) & \ref{lwo:rk:prop:Agaps}, \ref{lwo:rk:rem:fibres}\\
Well-ordering-based ultrafilter index & --- & --- & T12.1 (1265) & \ref{lwo:rk:thm:greedyultra}\\
Research questions & 10 topics (1573) & Q14.1 answered; Q14.2--Q14.7 & Q14.1--Q14.7 & \ref{lwo:sec:research}\\
Choice; dependency maps; checks & §14, App.\ A, B & §13 & §13 & \ref{lwo:sec:scope}, \ref{lwo:sec:choice}, \ref{lwo:app:checks}, \ref{lwo:app:sources}\\
Conclusions & --- & §15 & Conclusion & \ref{lwo:sec:conclusions}\\
Summary tables & Table 1 & App.\ A & §2 table & Tables \ref{lwo:tab:invariants}, \ref{lwo:sp:tab:objects}\\
\end{longtable}
}

\begin{thebibliography}{9}
\raggedright
\addcontentsline{toc}{section}{References}

\bibitem{KS}
Vladimir Kanovei and Saharon Shelah,
\emph{A definable nonstandard model of the reals},
Journal of Symbolic Logic \textbf{69} (2004), no.~1, 159--164.
\href{https://doi.org/10.2178/jsl/1080938834}{doi:10.2178/jsl/1080938834}.
\href{https://arxiv.org/abs/math/0311165}{arXiv:math/0311165}.
\href{https://shelah.logic.at/papers/825/}{Author-maintained paper archive, Sh:825}.

\bibitem{Knoblauch}
Vicki Knoblauch,
\emph{Lexicographic preference representation: Intrinsic length of linear
orders on infinite sets},
Journal of Mathematical Economics \textbf{105} (2023), article 102823.
\href{https://doi.org/10.1016/j.jmateco.2023.102823}{doi:10.1016/j.jmateco.2023.102823}.
\href{https://vickiknoblauch.com/research.html}{Author's publication list}.

\bibitem{HKKL}
Frederik S. Herzberg, Vladimir Kanovei, Mikhail G. Katz, and Vassily Lyubetsky,
\emph{Minimal axiomatic frameworks for definable hyperreals with transfer},
Journal of Symbolic Logic \textbf{83} (2018), no.~1, 385--391.
\href{https://doi.org/10.1017/jsl.2017.48}{doi:10.1017/jsl.2017.48}.
\href{https://arxiv.org/abs/1707.00202}{arXiv:1707.00202}.

\bibitem{ProveIt}
ProveIt repository, maintained by Vladimir Reshetnikov,
\emph{Every ordinal as a point-separating game value}, research report,
20 September 2026. File \texttt{ordinal\_separation.tex}, in the
repository's ordinal-order reports, under
\texttt{games-on-ordinals/point-separating-game-values/}.
\href{https://github.com/VladimirReshetnikov/ProveIt/blob/6ea60e3677bd847a00baf023c3c532c83884530e/SetTheory/Cardinals/docs/reports/ordinals-and-order-types/games-on-ordinals/point-separating-game-values/ordinal_separation.tex}{Immutable source at the inspected commit}.

\bibitem{Moore}
Justin Tatch Moore,
\emph{A universal Aronszajn line}, author's manuscript (2008), Section~2.
Cited by source~06 for the characteristic-function representation of a
linear order, not for its forcing-axiom results.
\url{https://pi.math.cornell.edu/~justin/Ftp/Aline_univ.pdf}.

\bibitem{Kuhlmann}
Salma Kuhlmann,
\emph{Isomorphisms of lexicographic powers of the reals},
Proceedings of the American Mathematical Society \textbf{123} (1995),
no.~9, 2657--2662.
Preprint: \url{https://d-nb.info/1102198234/34}.
Corollary~2.4 is the ordinal-exponent embedding theorem (cited by source~07).

\bibitem{Giarlotta}
Alfio Giarlotta,
\emph{On the representability number of lexicographic products in a
Dedekind-complete chain},
Applied Mathematical Sciences \textbf{7} (2013), no.~127, 6347--6353.
\href{https://doi.org/10.12988/ams.2013.39524}{doi:10.12988/ams.2013.39524}.
Used by source~07 for terminology, not as a source of the full-enumeration
results.

\bibitem{Ramakrishnan}
Janak Ramakrishnan,
\emph{Definable linear orders definably embed into lexicographic orders
in o-minimal structures}, arXiv:1003.5400v5 (2010), introduction and
references to Fleischer's 1960/1961 paper and 1963 correction.
\url{https://arxiv.org/abs/1003.5400}.

\bibitem{Tao}
Terence Tao,
\emph{245B, Notes 7: Well-ordered sets, ordinals, and Zorn's lemma
(optional)}, January 28, 2009. Background on canonical ordinal order types
and well-ordering methods (cited by source~06).
\url{https://terrytao.wordpress.com/2009/01/28/245b-notes-7-well-ordered-sets-ordinals-and-zorns-lemma-optional/}.

\bibitem{repo}
Vladimir Reshetnikov,
\emph{ProveIt}, repository snapshot
\nolinkurl{6ea60e3677bd847a00baf023c3c532c83884530e}, inspected 2 October 2026
by sources~06 and~07.
\url{https://github.com/VladimirReshetnikov/ProveIt/tree/6ea60e3677bd847a00baf023c3c532c83884530e}.

\bibitem{reports}
\emph{ProveIt research reports: index and status},
\nolinkurl{SetTheory/Cardinals/docs/reports/README.md}, at the snapshot above.
\url{https://github.com/VladimirReshetnikov/ProveIt/blob/6ea60e3677bd847a00baf023c3c532c83884530e/SetTheory/Cardinals/docs/reports/README.md}.

\bibitem{ProveItCardinals}
\emph{ProveIt Cardinals project README},
\nolinkurl{SetTheory/Cardinals/README.md}, at the snapshot above
(source~07's focused audit).
\url{https://github.com/VladimirReshetnikov/ProveIt/blob/6ea60e3677bd847a00baf023c3c532c83884530e/SetTheory/Cardinals/README.md}.

\bibitem{signsimplicity}
\emph{Simplest elements of convex sets of sign sequences},
\nolinkurl{Surreal/Foundations/SignSequenceSimplicity.lean}, in ProveIt's
\texttt{Algebra/SurrealNumbers} package, at the snapshot above.
\url{https://github.com/VladimirReshetnikov/ProveIt/blob/6ea60e3677bd847a00baf023c3c532c83884530e/Algebra/SurrealNumbers/Surreal/Foundations/SignSequenceSimplicity.lean}.

\end{thebibliography}
\end{document}
